% 本文件由 LaTeX 排版 Agent 根据有效 translation.json 自动生成。
% author=Ambrose; work_id=3409; title=信函集

\chapter{信函集}

% structure: p0001 source=h1 target=omitted reason=page-title-already-rendered-by-parent

信函 17\newline{}信函 18\newline{}信函 20\newline{}信函 21\newline{}信函 22\newline{}信函 40\newline{}信函 41\newline{}信函 51\newline{}信函 57\newline{}信函 61\newline{}信函 62\newline{}信函 63\par

\SourceRecord{Translated by H. de Romestin, E. de Romestin and H.T.F. Duckworth.；Nicene and Post-Nicene Fathers, Second Series；Vol. 10.；Edited by Philip Schaff and Henry Wace.；Buffalo, NY: Christian Literature Publishing Co.,；1896.；<http://www.newadvent.org/fathers/3409.htm>.}

% structure: p0001 source=h1 target=section reason=page-title

\section{第十七封信}

\EditorialNote{本函写于\Person{西马库斯}向\Person{瓦伦提尼安二世}递交请愿书之际。圣\Person{盎博罗削}力劝皇帝，捍卫宗教而非迷信乃其职责所在。该请愿书未获基督徒元老联署，故不代表整个元老院。他警告\Person{瓦伦提尼安二世}，若应允所求，不仅有愧于其父兄在天之灵，更将招致\OriginalTerm{教会}{Church}的谴责。}\par

\Person{盎博罗削}主教，致最蒙福的君王、最虔敬的基督徒皇帝\Person{瓦伦提尼安二世}\par

1. 既然生活在罗马治下的所有人在你们——世间的皇帝与君王——麾下服兵役，你们自身也理当为全能的天主和我们的神圣信仰效劳。因为除非人人都以真理朝拜真天主，即基督徒的天主，万物皆在祂权下的那位，否则\OriginalTerm{救恩}{salvation}便不稳固；因为唯有祂是真天主，应当从心底深处朝拜；因为正如\WorkTitle{圣经}所说：\BibleQuote{外邦人的神都是魔鬼}\BibleRef{咏 96:5}。\par

2. 如今，每个人都是这位真天主的士兵；凡在心灵深处接纳并朝拜祂的人，不会以虚伪或矫饰，而是以真诚的信与虔敬来事奉祂。即便他最终未能达到此境界，至少也不应赞同偶像崇拜和渎神的礼仪。因为无人能欺骗天主，一切——包括内心的隐密——在祂面前都昭然若揭。\par

3. 因此，最虔敬的基督徒皇帝，既然您理应向真天主献上\OriginalTerm{信德}{faith}与热忱、对信仰的关切与敬虔，我实在诧异为何竟有人期望，您会认为下令为外邦诸神恢复祭台，并为此渎神的祭祀提供所需款项，是您的职责；因为凡长久以来由皇室或城市国库所出之资，您若提供，便显得像是从您自己的财产中拨付，而非归还本属于他们的东西。\par

4. 他们如今抱怨自己的损失，可他们从不曾顾惜我们的鲜血，甚至拆毁了教堂的建筑物。他们如今请求您赐予特权，可正是他们，借着\Person{犹利安}的最后法律，剥夺了我们言说与教导的共同权利；那些特权也曾屡次欺骗基督徒；因为借着那些特权，他们试图诱捕一些人，部分出于疏忽，部分为了逃避公共义务的负担；又因并非所有人都意志坚定，即使在基督徒君王治下，也有许多人背弃了信仰。\par

5. 倘若这些事未曾被废除，我本可向您证明，理应由您的权威将其铲除；但既然它们已被许多君王在几乎整个世界下令禁止，且由您仁慈的兄长、可敬的\Person{格拉提安}皇帝出于对真信仰的考量，在\Place{罗马}予以废除，并通过诏书使其失效；我恳求您，既不要拔除已按信仰确立之事，也不要撤销您兄长的命令。在世俗事务上，他若确立了何事，无人会认为可以等闲视之，而关乎宗教的命令却被践踏脚下。\par

6. 不要让任何人利用您的年轻；若提出此要求者是个异教徒，他就不该以自己迷信的束缚捆绑您的心智；相反，他应以其热忱教导并劝诫您，如何为真信仰而大发热心，既然他以全部真情捍卫虚妄之事。我本人也劝您，固然要尊重显赫人物的功绩，但无疑，天主必须被置于万有之上。\par

7. 若我们须就军务进行商议，就应等候富有战争经验者发表意见，并听从其建议；若问题涉及宗教，则应思虑天主。无人会因天主被置于他之前而受损。他自有其见解。您不可强迫一个人违背其意愿，去朝拜他所厌恶的。皇帝啊，也请将同样的自由赋予您自己；若有人无法从皇帝处强求什么，而皇帝若想向他强求，他必会愤愤不平，那么愿他耐心承受。即使异教徒自己也厌恶反复无常的态度，因为每个人都应自由地捍卫并持守自己内心的信仰与决心。\par

8. 然而，若有任何名义上的基督徒，认为应颁布这样的法令，万勿让空洞的言辞误导您的心灵，也不要让虚浮的话语欺骗您。凡如此建议、如此下令者，便是献祭。但一人献祭，总比众人跌倒来得更可容忍。在此，整个基督徒元老院正面临危险。\par

9. 倘若今日有异教皇帝——愿天主禁止此事——为偶像建造祭台，并强迫基督徒聚集在那里，置身于献祭者之中，使得祭台的烟火、\OriginalTerm{亵圣}{sacrilege}的飞尘、焚烧的烟雾呛住信众的气息与咽喉；且在元老院作出判决，而议员们被迫在偶像祭台前起誓后投票（因为他们解释，祭台正是为此目的而设，好使一切会议都在其庇佑下进行——尽管如今元老院中基督徒已占多数）；若一个基督徒被迫如此选择而前往元老院，他会视此为迫害——事实上，他们常被暴力强迫到场。那么，当您作皇帝时，这些基督徒还会被迫在异教祭台前起誓吗？誓言是什么？无非是公开承认您所呼求并监察您诚信的那位拥有大能。可当您作皇帝时，竟有人要求并强求您下令建造祭台，并拨发渎神祭祀的费用。\par

10. 然而，颁布此令不可能不犯亵渎之罪，因此我恳求您不要颁布或命令此事，也不要签署任何此类法令。我，作为基督的司铎，凭着信德呼吁您；我们全体主教也必一同呼吁您，若非此消息如此突然且难以置信，竟说有人在您的御前会议中提出此事，或是元老院请求此事。但元老院决不可能为此事请愿；这只是一小撮异教徒假借公名而已。因为大约两年前，当他们企图作同样的事时，由天主的判断所拣选的\Place{罗马}教会主教圣\Person{达玛苏斯}，送给我一份备忘录，其中许多基督徒元老联署抗议，声明他们从未授予此等权力，不同意异教徒的此类请求，也不予赞同，并公开与私下宣布，若此项法令成立，他们将不会出席元老院会议。难道在您这基督徒的时代，为满足异教徒亵渎的意愿，竟要剥夺基督徒元老的尊严，这才符合您的尊严吗？这份备忘录我呈送给陛下您的兄弟，从中明显看出，元老院并未就迷信的开销作出任何决定。\par

11. 但也许有人会说，当他们提出这些请求时，他们为何先前没有出席元老院会议？他们不出席本身已充分表明了意愿，他们在向皇帝陈词时已说得够多。我们何必惊讶，那些人剥夺了\Place{罗马}个人公民的抵抗自由，还企图使您不能自由地拒绝下令您不赞同的事，也不能坚持您自己的意见？\par

12. 因此，我记起最近托付于我的使团使命，再次凭着信德呼吁您。我呼吁您自己的情感：不要依照异教徒的这项请愿作出答复，也不要在这样的答复上以您的签署附上那亵渎。请将此事呈报给您的父皇\Person{狄奥多西}皇帝，在几乎所有较重要的事务上您都惯于向他请教。没有比宗教更重大之事，没有比信德更崇高之事。\par

13. 若此事为民事诉讼，对方应有答辩之权；但这是宗教事务，而我以主教身份提出诉求。请将送来的备忘录副本交予我，以便我更详尽地答复，然后请教陛下您的父皇完整地审视整个议题，并恩准答复。当然，若另有裁决，我们主教们必不能默然忍受，置之不理；您固然可以来教堂，但会发现那里要么没有司铎，要么就有一位抵制您的司铎。\par

14. 您将如何回答一位司铎对您说的话呢？他会说：\Quote{教会不寻求你的馈赠，因为你用馈赠装饰了异教庙宇。基督的祭台拒绝你的奉献，因为你为偶像筑了祭坛，因为那声音出于你，那手属于你，那签署是你的，那行为是你的。主耶稣拒绝并摈弃你的事奉，因为你事奉了偶像，因为祂曾对你说：\BibleQuote{你们不能事奉两个主人。} \BibleRef{玛 6:24}。奉献于天主的贞女在你面前毫无特权，而\OriginalTerm{维斯塔贞女}{Vestal Virgins}却要求特权吗？你为何寻索天主的司铎，却更喜爱异教徒亵渎的请求？我们绝不能分担他人的错谬。}\par

15. 你对此有何话回答？说你跌倒了，不过是个孩子？在基督内，任何年龄都是成全的；任何年龄都充满了天主。在信德上，不分幼小，因为就连孩童也曾无畏地宣认基督，对抗迫害者。\par

16. 你对自己的兄弟又有何话回答呢？他岂不是要对你说：\Quote{我并不觉得被战胜，因为我留下你作皇帝；我死时并不悲伤，因为有你作我的继承人；我离开皇位时并不哀叹，因为我相信我的命令，特别是关于神圣宗教的命令，将永世长存。我已树立这些虔敬与德行的纪念，我献上这些取自世界的战利品，这些战胜魔鬼的胜利纪念品，这些我从万有之敌手中夺取的纪念，其中蕴含着永恒的胜利。我的仇敌还能从我这里夺走什么呢？你已废除了我的法令，而那举兵攻击我的人却未曾如此做。如今我遭受更可怕的创伤，我的法令竟被我的兄弟废止。我的更好部分因你而处于危险之中，那先前只是肉身的死亡，如今却是名誉的死亡。如今我的权力被撤销，更甚者，竟是被我的家人所撤销，而那甚至是我敌人也称赞的。若你是自愿同意，你就定了我的信德为有罪；若是出于被迫让步，你便出卖了你自己的信德。同样，更严重的是，在你身上我处于危险之中。}\par

16. 你又如何回答你的父皇呢？他会以更深的悲痛对你说：\Quote{孩子，你对我判断太差了，竟以为我能纵容异教徒。从无人告知我\Place{罗马}元老院议事厅内设有祭台，我从不相信如此邪恶之事，竟有异教徒在基督徒与异教徒共聚之处献祭，就是说，异教徒侮辱在场的基督徒，而基督徒被迫参加祭礼。我在帝位时，确曾发生过许多各式罪行。凡被察觉的，我均予惩处；若有人逃脱，岂能说我赞成这些无人禀告之事？你对我判断极差，若你认为是外来的迷信而非我自己的信德保全了帝国。}\par

17. 因此，皇帝陛下，既然您看出，若您颁布此类任何法令，将首先损害天主，其次损害您的父亲和兄长，我恳求您做出那在您看来于天主台前有益于您救恩的事。\par

\SourceRecord{Translated by H. de Romestin, E. de Romestin and H.T.F. Duckworth.；Nicene and Post-Nicene Fathers, Second Series；Vol. 10.；Edited by Philip Schaff and Henry Wace.；Buffalo, NY: Christian Literature Publishing Co.,；1896.；<http://www.newadvent.org/fathers/340917.htm>.}

% structure: p0001 source=h1 target=section reason=page-title

\section{第18封信}

\EditorialNote{\Person{圣安博罗削}对\Person{辛玛库斯}陈情书的答复，他在其中称赞\Person{瓦伦提尼安}之后，处理了陈情书中的三点。他以极其有力的方式回应了对手将\Place{罗马}拟人化的说法，并证明\Person{辛玛库斯}所说的饥荒与异教仪式的停止无关。}\par

\Person{安博罗削}主教致最蒙福的君主、最仁慈的皇帝\Person{瓦伦提尼安}·奥古斯都。\par

1. 既然显赫的城市长官\Person{辛玛库斯}已向陛下您呈递陈情书，要求将\Place{罗马}城\Place{元老院}中被移走的祭坛恢复原处；而您，皇帝，虽然年纪尚轻，经验尚浅，却在信德的力量上已是老手，不赞同异教徒的请求，我一听到消息便呈上了一份请求。在此请求中，虽然我陈述了似乎有必要提出的建议，但我请求将陈情书的副本提供给我。\par

2. 因此，我并非对您的信德有所怀疑，而是担忧风险，并且确信您会给予善意的考虑，所以我在本文中回应陈情书的各项主张，唯一请求是，您不要期待语言的优雅，而要期待事实的力量。因为，正如\WorkTitle{圣经}所教导的，智慧且勤学之人的舌头是金的；它被赋予闪耀的言辞，并因华丽辞藻的辉煌而发光，仿佛某种丰富的色彩，以美丽的外观迷住心灵的眼睛，并以景象眩惑。但这金子，若你仔细察考，外表有价值，内里却是贱金属。我恳求你仔细衡量，察考异教徒的派别：他们的言语，听起来响亮、凝重、宏大，却在为缺乏真理之事辩护。他们谈论天主，却崇拜偶像。\par

3. 那位显赫的城市长官在陈情书中提出了三项他认为有力的主张：一是如他所说，\Place{罗马}要求恢复她的仪式；二是应给她的祭司和维斯塔贞女发放酬金；三是在拒绝支付祭司薪俸后，一场大饥荒随之而来。\par

4. 在他的第一项主张中，\Place{罗马}以悲伤含泪的话语抱怨，如他所说，要求恢复她古老仪式的典礼。他说，这些神圣仪式曾将\Person{汉尼拔}击退于城下，将\Person{赛诺内斯人}击退于\Place{卡比托利欧山}外。于是，在宣扬这些神圣仪式能力的同时，也暴露了它们的软弱。正因为如此，\Person{汉尼拔}长期嘲讽\Place{罗马}仪式，尽管众神为他而战，他却作为征服者兵临城下。既然他们的神明明在为他们作战，他们为何还会陷入围城？\par

5. 我又何必提及\Person{赛诺内斯人}呢？若非一只鹅因受惊的嘎嘎声暴露了他们，\Place{罗马}的残余根本无法阻止他们进入\Place{卡比托利欧山}的最深处。看看吧，\Place{罗马}的神殿拥有怎样的保护者。那时候\Person{朱庇特}在哪里？他是在鹅身上说话吗？\par

6. 但我为何要否认他们的神圣仪式曾为\Place{罗马}人而战呢？因为\Person{汉尼拔}也崇拜同样的神祇。那么，让他们选择吧。如果这些神圣仪式在\Place{罗马}人身上得胜了，那么它们在\Person{迦太基人}身上就是失败的；如果它们在\Person{迦太基人}身上凯旋了，那么它们当然没有给\Place{罗马}人带来益处。\par

7. 那么，就让\Place{罗马}人民那个招人怨恨的抱怨到此为止吧。\Place{罗马}并没有发出这样的指责。她以别的话语说话。\Quote{你们为何每日以无害畜群无用的鲜血玷污我？胜利的纪念品不依赖于畜群的内脏，而依赖于战士的力量。我以不同的纪律征服世界。\Person{卡米卢斯}是我的士兵，他击杀了那些占据了\Place{塔尔佩亚岩}的人，并从\Place{卡比托利欧山}夺回了军旗；勇武击败了那些宗教未曾驱退的人。关于\Person{阿提利乌斯·雷古鲁斯}，我该说什么呢？他献出了自己死亡的事奉。\Person{阿非利加努斯}不是在\Place{卡比托利欧}的祭坛间，而是在\Person{汉尼拔}的战阵中寻得了他的凯旋。你们何必要提出祖先的仪式呢？我憎恨\Person{尼禄}的仪式。我又何必提起那些只称帝两个月的皇帝，以及那些开始与终结紧紧相连的统治者的结局？难道蛮族跨越其边界是件新鲜事吗？难道在他们那些悲惨而史无前例的案例中，这边皇帝被俘，那边世界被征服，就显明了那些承诺胜利的仪式是虚妄的？那时难道没有胜利祭坛吗？我哀悼自己的衰败，我晚年因那羞耻的流血而被玷污。我不以为耻：在晚年与整个世界一同皈依。毫无疑问，没有哪个年龄会因太老而不能学习。让那不能自我纠正的老迈感到羞耻吧。值得称颂的不是年岁的老迈，而是品格的老迈。转向更好的事物并不可耻。我与蛮族唯一共通的，就是从古时我不认识\OriginalTerm{天主}{God}。你们的祭献是一种以兽血洒身的仪式。你们为何在死去的动物中寻求天主的声音？来，在地上学习那天上的战斗；我们虽生活于此，但我们的战斗却在那里。让那创造我的天主亲自将天上的奥秘教导我，而不是那不认识自己的人。关于天主，我除了天主还应信谁？我怎能相信你们，你们承认自己根本不了解自己所崇拜的？}\par

8. 他说，借着一条道路，无人能抵达如此伟大的奥秘。你们所不知道的，我们靠着天主的声音而知道。你们凭幻想所寻求的，我们从天主的智慧和真理本身中寻得了。因此，你们的道路与我们的不同。你们为你们的神向皇帝祈求和平，我们则向\Person{基督}为皇帝们本身祈求和平。你们敬拜自己双手的作品，我们认为任何能被造之物都不应被视为天主。天主不愿在石头中受到敬拜。而且，说到底，你们自己的哲学家们也曾嘲笑过这些事。\par

9. 但若你们因不相信\Person{基督}死了而否认他是天主（因为你们不知道那死亡是肉身的死亡，而非\OriginalTerm{天主性}{Godhead}的死亡，这一事实已使如今所有信他的人都不再死亡），那么你们这些以侮辱表示敬重，又以敬重表示贬低的人，还有什么比这更轻率的呢？因为你们把一块木头当作你们的神。哦，充满侮辱的敬拜！你们不信\Person{基督}能死，哦，建立在尊敬上的扭曲！\par

10. 但是，他说，让祭台归还给偶像，装饰物归还给神龛。让那分享他们迷信的人提出这要求吧；一位基督徒皇帝已经学会只荣耀基督的祭台。他们为什么要求虔诚的手和忠信的唇舌服务于他们的亵渎之举呢？让我们的皇帝之声只发出基督之名，只谈论他所意识到的祂，因为，\BibleQuote{君王的心在主手中}。\BibleRef{箴 21:1}有哪个外教皇帝曾为基督建立过祭台吗？当他们要求归还过去之物时，藉他们自己的例子，他们向我们显示了基督徒皇帝应当对他们所追随的宗教怀有多大的尊敬，因为外教皇帝为他们的迷信奉献了一切。\par

11\Emph{a}. 我们很久以前开始，现在他们跟随那些他们曾排斥的人。我们以流血为荣，他们却被花费所打动。我们视这些为胜利，他们却认为是损失。他们从未给我们带来更大的益处，胜过他们下令殴打、剥夺公权、杀戮基督徒之时。信仰将那被不信视为惩罚之事变成赏报。看他们灵魂的伟大！我们通过损失、匮乏、惩罚而增长；他们却认为没有捐献，他们的礼仪便无法延续。\par

11. 让\OriginalTerm{维斯塔贞女}{Vestal Virgins}，他说，保留其特权。让那些不相信贞洁可以没有赏报而存在的人，让那些不信赖德行的人用利益来鼓励吧。但那些承诺的赏报为他们获得了多少贞女呢？几乎只有七个维斯塔贞女被接纳。看那所有数目：头带和花冠、紫色长袍的染料、随从簇拥的抬轿的排场、最大的特权、巨大的收益，以及规定的童贞期限，聚集在一起的有多少。\par

12. 让他们抬起灵魂和肉身的眼睛，看一个端庄之民，一个纯洁之民，一个童贞的集会。不是头带，而是普通却因贞洁而高贵的面纱装饰她们的头部；美丽的诱惑不是追求而是舍弃；没有那些紫色标志，没有美味奢侈，而是守斋；没有特权，没有收益；总之，一切如此，让人以为她们在履行职责时被禁止享乐。但在履行职责时，其中的乐趣便被唤起。贞洁因其自身的牺牲而增长。那不是童贞，若它用金钱买来，并非因热爱德行而持守；那不是纯洁，若它用金钱竞拍而来，竞标一段时期。贞洁的第一个胜利是战胜对财富的渴望，因为追求利益是对端庄的诱惑。然而，我们且假定应将慷慨的供应给予贞女。那会有多少洋溢到基督徒身上！什么样的库藏能提供如此财富？或者，如果他们以为只应给予维斯塔贞女礼物，他们难道不羞愧：那些在外教皇帝治下将全部据为己有的人，竟以为我们在基督徒君王治下不应有公共份额？\par

13. 他们也抱怨，认为他们的司祭和侍奉者不应被排除在公共资助之外。在这点上，言语的风暴何其汹涌！但在另一方面，新近的法律甚至否认我们继承私人财产的权利，却无人抱怨；因为我们不视之为伤害，因我们不为损失而悲伤。如果一位司铎寻求免除地方公务的特权，他就必须放弃祖传及一切其他财产。如果外教人遭遇此事，他们将会如何投诉：一位司铎必须通过失去全部祖产来购买他职分所需的空闲时间，并以牺牲所有私人财产为代价换取执行公共职务的权力；而且，他声称为了公共安全而守夜，却必须以家庭匮乏的代价来安慰自己，因为他不是出卖了一项服务，而是获得了一项恩惠。\par

14. 比较一下这两种情形。你们想要豁免一位\OriginalTerm{市议员}{decurio}，却不允许教会豁免一位司铎。遗嘱可以为神庙的侍奉者而立，没有任何世俗之人被排除，没有任何卑下之人，没有任何无耻不洁之人，唯独圣职人员被排除在共同权利之外，唯有他们为所有人献上共同的祈祷，并履行共同的服务；甚至连庄重寡妇的遗产遗赠也不被允许，任何礼物也不许。而且在品行上找不到过错时，惩罚仍然加于职务之上。基督徒寡妇遗赠给神庙司祭的财产是有效的，而她留给天主仆人的遗产却是无效的。我叙述这些并非为了抱怨，而是让他们知道我不抱怨什么；因为我宁愿在金钱上更贫乏，也不愿在恩宠上贫乏。\par

15. 但他们说，给予或遗留给教会的财产未曾被触动。他们也应该说明，是谁夺走了献给神庙的礼物，而这事临到了基督徒。如果这些事对外教人做，那非难便不是伤害，而是回报了。难道直到如今，正义才被当作借口，公平才被要求？这种感觉在哪里，当他们掠夺所有基督徒的财产之后，还吝惜他们生命的呼吸，禁止他们使用那任何死者都不被拒绝的最后安葬？大海归还了那些被外教人投入其中的人。这就是信仰的胜利：他们自己现在谴责他们祖先的行为，而他们也定那些行为的罪。但向那些他们谴责其行为的人寻求利益，是什么道理呢？\par

16. 然而，没有人拒绝给予神龛礼物，或给占卜者遗产；只有他们的土地被拿走了，因为他们没有以宗教方式使用他们以宗教权利所主张的东西。如果他们拿我们作例子，为什么他们不去实行我们所做的？教会除了信仰之外没有自己的财产。因此才有她的收获和增长。教会的财产是对穷人的供养。让他们数算一下，神庙赎回了多少俘虏，为穷人贡献了什么食物，为流亡者提供了什么生活资料。那么，他们是土地被拿走，而非权利被拿走。\par

17. 看哪，所发生的事：一场公共饥荒，正如他们所说，报复了那可悲的不虔敬——先前只对祭司们的安适有益的，开始对所有人的用度有益。为此，他们说，树皮从小树林中运来，垂死之人的口吮吸那淡而无味的汁液。为此，他们以谷物换取\Place{凯奥尼亚}的橡实，再度回到牲畜的食料和可怜食物的滋养，他们摇动橡树，在树林中慰藉那可怕的饥饿。这些，真是，地上的新异象，前所未有，正当异教迷信在世上盛行之时！的确，何时曾有过，庄稼用空空的禾秸嘲笑贪婪农夫的祈祷，而在犁沟中寻觅的禾穗令乡民们的希望落空？\par

18. 希腊人从何处获得了他们橡树的神谕，若非他们认为森林食粮的供养是天上宗教的恩赐？因为他们相信这些正是他们的神的恩赐。除了异教民，谁会敬拜\Place{多铎那}的树木，当他们尊崇那可怜的林地食物？他们的神不可能在愤怒时，将平息时赐予的恩物作为惩罚加给他们。若因忧伤于供养被拒绝给予少数祭司，便拒绝给予所有人，那有何公义可言？因为报复比过错更难以承受。因此，此因由不足以在衰败的世界带来如此苦难，以至于整年的希望竟在庄稼青绿时突然毁灭。\par

19. 诚然，许多年前，庙宇的灯火从世上被取走了；难道异教的神至今才终于想起要报复这伤害吗？\Place{尼罗河}未能按惯常的河道泛滥，是为了报复城中祭司的损失，而它却不报复自己的损失吗？\par

20. 但就算他们以为，他们所奉的神明所受的伤害在过去一年间遭到了报复。那为何在今年却未见报复呢？因为如今，乡民既不以拔起的根为食，也不从林中的浆果中寻求饱足，不从荆棘中摘取食物，而是欣喜于他们顺利的劳作，在惊叹收获之时，愿望全然实现，补偿了他们的禁食；因为大地连本带利地出产了。\par

21. 谁如此不谙世事，竟对年岁之差异感到惊奇？然而，即使去年，我们知道许多省区出产丰饶。对\Place{高卢}地区，我该说什么呢？它比平时更加多产。\Place{潘诺尼亚}诸省卖出了他们未曾耕种的谷物，而\Place{第二菲提亚}因其丰饶反受其害，因她习于贫乏以求安全，如今却因丰饶招来敌人。秋季的果实养活了\Place{利古里亚}和\Place{威尼斯}地区。如此，前一年并未因亵渎而枯萎，而去年则因信德之果而繁荣。让他们也否认，若他们有能，葡萄园曾出产极多的果实吧。于是，我们既收获了加倍的庄稼，又享有更丰盛的美酒之益。\par

22. 最后，也是最重要的一点尚存：皇帝们啊，你们是否应当恢复那些曾有益于你们的援助？因为他说：\Quote{让它们保卫你，也让我们敬拜它们。}最忠信的元首们，正是这一点我们无法容忍：他们竟嘲讽我们，说他们奉你们的名义祈求他们的神，未经你们命令，却犯下极大的亵渎，并将你们的闭目不见解释为赞同。让他们拥有自己的守护者吧，让这些神，若他们能够，保护他们的敬拜者。因为，若他们不能帮助那些敬拜他们的人，又怎能保护你们这些不敬拜他们的人呢？\par

23. 但是，他说，我们祖先的礼仪应当保留。然而，既然万物都已向着更美善的方向进步，那又怎样呢？世界本身，起初由元素之种聚结于虚空中，在柔软的球体内，或因工尚未有序而混沌无序，一片黑暗，难道它后来没有（在确立了天空、海洋与大地的区分之后），获得事物借以显得美丽的各种形式吗？大地脱去阴沉的黑暗，惊叹于新出的太阳。白昼并非在起初就闪耀，而是随时间推移，因光的增长而明亮，因热的增长而温暖。\par

24. 月亮本身，在先知预言中用以象征教会，当它再次初升、渐盈至月满之时，起初我们看不见它，隐于黑暗，然后逐渐盈满其角，以至在太阳对面圆盈，皎然明耀辉煌。\par

25. 大地在古时未经为生产果实而耕作的经验；后来，当勤恳的农夫开始主宰田野，用葡萄藤覆盖那无形的土壤时，它便脱去了野性，因家居的耕种而变得温驯。\par

26. 一年本身的最初时节，当万物开始生长时，它以自身的模样感染我们，随后转入春日，花朵即将凋谢，最终在果实中臻至成熟。\par

27. 我们也是如此，年岁尚稚，感官如同婴孩，但随着年岁渐长，逐渐摒弃了才干的初蒙。\par

28. 那么，让他们说，万物都该停留在最初的肇始中，那曾为黑暗笼罩的世界如今因太阳的照耀而明亮，便令人不悦吧。然而，驱散心灵的黑暗比驱散身体的黑暗更为愉快，信德之光比太阳之光更为照耀于人，如此岂不更美好？因此，世界如同万物，其原始的境界已经逝去，为的是让白发苍苍的信德的庄严晚年得以随从。让那些受此触动的人，因收成丰盛到来得晚而指责收成吧；因葡萄的收获在年终而指责葡萄吧；因橄榄是最迟的果实而指责橄榄吧。\par

29. 因此，我们的收获乃是灵魂的\OriginalTerm{信德}{faith}；教会的\OriginalTerm{恩宠}{grace}是功绩的收成，这恩宠自创世之初便在\OriginalTerm{圣人}{Saints}身上兴盛，但在末期却遍及万民，好让众人皆知基督的信德进入了并非愚钝的心灵（因为没有对手便无胜利的冠冕），但先前盛行的意见被推翻后，那真实的便理所当然地被优先选择了。\par

30. 若旧有的礼仪令人喜悦，\Place{罗马}为何还要采用外来的礼仪呢？我略过那被昂贵建筑遮掩的地面，以及闪耀着堕落金光的牧人茅舍。为什么，让我回应他们所抱怨的事吧，他们为何急切地接纳被攻占城市的雕像、被征服的众神，以及异邦迷信的外来礼仪呢？在模仿\Place{阿尔莫}河的溪流中为其战车洗濯的\Person{库柏勒}，其模式源自何处？\Place{弗里吉亚}的吟游诗人，以及常为罗马人所憎恨的不义\Place{迦太基}的众神，又源自何处？至于那位被非洲人尊为\Person{凯勒斯提斯}、波斯人尊为\Person{尼特拉}、多数人尊为\Person{维纳斯}的女神，仅是名称不同，而非神祇有别。因此他们相信\Person{维多利亚}是一位女神，但她实在是一种赐予，而非一种能力；是被赐予的，并不统治，源于军团的助力，而非宗教的力量。那么，靠士兵的人数所声称，或由战事结局所给予的那位女神，能算伟大吗？\par

31. 他们要求在\Place{罗马}城的元老院内为\Person{维多利亚}设立祭台，而那正是聚会的多数人为基督徒的地方！在所有庙宇中都有祭台，\Place{维多利亚女神庙}里也有一座祭台。既然他们喜欢人多势众，便在各处举行祭献。要求只在这一个祭台举行祭献，这不是侮辱信德又是什么呢？岂能容忍一个异教徒祭献而基督徒在场呢？他说：\Quote{让他们吸入吧，让他们吸入吧，即使违背他们的意愿，用他们的眼睛吸入烟雾，用耳朵吸入音乐，用喉咙吸入灰烬，用鼻孔吸入熏香，让从我们祭坛扬起的尘土覆盖他们的脸庞，尽管他们厌恶这些。}难道浴场、柱廊、街道中遍布的偶像还不够他们用吗？在那共同的集会中岂不将有共同的命运？元老院中信实的部分将被那些呼求诸神者的声音、向他们起誓者的誓言所捆绑。如果他们反对，就会显得他们虚伪；如果他们默许，就会承认\OriginalTerm{亵圣}{sacrilege}之事。\par

32. 他说：\Quote{我们该在哪里宣誓服从陛下您的法律和法令呢？}那么，那体现在法律中的您的意旨，难道要靠异教的仪式来获得认同并约束忠诚吗？信德受到攻击，不仅是那些在场者的信德，还有那些不在场者的信德，而且，皇帝们啊，你们的信德也受到攻击，因为你们若下令，便是在强迫。可敬的\Person{君士坦提乌斯}，虽尚未领受\OriginalTerm{神圣奥迹}{sacred Mysteries}，却认为若见到那祭台，自己便会受玷污。他曾下令将它移走，却未下令将它放回。这移走有法令的权威，这放回却没有命令的权威。\par

33. 不要因人不在场便自欺。那内心与他人相连的人，比身体在场而表示赞同的人更为在场。因为心灵相通胜于身体相连。元老院以你们为召集会议的主席，它为你们而聚集；它把良心交给你们，而非交给异教徒的神明；它爱你们胜于自己的儿女，但不会超过它的信德。这是一种值得向往的爱，这爱比任何统治权都大，倘若那维护统治权的信德得以稳固的话。\par

34. 但或许有人会因此动摇，认为果真如此，一位最忠信的皇帝竟被离弃，仿佛功绩的赏报竟要用眼前事物的短暂尺度来衡量。哪个智者不知道，人间事务是在一种循环往复中运行的，因为它们并非总是一帆风顺，而是境况多变，经历兴衰更替。\par

35. 罗马的庙宇曾派遣何人比\Person{格奈乌斯·庞培}更为成功？然而，当他以三次凯旋环绕大地之后，却在战斗中被击败，成为战争中的逃亡者，被流放至自己帝国的疆域之外，最终倒毙于\Place{卡诺帕斯}一个宦臣之手。\par

36. 整个\Place{东方}之地曾赐予世人何等人物比波斯王\Person{居鲁士}更为尊贵？他也一样，征服了那些最强大的反对他的君主，并将这些被征服者作为俘虏囚禁，之后却败亡于一位妇人的武力之下。那位素以礼待战败者而闻名的君王，竟被砍下头颅，置于盛满鲜血的器皿中，并被命令饱饮，如此屈服于妇人之力的嘲弄。因此，在他的一生中，并未受到同等的回报，而是大相径庭。\par

37. 我们又能找到谁比\Place{迦太基}人的领袖\Person{哈米尔卡}更热衷于祭献呢？在整场战役期间，他始终在队伍之间祭献，当他看到己方战败时，竟投身于自己正在照看的火堆，好让他发现对自己毫无益处的那些火焰，用自己的身体一并熄灭。\par

38. 那么，关于\Person{尤利安}，我又该说什么呢？他轻信占卜者的答复，结果毁掉了自己的退路。因此，即便在类似的情形中，过失也不尽相同，因为我们的许诺未曾欺骗过任何人。\par

39. 我已经回答了那些挑衅我的人，仿佛我未曾被挑衅一般，因为我的目的是驳斥那\WorkTitle{陈情书}，而非揭露迷信。但，皇帝啊，让他们的陈情书本身使你更加谨慎。因为它在叙述了从前的君主们——其中较早的遵行他们祖先的礼仪，较晚的也没有废除它们——并且进一步说，如果早期君主的宗教实践不构成先例，那么晚期君主的默许便构成了先例；这就清楚地表明了，你对自己的信仰应负的责任，即你不应追随异教礼仪的榜样，以及你对兄弟之情应尽的本分，即你不应废除你兄弟的谕令。因为，如果他们仅为了他们一方而赞扬那些君主的默许——这些君主虽是基督徒，却丝毫未曾废除异教的法令——那么，你岂不更应当顺从兄弟之爱？这样，你，即使有些事你并不赞同，也应当置之不问，以免减损你兄弟的法规，而如今你更应维护那些你认为既符合自己的信仰，也符合兄弟情谊之事。\par

\SourceRecord{Translated by H. de Romestin, E. de Romestin and H.T.F. Duckworth.；Nicene and Post-Nicene Fathers, Second Series；Vol. 10.；Edited by Philip Schaff and Henry Wace.；Buffalo, NY: Christian Literature Publishing Co.,；1896.；<http://www.newadvent.org/fathers/340918.htm>.}

% structure: p0001 source=h1 target=section reason=page-title

\section{书信二十}

\EditorialNote{\Person{圣安博罗削}向他妹妹叙述在\Place{米兰}发生的事件，与亚略派要求一座大殿有关，以及民众如何起来反对。随后第二天，大殿被士兵占领，但士兵却与天主教徒打成一片。他简述了自己的讲话内容，将他们所受的考验比作\Person{约伯}的，尤其是由他妻子造成的考验，以及其他因女人而起的案例。虽然大殿被交出，但他本人曾受一名书记官的威胁，但这并未使他忧虑。他根据当时情形改编了\Person{约纳}的故事，叙述了民众收回他们教会时的喜悦，\Person{瓦伦提尼安}对其朝臣所说的话，以及\Person{卡利戈努斯}对他的行为。这封信的日期是主历385年复活节。}\par

1. 既然在几乎所有的信中，你都关切地询问教会的情况，你就听听现在发生的事吧。我收到你的信，在信中你说你因梦而不安，就在次日，沉重的忧患便开始令人感受到了。而这一次，被要求交出的不是波尔齐亚大殿，即城外的那座，而是新大殿，即城内那座更大的。\par

2. 首先，一些显要人物，即国家的顾问，恳求我交出大殿，并设法让民众不引发骚乱。我自然回答说，天主的圣殿不能由一位主教交出。\par

3. 次日，这一答复在教会中得到了民众的赞同；\OriginalTerm{总督}{Prefect}来到那里，开始劝说我们至少交出波尔齐亚大殿，但民众大声反对。然后他离开了，示意他将向皇帝报告。\par

4. 第三天，即主日，在诵读和讲道之后，慕道者被遣散时，我正在大殿的洗礼堂向一些候洗者教导信经。就在那里，有人向我报告说，他们从皇宫派来了\OriginalTerm{十夫长}{decani}，正在悬挂帷幔，一部分民众正往那里去。然而，我继续我的职务，并开始举行弥撒。\par

5. 在呈献祭品时，我听说有个名叫\Person{卡斯图鲁斯}的人，亚略派说他是位司铎，被民众抓住了。路人是在街上碰见他的。我开始痛哭，并在祭献中恳求天主来援助我们，不要为教会的事流任何人的血，或者至少，流我的血吧，不仅为了我民众的益处，也为了不信者本身。不多说了，我派了司铎和执事，从暴力中救出了那人。\par

6. 于是，最严厉的判决被宣布了，首先针对整个商人团体。这样，在圣周期间，本应解开欠债者的束缚的时候，却听到了锁链的摩擦声，它们被戴在无辜者的颈上，而且要求在三天之内交出二百磅黄金。他们回答说，只要能够保全信德，即使要求再多或再多一倍，他们也愿意给。监狱里挤满了商贩。\par

7. 皇宫的所有官员，即书记官、专员以及各级官吏的差役，都奉命远离正在发生的事，借口是他们被禁止参与任何骚乱；对于有地位的人，若不交出大殿，则威胁以许多极重的刑罚。迫害正猖獗，如果他们一旦打开闸门，似乎就要爆发出各种暴行。\par

8. \OriginalTerm{将领和保民官}{Counts and Tribunes}来了，催促我尽快交出大殿，说皇帝是在行使自己的权利，因为一切都在他的权力之下。我回答说，如果他向我索取属于我的东西，即我的土地、我的金钱，或任何这类属于我私人的东西，我不会拒绝，尽管我所拥有的一切都属于穷人，但那些属于天主的事物，却不受皇权的支配。\Quote{如果要求我的祖产，就拿去吧；如果要求我的身体，我立刻就去。你们想给我戴上锁链，还是置我于死地？这对我都是一种荣幸。我不会聚众为自己辩护，也不会紧抱祭台哀求饶命，反而更乐意自己为祭台而被杀。}\par

9. 当我得知武装人员被派来占据大殿时，我确实惊恐万分，唯恐民众在保卫大殿时，发生可能导致全城受损的屠杀。我祈求自己不要在如此伟大城市的毁灭后苟活，甚或是整个意大利的毁灭。我害怕流血之恶名，便献出了自己的颈项。有些哥特保民官在场，我向他们搭话，说：\Quote{你们领受了罗马权利这份恩赐，难道是为了成为公众和平的扰乱者吗？如果这里的一切都被毁灭，你们将去向何方？}\par

10. 接着，他们要求我约束民众；我回答说，不去煽动他们，这是我力所能及的；但要平息他们，却在天主手中。如果他们以为我在鼓动他们，就应立即惩罚我，或者把我流放到他们选定的任何荒僻之地。我说完这些话，他们便离去了。我整天待在旧大殿里，然后回家就寝，这样如果有人想抓我，就会发现我已准备好了。\par

11. 天亮前，我离开家时，大殿已被士兵包围。据说，士兵曾向皇帝暗示，如果他愿意出来，他可以出来；他们如果看见他去加入天主教徒，就会随行；否则，他们就会去参加\Person{安博罗削}所召集的集会。\par

12. 没有一个亚略派的人敢出来，因为在公民中没有一个，只有少数皇室成员和一些哥特人。他们从前以车为屋，如今以教会为车。那个女人无论去哪里，都带着她所有的随从。\par

13. 我听到大殿四周充满民众的哀叹声，但在诵读进行时，有人告知我，新大殿也挤满了人，人群似乎比他们全获自由时更多，而且有人正呼唤一名读经员。简言之，那些似乎占据了大殿的士兵，得知我命令民众不要与他们共融，便开始来到我们的集会中。妇女们见此情形，心神不安，有一人冲了出去。但士兵们自己说，他们是为祈祷而来，不是为争斗。民众发出了一些呼喊。他们以极大的节制、极大的恳切、极大的信实，请求我们前往那座大殿。据说，那座大殿里的民众也在要求我前往。\par

14. 于是我开始发表如下讲话。我的孩子们，你们已经听到《\BibleBook{约伯}{传}》的诵读，按照既定的次序和时期，正被宣读。魔鬼也从经验中知道这部书将被解释，书中展示了并阐明了他诱惑的全部力量，所以他今天更猛烈地鼓动起来。但感谢我们的天主，祂以信德和忍耐如此坚固了你们。我本登上讲道台，只为赞扬约伯，而我却发现你们人人都是可赞美的约伯。在你们每个人身上，约伯重活了，在每个人身上，那位圣人的忍耐和勇敢再次闪耀出来。因为基督徒能说出什么比圣神今天在你们内所说的更果断的话呢？奥古斯都啊，我们请求，我们不战斗，我们不惧怕，我们只是请求。这正符合基督徒：既渴望和平与安宁，又不让信德和真理的恒心被恐惧所抑制。因为主是我们的领袖，\Quote{凡仰望祂的，祂就是他们的救主。}\par

15. 但让我们回到眼前的读经。你们看到，天主允许魔鬼试探善者。那恶者嫉妒一切在善中的进步，他以各种方式试探我们。他用财产、儿女和身体的痛苦试探圣约伯。更强的人在自己身上受试探，较弱的人在他人的身上受试探。他也渴望夺去我在你们身上拥有的财富，想要驱散你们安宁的这份祖产。他竭力把你们也从我身边夺走，我善良的孩子们——我为你们每天重新举行\OriginalTerm{祭献}{Sacrifice}——他企图将你们卷入公共动乱的毁灭中。因此，我已经遭受了两种试探。或许因为主我们的天主知道我太过软弱，尚未赐给他控制我身体的权柄。尽管我自己也许渴望，尽管我奉献自己，天主认为我可能不敌这场战斗，而用各种劳苦来锻炼我。约伯不是从这场战斗开始，而是以它结束。\par

16. 但约伯受到接踵而至的凶信的试探，他也受他的妻子的试探，她说：\BibleQuote{你出言攻击天主，然后死去！}\BibleRef{约 2:9} 你们看到，突然激起何等可怕的事：哥特人，武装的人们，异民，对商人们的罚款，圣徒们的苦难。你们注意到，当时下了命令：\Quote{交出大殿；} 那便是说 \Quote{出言攻击天主而死。还不只是出言攻击天主，} 更要行事反对祂。因为命令是：交出天主的祭台。\par

17. 就这样，我们为皇帝的命令所准备，却由圣经的话得以坚固，圣经回答说：\BibleQuote{你说话真像一愚妇。} 那么，这种试探并非轻微，因为我们知道，那些经由女人而来的试探更为猛烈。因为连亚当也因厄娃而跌倒，\BibleRef{创 3:6} 由此他偏离了天主的诫命。当他认识自己的错误，感到负罪的良心所斥责时，本想隐藏自己，却不能隐藏，于是天主对他说：\BibleQuote{亚当，你在哪里？}\BibleRef{创 3:9} 那意思是：你从前是什么？你现在开始在了何处？我曾把你安置在哪里？你流荡到了何方？你承认自己是赤身露体的，因为你失去了良好信德的外袍。那些树叶，是你现在设法用来遮掩自己的。你拒绝了果实，你想要隐藏在法律的树叶下，但你被暴露了。你因一个女人而愿意离开主你的天主，所以你躲避你从前切望看见的那一位。你选择和那女人一起隐藏自己，舍弃了世界的明镜、乐园的住所和基督的恩宠。\par

18. 我何必讲述\Person{依则贝耳}也曾以嗜血的方式迫害\Person{厄里叟}呢？或者\Person{黑落狄雅}\BibleRef{玛 14:3} 使\Person{若翰洗者}被杀害呢？个人迫害个人；但对于我，功劳远不及他们，试炼却更为艰难。我的力量更小，却面临更大的危险。关于女人，变故接踵而至，她们的仇恨交替，她们的虚假变幻莫测，长老们聚集起来，对皇帝的冒犯被用作借口。这样凶猛的试探临到我这个区区虫子的头上，原因何在？除非他们攻击的并非我，而是教会。\par

19. 最后，命令下达：交出大殿。我的回答是：我交出它是非法的，皇帝啊，你接收它也并非益处。你毫无权利侵犯私人的住宅，你以为天主的殿就可以被夺去吗？有人声称，皇帝凡事都可为，万物都属于他。我的回答是：皇帝啊，千万不要让你自己负上这样一种想法的重担，即你对属于天主的事物拥有任何皇权。不要高举自己；但若你愿意长久为王，就当顺服于天主。经上记载：\BibleQuote{把属于天主的，应归还天主；属于凯撒的，应归还凯撒。}\BibleRef{玛 22:21} 宫殿属于皇帝，圣堂属于主教。交给你的权柄只管理公共建筑，而非神圣建筑。又据说皇帝宣称：我也该拥有一座大殿。我的回答是：你拥有它是不合法的。你与淫妇有什么相干？因为凡不以合法婚礼与基督结合的，就是淫妇。\par

20. 当我正论及此事时，有人给我带来消息，说皇家帷幕被取下，大殿中挤满了人，他们呼唤我出现。于是我立刻把话锋转向此事，说道：圣神的谕示何其高深！弟兄们，你们记得，我们在晨祷时曾极之悲痛地答唱：\BibleQuote{天主啊，异民侵入你的产业。} 事实上，异民确实来了，甚至比异民更坏的来了；因为哥特人和各种外族的人都来了；他们带着武器，包围并占据了大殿。我们因未能认识你的伟大而哀叹此事，但我们因缺乏远见而陷入了错误。\par

外邦人来了，实在来到了你的产业中；那些作外邦人而来的，已成了基督徒。那些来侵犯你产业的，已被立为与天主同为承继者。先前我视为仇敌的，如今作了我的保护者。先知论及主耶稣所歌咏的，已经应验：\BibleQuote{他的居所是在平安中，} 且 \BibleQuote{他在那里折断了弓上的火箭、盾牌、刀剑与战事。} 此恩此工，若非你的，又属谁呢，主耶稣？你看见武装的人来到你的圣殿；一方面，百姓哀号聚集，为免显得他们交出了天主的大殿；另一方面，兵士奉命施暴。死亡临于眼前，惟恐混乱中疯狂得以立足。主，你来到中间，使双方合而为一\BibleRef{弗 2:15}。你遏止了武装的人，说：若你们蜂拥操戈，若闭在我殿中的人受惊扰，\BibleQuote{我的血有何益处？} 因此，基督，感谢归于你。不是使者，不是信差，而是你，主，拯救了你的子民，\BibleQuote{你已脱去我的苦衣，给我披上喜乐。}\par

我说了这些话，心中惊奇皇帝的心意竟能被士兵的热忱、伯爵们的恳求及民众的哀告所软化。这时有人告诉我，已派了一名书记官前来，向我传达命令。我稍退后，他便向我传达了命令。他问道：你在想什么，竟敢违抗皇帝的旨令？我回答说：我不知道颁布了什么命令，也不知道有何未经审虑便行之事。他问：你为何派遣司铎去大殿？你若是个暴君，我要知道，好知道如何对付你。我回答说，关于教会的事，我从未仓促行事。当我听说大殿被士兵占据时，我只是更自由地发出叹息；许多人劝我前往那里时，我说：我不能交出大殿，但我也不可争斗。但当我听说皇家的帷幔已被取下，民众催促我前往时，我派遣了几位司铎；我本人没有去，但说：我信基督，皇帝自会与我们交涉。\par

如果这些行为看似暴政，那么我拥有武器，但仅在基督的名内；我有能力献上我自己的身体。我问道：如果他视我为暴君，为何迟迟不打击我？古制，皇权是由主教们授予的，从未自行僭取；而且人们常说，皇帝们渴望司铎职，而非司铎们渴望皇权。基督退避，免得被立为王。我们拥有自己的权力；因为主教的权力就是他的软弱。宗徒说：\BibleQuote{几时我软弱，正是我有能力的时候。}\BibleRef{格后 12:10} 但那天主未曾给他兴起对手的人，当心别为自己制造暴君。\Person{马克西穆斯}并未说我是\Person{瓦伦提尼安}的暴君，他只抱怨说，由于我使节的干预，他未能进入\Place{意大利}。我又加上一句：司铎们从未是暴君，却常遭受暴君之苦。\par

我们在哀伤中度过了那一整天，皇家的帷幔被孩童们嘲笑地撕碎。我无法回家，因为守卫大殿的士兵遍布四周。我们与弟兄们在教堂的小圣堂内咏唱\BibleBook{}{圣咏}。\par

次日，按常规宣读了\BibleBook{约纳}{书}，读毕，我便开始这篇讲道。弟兄们，我们读了一段经文，其中预言罪人将归依。他们的被接纳，是因所将要发生的事，如今被人期盼。我补充说，那位义人也甘愿承受责难，只为不愿看见或宣告城市的毁灭。正因为判决可悲，他也因蓖麻枯萎而难过。天主对先知说：\BibleQuote{你因这蓖麻而悲伤吗？} \Person{约纳}回答：\BibleQuote{我悲伤。}\BibleRef{纳 4:9} 主然后说，若他为蓖麻枯萎而伤痛，祂自己岂不更该关心如此众多人民的得救？因此祂撤回了为全城所预备的毁灭。\par

未再迟延，便有消息传来，皇帝已命令士兵撤出大殿，且先前向商人强索的款项必须归还。于是，全体民众何其欢乐！他们的喝彩何其正当！他们的感谢何其洋溢！那天正是主为我们被交付的日子，教会中补赎得以松弛的日子。士兵们争先恐后地传送这些消息，奔向祭台，互吻，以示和平的标记。那时我明白，天主击打了那清晨的虫子，为保全全城。\par

这些事已过，但愿一切止息！但皇帝激动的言辞预示更严重的未来患难。我被称为暴君，且甚于暴君。当伯爵们恳请皇帝前往教堂，并说他们是应士兵之请求这样做时，他回答：若\Person{安博罗削}命令你们，你们便会将他锁链交予他。这话之后，你们可想象还会发生什么。众人闻之无不战栗，但他身边有些人煽动他。\par

后来，连内侍长\Person{加里贡}，也用奇特的语言大胆向我说话。他说：只要我还活着，你就敢轻视\Person{瓦伦提尼安}？我要取下你的头。我回答说：愿天主允你实现你的威胁；因为那样，我将如主教们一样受苦，而你则行阉宦之事。但愿天主使他们远离教会，让他们将所有武器对准我，用我的血满足他们的饥渴。\par

\SourceRecord{Translated by H. de Romestin, E. de Romestin and H.T.F. Duckworth.；Nicene and Post-Nicene Fathers, Second Series；Vol. 10.；Edited by Philip Schaff and Henry Wace.；Buffalo, NY: Christian Literature Publishing Co.,；1896.；<http://www.newadvent.org/fathers/340920.htm>.}

% structure: p0001 source=h1 target=section reason=page-title

\section{书信二十一}

\Emph{圣\Person{安博罗削}为自己未应召赴\OriginalTerm{御前会议}{consistory}辩解，理由是在信德问题上，唯有主教才能正当地裁断，且他并非顽抗，因他不愿自己的职级遭受不公。他又补充说，\Person{奥克森提乌斯}或许会选择犹太人或不信者、亦即敌视基督的人作审判者。他还表明愿在\OriginalTerm{教务会议}{synod}上讨论争议事项，并说他本可当面告知皇帝今日信中所言，但其他主教及人民不容他如此行。}\par

\Person{安博罗削}主教致最仁慈的皇帝、蒙福的\OriginalTerm{奥古斯都}{Augustus}\Person{瓦伦提尼安}。\par

1. \Person{达尔马提乌斯}，那位\OriginalTerm{保民官}{tribune}兼\OriginalTerm{书记员}{notary}，奉陛下之命召见我，据他所称，要求我也应选择审判者，一如\Person{奥克森提乌斯}所做的那样。他没有提及被要求者的名字，但补充说，将在\OriginalTerm{御前会议}{consistory}中进行辩论，并由陛下的虔敬作出裁决。\par

2. 对此，我作出如下我认为恰当的答复。当我说您那威名显赫的先父不仅口头上如此答复，且以其法律予以确认，即在信德问题或任何教会训令上，应由职位合宜且公正不偏者审判，因为\OriginalTerm{敕令}{rescript}的原话如此——也就是说，他期望司铎审判司铎。此外，若一位主教被指控其他事宜，且需调查其品行，他也意将此保留于主教审判之下。\par

3. 那么，谁以违抗态度回答了陛下呢？是那愿您像您父亲的人，还是那愿您不像他的人呢？除非，或许，如此伟大皇帝的判断在某些人看来微不足道——他的信德已由其宣认的恒心所证实，他的智慧已由国家的持续改善所彰显。\par

4. 最仁慈的皇帝啊，您可曾听说平信徒在信德问题上审判主教？我们岂因某些人的谄媚而如此卑屈，以致忘却\OriginalTerm{司铎职}{priesthood}的权利？我岂会认为能将天主赐予我者托付他人？若主教须受平信徒教导，后果将如何？让平信徒辩论，主教聆听；让主教向平信徒学习。然而，毫无疑问，无论我们查阅\WorkTitle{圣经}或古老年代，谁能否认，在信德问题上——我说的是信德问题——习惯上是主教审判基督徒皇帝，而非皇帝审判主教。\par

5. 您将靠天主的恩宠达到更成熟的年龄，那时您便可判断，那种将司铎职权利拱手让给平信徒的主教是何等样人。您的父亲，靠天主的恩宠，是一位年龄更成熟的人，他常说：审判主教不是我的事。如今陛下却说：我应当审判。他虽在基督内受洗，仍自认不堪承担如此审判的重任；而陛下，您尚未为自己赢得\OriginalTerm{圣洗圣事}{sacrament of baptism}，却将关于信德的审判据为己有，尽管您对那信德的\OriginalTerm{圣事}{sacrament}尚一无所知？\par

6. 他选择了何种审判者，可以留待想象，因为他不敢公布他们的名字。让他们径直来到教会，若有人愿来；让他们与民众一同聆听，并非让每个人坐堂审判，而是各自省察自己的心意，选择跟随谁。这事关系到该教会的主教：若民众聆听他，认为他的论证有力，就跟随他，我不会嫉妒。\par

7. 我且不提民众自己已作出判断。我缄默不谈他们曾向您父亲要求现在这位主教的事实。我缄默不谈您父亲的许诺：若当选者愿意接受主教职务，便得安宁。我正是依据这些许诺行事的。\par

8. 但若他夸耀某些\OriginalTerm{外地人}{foreigners}的认可，就让他到那里去做主教，那里的人认为他应得主教之名。因为我既不承认他是主教，也不知他从何而来。\par

9. 皇帝啊，当您已宣布判断且已颁布法律，以致不容任何人另作判断时，我们又如何能解决此事？而当您为他人立下此法时，您自己也同样受约束。因为皇帝首先当遵守自己所颁布的法律。那么，您是否希望我尝试，看那些被选为审判者的人，如何要么违背您的决定而来，要么至少推辞说，他们不能违背皇帝如此严厉、如此严格的法律？\par

10. 但这岂是违抗者所为，非明理者所为。看哪，皇帝，您自己已在部分地撤销您的法律；但愿不仅是部分地，而是完全撤销！因为我不愿您的法律被置于天主的法律之上。天主的法律已教导我们当遵循什么；人类的法律无法教导我们此事。它们通常能从惧怕者那里苛索改变，却无法激发信德。\par

11. 那么，谁读了这样的命令——转瞬间，敕令传遍众多行省，凡违背皇帝者将被斩首，凡不交出天主殿堂者将立即处死——谁会，我说，谁能独自或与少数几人，对皇帝说：我不赞许您的法律？司铎不得如此说，难道平信徒就可以吗？而那指望得恩宠或害怕得罪的人，竟能审判信德吗？\par

12. 最后，我自己要选择平信徒做审判者吗？他们若坚持信德的真理，就会按那有关信德的法律而被放逐或处死。那么，我岂要让这些人要么否认真理，要么遭受惩罚？\par

13. \Person{盎博罗削}本人并非重要到足以降低\OriginalTerm{司铎职}{priesthood}的地步。一个人的生命远不如全体\OriginalTerm{司铎}{priests}的尊严宝贵。我正是遵照他们的建议发出了那些指示，当时他们指出：\Person{奥森提乌斯}有可能推选某个\OriginalTerm{异教徒}{heathen}或\OriginalTerm{犹太人}{Jew}，若是将关于\Person{基督}的判断交托给这样的人，我便无异于让他们对\Person{基督}得胜。除了听到对\Person{基督}的侮辱，还有什么能取悦他们呢？除非（但愿天主不许）\Person{基督}的\OriginalTerm{天主性}{Godhead}遭到否认，还有什么能让他们欢喜呢？很明显，他们与\OriginalTerm{亚略派}{Arian}沆瀣一气，后者声称\Person{基督}是\OriginalTerm{受造物}{creature}，而\OriginalTerm{异教徒}{heathen}和\OriginalTerm{犹太人}{Jews}也极为乐意承认这一点。\par

14. 这正是在\OriginalTerm{里米尼会议}{Synod of Ariminum}中所决议的，而我理所当然地憎恶那次会议，因为我遵循\OriginalTerm{尼西亚大公会议}{Nicene Council}的准则，死亡或刀剑都不能使我与之分离。陛下之父、至福的皇帝\Person{狄奥多西}也曾赞同并遵循这信仰。\Place{高卢}和\Place{西班牙}持守这信仰，并以虔敬的\OriginalTerm{圣神}{Divine Spirit}宣认加以维护。\par

15. 若需要讨论什么，我学会在\OriginalTerm{教会}{church}中讨论，一如我之前的先辈所作。若要就\OriginalTerm{信仰}{faith}举行会议，便应当有\OriginalTerm{主教}{Bishops}的集会，如同在可敬的君主\Person{君士坦丁}治下所行的那样，他并未预先颁布任何法律，而是将决定留给主教们。在可敬的皇帝\Person{君士坦提乌斯}治下也是如此，他是父王尊位的继承者。然而，开端虽好，结局却不然，因为主教们起初签署了一份纯正的\OriginalTerm{信仰宣认}{confession of faith}，但后来有些人渴望在宫殿内就信仰做出决定，便设法以欺诈手段修改主教的决议。不过，他们随即撤回了这一颠倒的决定，而且毫无疑问，在\Place{里米尼}，多数主教都赞同\OriginalTerm{尼西亚大公会议}{Nicene Council}的信仰，谴责了\OriginalTerm{亚略派}{Arian}的主张。\par

16. 倘若\Person{奥森提乌斯}诉诸\OriginalTerm{会议}{synod}，以讨论有关\OriginalTerm{信仰}{faith}的要点（尽管没有必要因一人而惊动如此众多的\OriginalTerm{主教}{Bishops}，即使他是一位来自天上的天使，也不应置于\OriginalTerm{教会}{Church}和平之上），当我听闻有会议即将召开时，我也绝不会缺席。那么，你若想进行\OriginalTerm{辩论}{disputation}，就请废除那条法律。\par

17. 皇帝陛下，我本愿前往您的\OriginalTerm{御前会议}{consistory}，并在您面前陈明这些话，但\OriginalTerm{主教}{Bishops}们或民众都未允许我，因为他们说，有关\OriginalTerm{信仰}{faith}的事务应当在\OriginalTerm{教会}{church}中、在民众面前处理。\par

18. 皇帝陛下，我但愿您未曾下旨判我\OriginalTerm{流放}{banishment}，并任由我随意前往。我每日外出，无人看守我。您本当为我指定一处您所愿的地方，因为我甘愿献身于任何事。然而，如今\OriginalTerm{圣职人员}{clergy}却对我说：\Quote{无论是自愿离开\Person{基督}的\OriginalTerm{祭台}{altar}，还是背弃它，都差别不大，因为若你离开，便是背弃了它。}\par

19. 我但愿能明确确知，\OriginalTerm{教会}{Church}绝不会被交给\OriginalTerm{亚略派}{Arians}。那样的话，我便会甘愿顺从您的\OriginalTerm{虔敬}{piety}旨意。但如果只是我一人犯了\OriginalTerm{扰乱}{disturbance}之罪，为何竟有命令要侵占所有其他的\OriginalTerm{教会}{churches}呢？但愿能确立这样一条规定：无人可扰乱众\OriginalTerm{教会}{churches}，然后无论何种判决看起来对我合适，我都甘愿接受。\par

20. 那么，皇帝陛下，请俯允接纳我未能前往\OriginalTerm{御前会议}{consistory}的理由。除了为您的事务，我从未学会在御前会议中露面，而且我也不能在宫殿内辩论，因为我既不知道，也不想知道宫殿的秘密。\par

\Signature{我，主教\Person{盎博罗削}，谨向至为仁慈的皇帝、至为有福的奥古斯都\Person{瓦伦提尼安}呈上此\OriginalTerm{奏章}{memorial}。}\par

\SourceRecord{Translated by H. de Romestin, E. de Romestin and H.T.F. Duckworth.；Nicene and Post-Nicene Fathers, Second Series；Vol. 10.；Edited by Philip Schaff and Henry Wace.；Buffalo, NY: Christian Literature Publishing Co.,；1896.；<http://www.newadvent.org/fathers/340921.htm>.}

% structure: p0001 source=h1 target=section reason=page-title

\section{第二十二封信}

\EditorialNote{\Person{圣盎博罗削}在给他姊妹的信中，记述了发现\Person{圣格瓦修斯}和\Person{圣普罗塔修斯}遗体的经过，以及他当时对民众的讲话。他根据圣咏第十九篇讲道，寓意地讲解高天代表殉道者和宗徒，而日则代表他们的宣认。他们曾为天主所贬抑，然后又被举扬。接着他记述了他们遗体被发现时的状况，以及迁葬至圣殿的情形。在另一篇讲道中，他谈到公教徒的喜乐和亚略派的恶意，后者如同犹太人惯常所做的那样，否认所行的奇迹；并指出他们的信德与殉道者们的信德截然不同，且既然魔鬼承认天主圣三，而他们不承认，那么他们比魔鬼本身更不如。}\par

致那位女士，我的姊妹，比我的眼目和生命更亲爱的，\Person{盎博罗削}主教。\par

1. 我不愿在你离开期间这里所发生的事不为你的圣德所知，你当知道我们找到了几具圣殉道者的遗体。因为我祝圣了圣殿之后，许多人仿佛异口同声地开始对我说：\Quote{请像你祝圣罗马圣殿那样祝圣这座圣殿。} 我回答说：\Quote{当然，如果我找到任何殉道者的圣髑的话。} 于是，立刻有一种先知般的热情似乎进入了我的心。\par

2. 我何必多言？天主眷顾了我们，因为就连奉命在\Person{圣斐理斯}和\Person{圣纳博尔}圣所屏栏前清除泥土的神职人员也感到害怕。我找到了适当的标记，当带进一些将要接受覆手的人时，圣殉道者的能力便如此显赫，以致我还未开口，就有一人被抓住并摔倒在神圣的安葬之处。我们发现了两具身材魁梧、实属古时伟人的遗体。所有骨骸完好，且有大量血迹。在那整整两天里，人群极多。我们简略地将一切安排就绪，到了傍晚，便将他们迁至\Place{福斯塔}圣殿，在那里守夜，有些人接受了覆手。次日，我们将圣髑迁至名为\Place{盎博罗削}的圣殿。在迁移途中，一位盲人获得了痊愈。当时我对民众讲话如下：\par

3. 当我看到在此聚集的、如此众多且前所未有的人数，以及闪烁于圣殉道者身上那位天主恩宠的恩赐时，我必须承认，我感到自己不堪担当此任，无法用言语表达我们几乎不能以心灵领会、以眼目目睹的事。但当圣经开始宣读时，那曾在先知们内发言的圣神，恩赐我说出一些堪当如此盛大的聚会、你们的期待以及圣殉道者的功劳的话。\par

4. \BibleQuote{高天陈述天主的光荣}，经上说。当这篇圣咏被宣读时，人便会想到，似乎不是物质的元素，而是天上的功劳，呈献出堪当天主的赞颂。且因今日诵读的经文的巧合，清楚地显示出什么是\BibleQuote{高天}在陈述天主的光荣。请看在我左右两边的圣髑，请看这些与上天交往的人，请看这天上心灵的奖赏。这些就是陈述天主光荣的高天，这些就是穹苍所宣扬的祂的杰作。因为不是世俗的诱惑，而是天主恩宠的化工，将他们提升至至圣苦难的穹苍，且早在以前，便以他们的品格和德行的见证，证明他们对抗这世界的危险时，屹立不摇。\par

5. \Person{保禄}就是一个高天，当他说：\BibleQuote{我们的家乡本是在天上}。\BibleRef{斐 3:20} \Person{雅各伯}和\Person{若望}也是高天，他们被称为\BibleQuote{雷霆之子}；\BibleRef{谷 3:17} 而\Person{若望}，仿佛就是一个高天，看见了与天主同在的圣言。\BibleRef{若 1:1} 主耶稣自己就是永恒之光的高天，当祂在宣扬天主的光荣，那从前无人见过的光荣时。因此祂说：\BibleQuote{从来没有人见过天主，只有那在父怀里的独生子，身为天主的，祂给我们详述了}。\BibleRef{若 1:18} 如果你寻求天主的杰作，请听\Person{约伯}说：\BibleQuote{天主的神造了我}。\BibleRef{约 33:4} 这样，他坚强抵抗魔鬼的诱惑，步步谨慎不犯过。但让我们继续看下文。\par

6. \BibleQuote{日与日侃侃而谈}，经上说。看，这就是真正的日，其间没有黑夜的黑暗介入。看，这些充满生命和永恒光辉的日子，他们宣讲了天主的话，不是用那消逝的言语，而是以他们内心深处，藉着宣认的恒心和见证的坚忍。\par

7. 另一篇读过的圣咏说：\BibleQuote{谁能相似主，我们的天主？祂坐在高天之上，俯视天上地下的卑微事物}。主确实垂顾了卑微的事物，当祂向祂的教会启示了那些隐藏在不为人知的草地下、圣殉道者的圣髑时，他们的灵魂已在天堂，身体则在地下：\BibleQuote{从尘埃中提拔卑微者，从粪土中扶起贫穷人}，你且看祂如何\BibleQuote{使他们与祂子民的诸侯同席}。除了圣殉道者，我们还能认谁为民众的诸侯呢？在他们之中，\Person{普罗塔修斯}和\Person{格瓦修斯}，长久以来默默无闻，如今已加入其中，他们使迄今没有殉道者的\Place{米兰}教会，如今如同多子的母亲，因她自身苦难的殊荣和实证而欢欣。\par

8. 这也不可与真信德相悖：\BibleQuote{日与日传递言语}；灵魂与灵魂，生命与生命，复活与复活；\BibleQuote{夜与夜传授知识}；就是肉身与肉身，即那些以苦难向众人显示了信德的真知识的人。这些夜是美好的，明亮的夜，并非没有星辰：\BibleQuote{因为星辰与星辰的光辉有区别，死人的复活也是这样}。\BibleRef{格前 15:41}\par

9. 因为许多人并非无缘无故地称这为殉道者的复活。我不说他们是为自己复活了，对我们而言，殉道者确实复活了。你们知道——不，你们已亲眼看见——许多人摆脱了邪魔，又有许多人，伸手触摸了圣人的外袍，便从困扰他们的疾病中获得解脱；你们看见古时的奇迹得以重现，当主耶稣降临时，恩宠更丰沛地倾注于大地，许多身体似乎因接触圣躯的阴影而痊愈。多少手帕在传递！多少衣服，放在圣髑上并获得治愈的能力，被人索求！所有人都乐于触摸哪怕外面的线缝，凡触摸的都将获得痊愈。\par

10. 感谢你，主耶稣，因为你在此刻为我们激发了圣殉道者的灵魂，正当你的教会需要更大护佑之时。愿众人皆知我渴望怎样的勇士，他们能防守，却不欲攻击。圣洁的子民，我为你们赢得了这样的勇士，他们能帮助众人，却不伤害任何人。我渴望的正是这样的守护者，我拥有的正是这样的士兵，不是这世界的士兵，而是基督的士兵。我因他们不畏惧任何恶意，他们的代祷越是有力，安全就越大。我也愿那些嫉妒他们赐给我的人，得到他们的护佑。让他们前来，看看我的侍从。我不否认，我被这样的刀剑环绕：\BibleQuote{有人信赖战车，有人信赖战马，我们却要凭上主天主的名号夸耀。}\par

11. 圣经的记述告诉我们：\Person{厄里叟}被阿兰人的军队包围时，对他那害怕的仆人说，不要怕；\BibleQuote{因为，}他说，\BibleQuote{与我们同在的，比反对我们的更多。}为了证明这话，他祈求开启\Person{革哈齐}的眼目，眼目一开，\Person{革哈齐}便看见无数天使军旅在场。我们虽然不能看见他们，却能感受到他们的临在。我们的眼目曾紧闭，因为圣人们的遗体隐藏不露。主开启了我们的眼目，我们便看见那常保护我们的援助。我们过去未曾看见它们，却早已拥有它们。就这样，仿佛主在我们恐惧时对我们说：\Quote{看，我赐给了你们何等伟大的殉道者！}于是我们以开启的眼睛看见主的光荣，这光荣已见于殉道者的苦难，并显于他们的德能。弟兄们，我们已摆脱不小的羞耻；我们曾有主保，却不自知。我们发现了一件似乎超越前辈的事：他们失去的关于殉道者的知识，我们重新获得了。\par

12. 光荣的圣髑从不体面的葬处取出，凯旋的标记展现在穹苍之下。墓穴被血润湿。血腥胜利的痕迹现前，圣髑按原位安放未乱，头颅与身体分离。老人们如今反复说，他们曾听说过这些殉道者的名字，诵念过他们的尊号。那曾掠走各处殉道者的城市，却失去了自己的殉道者。这虽是主的恩赐，我仍不能否认主耶稣将恩宠赐予我司祭职的时期；既然我本人不堪为殉道者，我便为你们获得了这些殉道者。\par

13. 愿这凯旋的牺牲祭品被带到基督为祭品的地方。他是在祭台之上，为众人受苦的那一位；他们是在祭台之下，因他的苦难得救赎的人。我原为自己预定了这地方，因为司铎理当安息在他惯常献祭之处；但我将右方的位置让给这神圣的祭品；那地方本属于殉道者。那么，让我们安放圣髑，将它们安置在相称的安息之所，并以忠信的热忱庆祝整日。\par

14. 民众呼喊，要求将殉道者的安葬礼推迟到主日，但最终商定在次日举行。次日，我又向民众讲了一篇这样的道理。\par

15. 昨天，我按自己的能力解释了这句经文：\BibleQuote{日复一日传出言语；}今天，圣经在我看来不仅在过去预言，甚至仍在当下预言。因为我看见你们日夜延续这圣善的庆典，先知诗歌的宣示便指明，昨天与今天正是最适宜说这话的日子：\BibleQuote{日复一日传出言语；}而昨夜与今夜，正是最适宜说这话的夜晚：\BibleQuote{夜复一夜透露知识。}这两天，你们以心灵深处除了宣扬天主的圣言外，还做了什么？你们证实了自己拥有信德的知识。\par

16. 那些惯常如此的人对你们的这庆典心怀怨恨；出于嫉妒，他们不能忍受这庆典，便憎恨其根源，甚至疯狂到否认殉道者的功劳，而邪魔都承认他们的作为。但这不足为奇，因为不信者的无信如此之深，常使得魔鬼的宣认更易忍受。魔鬼曾说：\BibleQuote{天主子，耶稣，时候未到，你就来这里苦害我们吗？}\BibleRef{玛 8:29} 而犹太人听见了，却自己否认他是天主子。此时你们也听见魔鬼喊叫，向殉道者承认他们不能忍受其苦难，说：\Quote{你们为什么这样严厉地折磨我们？}亚略派却说：\Quote{他们不是殉道者，不能折磨魔鬼，也不能救助任何人；}但魔鬼的痛苦由它们自己的话证实，而殉道者的恩惠由痊愈者的复原和受缚者的解脱得以彰显。\par

17. 他们否认那盲人复明，但他不否认自己痊愈了。他说：我从前看不见，现在能看见了。他说：我摆脱了目盲，以事实为证。那些无法否认事实的人，却否认恩惠。这人是众所周知的：他健康时曾在公务中服务；名叫\Person{塞维鲁}，职业是屠夫。当这障碍临到他时，他已放弃了这行业。他请那些曾资助过他的人作证；他提出那些能证明他复明事实的人，就像他曾有目盲的证人一样。他声称，当他触摸那覆盖圣髑的殉道者衣袍的边缘时，他的视力就恢复了。\par

18. 这不正与我们在\WorkTitle{福音}中所读到的相似吗？在每一情形中，我们赞颂同一创造主的德能；工作与恩赐并不相悖，因为祂在工作中赋予恩赐，又在恩赐中施展工作。祂让别人去完成的，正是祂的名在别人的工作中所实现的。所以我们在\WorkTitle{福音}中读到，犹太人看到瞎子得到治愈的恩赐后，叫来他的父母作证，问道：\BibleQuote{你的儿子怎么能看见？}他回答说：\BibleQuote{我以前是瞎子，现在能看见了。}\BibleRef{若 9:25} 在这个事例中，这人说：\BibleQuote{我是瞎子，现在能看见了。}你们若不信我，就去问别人；若以为他父母与我串通，就去问陌生人。这些人的顽固比犹太人更可恨，因犹太人虽怀疑，至少还问了他的父母；而这些人却暗中调查，公开否认，不信的不是工作，而是工作的创造主。\par

19. 但我要问，他们不信的是什么？是怀疑殉道者能帮助人吗？这等于不信基督，因为祂自己说过：\BibleQuote{你们将要作比这些更大的事。}\BibleRef{若 14:12} 如何作呢？藉着那些殉道者，他们的功绩早已生效，他们的遗体早已寻获？在此我要问，他们是怨恨我，还是怨恨圣殉道者？若怨恨我，难道我行了什么奇迹？是藉我或凭我的名吗？那么，为何怨恨于我，而这些事并非出于我？若是怨恨殉道者（因为若他们不怨恨我，那就只能是怨恨殉道者），这表明殉道者的信仰与他们所信的并非同一信仰。否则，他们不会反对殉道者的工作，除非他们认为殉道者没有那曾存于他们内的信仰，那依祖先传统所立的信仰，这信仰连魔鬼自己都不能否认，亚略派却加以否认。\par

21. 我们今天听见那些经覆手的人说，除非信圣父、圣子和圣神，无人能得救；凡否认圣神、不信天主圣三全能德能的人，是已死被埋葬的人。魔鬼承认这一点，亚略派却拒绝承认。魔鬼说：凡否认圣神天主性的人，就让他像魔鬼自己被殉道者折磨那样，受同样的折磨。\par

22. 我不接受魔鬼的作证，而是接受他的招认。魔鬼是不情愿说话，是在被迫和受折磨下说的。邪恶所压制的事，酷刑会逼露出来。魔鬼屈服于鞭打，亚略派却还未学会屈服。他们受了多少苦，却仍像法郎一样，因患难而心硬！如经上所载，魔鬼说：\BibleQuote{我知道你是谁，你是永生天主之子。}\BibleRef{谷 1:24} 而犹太人说：\BibleQuote{我们不知道祂是从哪里来的。}\BibleRef{若 9:30} 恶魔今天、昨天和夜间说：我们知道你们是殉道者。而亚略派却说：我们不知道，我们不愿明白，我们不愿相信。恶魔对殉道者说：你们来是为毁灭我们。亚略派却说：魔鬼的折磨不是真实的，而是虚构捏造的故事。我听说过许多捏造的事，但从未有人能假装自己是恶魔。我们在那些受覆手的人身上看到的折磨是什么意思？这里有何欺诈的余地？有何伪装的嫌疑？\par

23. 但我不会利用恶魔的声音来支持殉道者。他们神圣的苦难，由他们施予的恩惠得以证实。这些事有判断者，即那些被洁净的人，也有见证者，即那些被释放的人。那声音比魔鬼的声音更好，是那些来时患病、离开时健全的人所发出的；更好的声音是由血发出的，因为血有洪亮的声音，从地上直达天上。你们读过天主怎么说：\BibleQuote{你弟弟的血由地上向我喊冤。}\BibleRef{创 4:10} 这血以其颜色呼喊，这血以其效力的声音呼喊，这血以其苦难的胜利呼喊。我们已答应了你们的要求，将原定昨天举行的安放圣髑典礼推迟到今天。\par

\SourceRecord{Translated by H. de Romestin, E. de Romestin and H.T.F. Duckworth.；Nicene and Post-Nicene Fathers, Second Series；Vol. 10.；Edited by Philip Schaff and Henry Wace.；Buffalo, NY: Christian Literature Publishing Co.,；1896.；<http://www.newadvent.org/fathers/340922.htm>.}

% structure: p0001 source=h1 target=section reason=page-title

\section{第40封信}

\EditorialNote{圣盎博罗削恳请德奥多西听他进言，因为他若保持沉默，则双方都将面临极大危险。他指出，德奥多西虽敬畏天主，仍可能被误导，并指出关于重建犹太人会堂的决定充满危险，会使主教面临背弃真理或死亡的危险。文中提到了尤利安的案例，并回应了皇帝批示的理由，特别是犹太人已烧毁多座教堂的辩护。圣盎博罗削提及瓦伦廷派的庙宇，他宣称其异端更甚于外教人，并指出这将为犹太人的诽谤和对基督本人的嘲弄敞开大门。最后，以马克西穆斯为例警告皇帝不要偏袒犹太人或异端者，并劝勉他施行仁慈。}\par

\Person{盎博罗削}主教致最仁慈的君王、真福皇帝\Person{德奥多西}·奥古斯都。\par

1. 真福皇帝，我虽时常为几乎不断的忧虑所困扰，但从未像现在这样焦虑，因为我看到我必须小心，免得有任何事可归咎于我，甚至略有亵圣的意味。所以我恳求你耐心听我所言。因为，如果我不配蒙你俯听，那我也不配为你献祭，我原是被你托付了你的誓愿和祈祷的。你难道不愿意亲自听他，而你却希望他为你而蒙俯听？你难道不愿意听他为自己辩护，而你却为他人的缘故听了他？你难道不担心自己的决定，怕因你认为他不配蒙你俯听，反而使他也不配为你而蒙俯听？\par

2. 然而，拒绝言论自由，既非皇帝所为，而不说心中所想，也非司铎之职。因为在你们皇帝身上，最受欢迎和最可贵的，莫过于尊重自由，即使对那些在军事上服从你的人们也是如此。因为善君与恶君的区别正在于此：善君爱自由，恶君好奴役。而在一位司铎身上，对天主而言最危险的事，或在人看来最卑鄙的事，莫过于不自由地说出心中所想。因为经上记载：\BibleQuote{我在君王面前传述你的法度，并不羞耻；}另一处又说：\BibleQuote{人子，我立你作以色列家族的守卫，为的是，}经上说，\BibleQuote{如果义人离弃他的义德而行不义，因为你没有警告他，}就是没有告诉他要提防什么，\BibleQuote{他的义德将不被记念，我要向你追讨他的血债。但如果你警告义人，使他不犯罪，而他没有犯罪，他必因受到警告而生存，你也救了自己的灵魂。}\par

3. 皇帝啊，我宁愿在善事上与你有份，也不愿在恶事上与你有份，因此司铎的沉默应使你的仁慈不悦，而他的直言应使你喜悦。因为你被卷入我沉默的危险中，却因我直言的益处而得助。所以，我并非多管闲事，闯入不该去的地方，也不是干涉他人的事务。我是在服从天主的命令。我这样做，首先是出于对你的爱，对你的善意，以及保持你善行的愿望。如果我这话不被相信，或被禁止凭这种情感行事，那么我确确实实是出于害怕得罪天主而说话。因为如果我的危险能解救你，我会耐心地为你献上自己，虽然并非情愿，因为我更愿意你没有我的危险也能蒙天主悦纳并得荣耀。但如果我缄默和掩饰的罪咎既会压垮我，又不能解救你，那么我宁愿你认为我过于急切，而不愿被认为无用和卑鄙。因为经上记载，正如圣宗徒保禄所说，他的教导你们无法否认：\BibleQuote{务要宣讲真道，不论顺境逆境，总要坚持不变；以百般的忍耐和各样的教训去反驳，去斥责，去劝勉。} \BibleRef{弟后 4:2}\par

4. 那么，我们也有了一位更不可得罪的主，尤其是因为即便皇帝也不讨厌各人履行自己的职责，你也耐心听取每个人在其本分内提出建议，不，若有人不按他职务的规矩行事，你还要责备他。那么，司铎们所做的这件事，既然你乐意接受那些侍奉你的人的劝说，难道在司铎身上就显得讨厌吗？因为我们说话，不是凭自己的意愿，而是受命而说。因为你知道这段经文：\BibleQuote{当你们被带到君王和官长面前时，不要思虑怎样说或说什么，因为在那时刻，会赐给你们应说的话；因为说话的不是你们，而是你们父的圣神在你们内说话。} \BibleRef{玛 10:19}\par

5. 我知道你敬畏天主，仁慈、温和、沉稳，心中怀着对天主的信德和敬畏，但往往有些事我们不曾留意。\BibleQuote{他们有些人对天主有热心，但不合乎真知超见。} \BibleRef{罗 10:2} 我认为我们应当小心，免得这事也临到忠信的灵魂身上。我知道你对天主的虔诚，对人的宽厚，我自己也蒙受你恩惠的眷顾。因此我越加恐惧，越加焦虑；生怕连你日后也会凭自己的判断来指责我，因为由于我缺乏坦率或过于奉承，你没有避免某个过失。如果我看到你得罪了我，我不该保持沉默，因为经上记载：\BibleQuote{如果你的弟兄得罪了你，你要趁你和他单独在一起的时候规劝他；如果他不听，你就另外带一两个人，好使任何事情凭两三个见证人的口得以成立。若是他仍不听从他们，你要告诉教会。} 那么，在天主的事上，我能保持沉默吗？因此，让我们想想我所必须恐惧的事。\par

6. 东方军务伯爵报告说，一座犹太人会堂被烧毁，这事是出于主教的指使。你下令惩罚其他人，并让主教本人重建会堂。我不坚称本该先听取主教的陈述，因为司铎本是平息骚乱、渴望和平的人，除非他们自己也被某些冒犯天主或侮辱教会的事所激动。我们姑且假定那位主教在烧毁会堂的事上过于急切，在审判座前又过于胆怯；皇帝啊，你难道不担心，他若遵从你的判决，会背弃他的信德吗？\par

7. 难道你不怕，正如将要发生的那样，他会以拒绝来反对你的伯爵吗？那时伯爵将被迫使他要么成为\OriginalTerm{背教者}{apostate}，要么成为\OriginalTerm{殉道者}{martyr}，这两种情形都不合时宜，都等同于迫害，如果他被迫要么背教，要么忍受殉道。你看得出来，这事的结局倾向于何方。如果你认为主教是坚定的，就要防范使一个坚定的人成为殉道者；如果你认为他动摇不定，就要避免造成一个软弱者的跌倒。因为谁使软弱的人跌倒，这人负有重大的责任。\par

8. 既然我已这样陈述了事情的两面，假设这位主教说，他亲手点燃了火，聚集了群众，召集了百姓，为的是不失去\OriginalTerm{殉道}{martyrdom}的机会，并代替软弱者推出一个更强壮的运动员。啊，幸运的谎言，借此一人为别人获得开释，为自己获得\OriginalTerm{恩宠}{grace}！这就是我，皇帝啊，也曾请求的：愿你宁愿向我报复，并如果你认为这是罪行，就归咎于我。为什么下令审判一个不在场的人呢？有罪之人就在你面前，你听到了他的承认。我声明我放火烧了\OriginalTerm{会堂}{synagogue}，至少是我命令了那些放火的人，为使那里不再有一个否认基督的地方。如果对我提出异议，说我没有在这里放火烧会堂，我回答，它开始被烧是出于天主的审判，我的工作便结束了。如果问及实情，我所以更加懈怠，是因为我未曾料到它会受到惩罚。我为什么要做那件既不受报复便也没有赏报的事呢？这些话虽伤害谦逊，却能挽回恩宠，以免做出冒犯至高者天主的事。\par

9. 然而，且假定没有人会召主教去执行这项任务，因为我已向你的仁慈请求了此事，而且虽然我尚未读到这道法令被撤回，但我们还是假定它已被撤回。如果其他更胆怯的人提议由他们出资重建会堂；或者伯爵发现此事早有定论，便自行下令用基督徒的资金重建它，那又如何？皇帝啊，你将会有一个背教的伯爵，你还要把胜利的军旗托付给他吗？你愿将因基督之名而被祝圣的\OriginalTerm{拉布兰旗}{labarum}，托付给那重建不认识基督的会堂的人吗？下令把拉布兰旗抬进会堂，让我们看看他们会不会反抗。\par

10. 那么，难道要用教会的掠物为犹太人的不信建造一个地方，并将那因基督的恩惠而为基督徒获得的产业，转交给不信者的库房吗？我们读到，古时曾用从辛布里人那里夺来的掠物和其他敌人的掠物为偶像建造殿宇。难道犹太人要在他们会堂的门面上题写这样的字样：\Quote{不虔敬的殿，用基督徒的掠物建造}吗？\par

11. 但是，皇帝啊，或许触动你的是纪律的理由。那么，哪一样更为重要，是纪律的表象，还是宗教的事业？审判必须让位于宗教。\par

12. 皇帝啊，你没有听说过吗，当\Person{尤利安}下令重建\Place{耶路撒冷}圣殿的时候，那些清除瓦砾的人是如何被火烧灭的？你难道不会谨慎，免得这事如今再次发生吗？因为你本不该命令尤利安所命令过的事。\par

13. 但你的动机是什么呢？是因为任何种类的公共建筑被烧毁了，还是因为那是会堂？如果你是因一座无足轻重的建筑被烧而受触动（因为在如此卑微的城镇里能有什么呢？），皇帝啊，你难道不记得，在\Place{罗马}，有多少长官的住宅被烧毁，却无人因此受罚？事实上，若有哪位皇帝想要严厉惩罚此事，他就会损害那遭受如此重大损失之人的事业。那么，哪一样更适宜呢：是\Place{卡利尼库姆}某处建筑的火灾，还是\Place{罗马}城的火灾，该受惩罚呢——如果确实应该惩罚的话？最近在\Place{君士坦丁堡}，主教的住宅被焚毁，你的儿子向他的父亲求情，恳求你不要为施加于他——即皇帝之子——的侮辱，以及主教宅邸的被焚而施行报复。皇帝啊，你难道没有考虑，如果你命令惩罚此事，他会再次干预，反对惩罚吗？然而，儿子从父亲那里获得了这一恩惠是合宜的，因为他首先宽恕加于己身的伤害，是与他相称的。在恩惠的分配上，这是一项好的分工：儿子为自己的损失受恳求，父亲为儿子的损失受恳求。在这里，没有什么需要你为儿子保留的。所以，要谨慎，免得你减损属于天主的任何事物。\par

14. 因此，没有任何充分的理由闹出如此大的骚动，使民众因一座建筑的焚毁而遭受如此严厉的惩罚，更何况这是烧毁一座会堂，一个不信之家，一座不虔的房屋，一个愚昧的收容所，这一切是天主亲自所定罪的。因为我们在上主我们的天主藉\Person{耶肋米亚}先知的口所说的话中读到：\BibleQuote{我必对付这称为我名下的殿，就是你们所信赖的，并对付这地方，就是我赐给你们和你们祖先的，如同我怎样对付了史罗一样。我必从我的面前将你们抛弃，如同我抛弃了你们的兄弟，全体厄弗辣因的后裔。你不要为这百姓祈祷，不要为他们呼吁哀求，也不要向我为他们求情，因为我不会听你。难道你没有看见他们在犹大各城所做的事吗？}\BibleRef{耶 7:14} 天主禁止为那些人代祷。\par

15. 的确，如果我按万民法进行辩护，我能告知在\Person{尤利安}皇帝时代，犹太人烧毁了多少座教会的\OriginalTerm{大殿}{basilica}：在\Place{大马士革}有两座，其中一座至今几乎尚未修复，且是用教会的费用，而不是会堂的费用；另一座大殿仍是一堆不成形的乱石废墟。在\Place{迦萨}、\Place{阿什克伦}、\Place{贝鲁特}以及那些地区的几乎每一个地方，大殿都被烧毁，却无人要求惩罚。在\Place{亚历山大}，一座大殿被外教人和犹太人烧毁，它超过了所有其他的。教会未曾得到报复，会堂就该得到吗？\par

16. 那么，瓦伦廷派庙宇的被焚也应受到报复吗？但那不过是一座庙宇，里面聚集的是异教徒。虽然异教徒呼求十二位神祇，瓦伦廷派却崇拜三十二个他们称为神的\OriginalTerm{永世}{Æons}。我还发现，关于此事，有报告并下令说，一些隐修士当瓦伦廷派堵住道路——他们按习俗和古老惯例，正咏唱圣咏前往庆祝玛加伯节——因愤于他们的无礼，在某个乡村烧毁了瓦伦廷派仓促建造的庙宇，这些隐修士应受惩罚。\par

17. 有多少人不得不面临这样的抉择，当他们想起在\Person{尤里安}时代，那推倒祭台、扰乱祭献的人被法官定罪并遭受殉道？因此，审理他的法官从未被视作别的，只被看作控告者，因为没有人认为配与他交往或给他亲吻。若不是他现在已经死了，皇帝啊，我恐怕你会向他报仇，尽管他未能逃脱上天的报复，比自己的继承人活得更久。\par

18. 但据说，法官受命审理此事，且有书面指示，他不应只报告此事，而应加以惩罚，并应追索被拿走的钱箱。其他事我就不提了。我们的教堂建筑被犹太人焚烧，却没有任何归还，没有任何索回，没有任何要求。但在一个偏远城镇，犹太会堂又能拥有什么呢？那里的一切都不多；没有值钱的东西，也没有富余。那么，诡诈的犹太人又能因火灾损失什么呢？这些都是犹太人的诡计，他们想要诽谤我们，好因他们的抱怨而下令进行特别的军事调查，并派出一名士兵；这名士兵或许会像从前在这里有人曾说过的那样，皇帝啊，在您即位之前说的话：\Quote{基督怎能帮助我们这些为犹太人而与基督争战的人呢？我们被派去为犹太人报仇。他们已毁了自己的军队，还想毁掉我们的军队。}\par

19. 再者，他们既用假见证诽谤过基督，还有什么诽谤的话说不出口呢？这些人在属于天主的事上尚且说谎，还有什么诽谤的话说不出口呢？他们会不说谁是那场骚乱的煽动者吗？他们不会攻击谁呢，甚至那些他们不认识的人，好让他们能观看无数基督徒被锁链捆绑的队列，让他们看到忠信子民的颈项俯首为囚，让天主的仆人被藏在黑暗中，被斩首，被投入火中，被送进矿场，好让他们的苦难不会迅速消逝？\par

20. 难道您要把这战胜天主的教会的胜利送给犹太人吗？皇帝啊，要把这战胜基督子民的战利品、这欢腾，送给不信者吗？把这喜乐送给犹太会堂，把这悲痛送给教会吗？犹太人民将把这庆典列入他们的节日，毫无疑问地将它算入那些他们曾战胜阿摩黎人或客纳罕人，或从埃及王法郎手中被解救，或从巴比伦王拿步高手中被解救的日子之中。他们将增添这个庆典，以纪念他们战胜了基督的子民。\par

21. 他们否认自己受罗马法律的约束，并认为那些法律是有罪的，然而现在他们却认为他们应如同受到罗马法律的报复。当他们自己放火焚烧圣堂的屋顶时，那些法律在哪里？如果\Person{尤里安}因为他是背教者而不为教会报仇，皇帝啊，难道您因为您是基督徒，就要为犹太会堂所受的伤害报仇吗？\par

22. 之后基督将对您说什么呢？您不记得祂藉先知\Person{纳堂}对圣\Person{达味}说的话吗？\BibleQuote{我从你的弟兄中拣选了你最年幼的，并使你由平民成为皇帝。我使你的后裔坐在皇位上。我使蛮族臣服于你，赐予你和平，将你的仇敌交在你手中成为俘虏。你的军队没有粮食供应，我为你打开城门，借仇敌之手为你打开他们的仓库。你的仇敌将他们为自己预备的粮食给了你。我扰乱了你的仇敌的计谋，使他赤露敞露。我如此束缚了帝国篡位者的心思，使他的心智被捆绑，以致虽然他仍有逃脱之路，却连同他所有的一切，好像生怕有人逃脱你一样，将自己关锁起来。他另一部分的军官和军队，我先前将他们驱散，使他们不能联合起来与你争战，我又将他们聚集起来，以完成你的胜利。你的军队，从许多未驯服的民族中聚集而来，我命他们保持信实、安宁与和谐，如同一国。当最大的危险存在，即蛮族的诡诈计谋可能穿透\Place{阿尔卑斯山}时，我就在\Place{阿尔卑斯山}墙之内赐予你胜利，使你毫无损失地得胜。如此，我使你战胜了你的仇敌，而你却使我的仇敌胜过我的子民。}\par

23. 难道不正是由于这个原因，\Person{马克西穆斯}才被抛弃的吗？他在出征之前，听说\Place{罗马}有一座犹太会堂被烧，便向\Place{罗马}发出一道敕令，仿佛他是公共秩序的维护者。因此基督徒说：他不会有好事。那位国王成了犹太人，我们听说他是秩序的维护者，而基督，那位为罪人而死的，很快就考验了他。如果连这样的话都招致如此议论，那么实际的惩罚又将招来什么议论呢？于是，他立刻被法兰克人和撒克逊人击败，在\Place{西西里}、在\Place{西斯基亚}、在\Place{佩塔维奥}，总之，在各地都吃了败仗。信徒与不信者有什么相通呢？他背信的事例理应与那不信者一同被铲除。那使他受害的事，那战败者所犯的过错，征服者不应去效法，而应加以谴责。\par

24. 我重述这些事，并非对一位不知感恩的人，而是列举了那些正确赐予的恩惠，好使你因它们的警示——你既蒙受更多——能爱得更深。当\Person{西满}这样回答时，主耶稣说：\BibleQuote{你判断得正对。}\BibleRef{路 7:43} 随即转身向着那位用香液抹祂脚的女人，以她作为教会的预像，对\Person{西满}说：\BibleQuote{因此，我告诉你：她那许多罪得了赦免，因为她爱的多；但那少得赦免的，是爱的少。}\BibleRef{路 7:47} 这就是那位进入法利塞人家里，抛弃了犹太人，却获得了基督的女人。因为教会已将\OriginalTerm{会堂}{Synagogue}排除在外，为何如今又要试图，在基督的仆人身上，让会堂将教会从信仰的怀抱中，从基督的家中排除出去呢？\par

25. 皇帝啊，我出于对你的爱和热诚，在这篇讲话中汇集了这些事。因着你的仁慈（藉着我的请求，你已将许多人从流放、监牢、死刑的极刑中释放），我应当不惧怕为了你自己的救恩而冒犯你的情绪（没有人比那发自内心去爱的人更有信心，当然也不应伤害那为他着想的人）；我不应在顷刻间失去那赐予每一位司铎、我多年来所领受的恩宠；然而我所恳求避免的，并不是失去恩宠，而是救恩的危险。\par

26. 然而，皇帝啊，这是何等重大的事：你竟认为对于一件迄今无人询问、无人曾经施加惩罚的事，没有必要调查或惩罚。为犹太人而危及你的救恩，这是一件严重的事。当\Person{基德红}杀了那头圣牛时，异教徒说：\Quote{神明自己会为所受的伤害复仇。} 谁来为会堂复仇呢？是基督——他们杀害了祂，他们否认了祂吗？难道天主圣父会为那些不接待圣父的人复仇，因为他们没有接待圣子？谁来为\OriginalTerm{瓦伦提尼安派}{Valentinians}的异端复仇？你的虔诚怎能为他们复仇，既然你已下令将他们排除在外，并禁止他们集会？如果我向你举出\Person{约史雅}，一位为天主所悦纳的君王，难道你会谴责在他们身上所发生的事，就是在约史雅身上所悦纳的事吗？\par

27. 然而，无论如何，如果对我缺乏足够的信任，就请你召集那些你认为合适的主教，让他们来讨论，皇帝啊，在不损害信德的前提下应当做什么。如果你在金钱事务上咨询你的官员，那么在宗教事务上咨询天主的司铎，岂不更为公正？\par

28. 愿你的仁慈思量，教会遭受了多少阴谋者、多少暗探的困扰。他们若发现一丝裂痕，便将毒箭射入其中。我是按人之常理说话，但人对天主的敬畏应超过对人，天主确实被置于君王之上。若有人认为对朋友、父母或邻人应当表示尊重，那么我断定，对天主应当表示尊重，并应将祂置于一切之上，这无疑是正确的。皇帝啊，请你为自己的利益着想，或容我为我自己的利益着想。\par

29. 如果日后发现，因着由这里发出的授权，基督徒被刀剑、棍棒或加铅的鞭子所杀，我将如何回答？我该如何解释这样的事？我又该如何向那些主教们辩解？他们如今正痛苦哀叹，因为一些已履行司铎职务三十年甚至更久的人，或其他教会的圣职人员，被调离他们的圣职，并被派去履行市政职务。因为，那些为你作战的人尚且按一定的服役期服役，那么，你岂不更应考虑那些为天主作战的人？我说，我该如何向那些主教们辩解？他们正为圣职界而抱怨，并写信说教会正因一场严重的攻击而荒废。\par

30. 我曾切愿此事为你的仁慈所知。当你乐意时，请你俯允考虑，并按你的旨意下达命令，却要排除并抛弃那使我忧虑、且使我忧虑得合情合理的事。你个人便是在实行你所命令的一切，即使你的官员未能执行。我非常愿意你显示仁慈，胜过他不执行他所接受的命令。\par

31. 你尚有那些人，为了他们你仍应在罗马帝国面前寻求并赢得主的仁慈；你尚有那些人，你为他们所怀的希望更甚于为你自己；愿天主的恩宠临于他们，愿他们的救恩藉着我的这些话向你呼吁。我担心你会将你的事由交由他人判断。在你面前，一切尚无成见。在这一点上，我为你向我们的天主担保，不要害怕你的誓言。难道为荣耀天主而修正的事，会令天主不悦吗？你无需更改那封信中的任何内容，无论它已发出还是尚未发出。请下令另写一封，一封充满信德、充满虔诚的信。对你而言，改过尚有可能；对我而言，隐瞒真理却不可能。\par

32. 你曾宽恕\Place{安提约基雅}人对你的侮辱；你曾召回你敌人的女儿，并将她们交给一位亲戚抚养；你曾从自己的府库中送钱给你敌人的母亲。如此伟大的虔诚，如此伟大的对天主的信德，却会被这一行为所蒙蔽。因此，我恳求你，你既饶恕了持武器作战的敌人，保全了你的对头，就不要认为基督徒应当如此急切地遭受惩罚。\par

33. 现在，皇帝啊，我恳求你，不要不屑于聆听我——我既为你也为自己恐惧——因为这是一位圣人的话，他说：\BibleQuote{我为什么生在世上，为见我民族的苦难？}\BibleRef{加上 2:7} 致使我得罪天主。我确实已做了在尊重你的前提下所能做的一切，为的是你愿意在王宫里听从我，免得万不得已时，你需要在教会中听从我。\par

\SourceRecord{Translated by H. de Romestin, E. de Romestin and H.T.F. Duckworth.；Nicene and Post-Nicene Fathers, Second Series；Vol. 10.；Edited by Philip Schaff and Henry Wace.；Buffalo, NY: Christian Literature Publishing Co.,；1896.；<http://www.newadvent.org/fathers/340940.htm>.}

% structure: p0001 source=h1 target=section reason=page-title

\section{信函 41}

\EditorialNote{圣盎博罗削在这封写给他姐妹的信中，继续叙述了他写给狄奥多西皇帝的信中所包含的事项，以及他随后在皇帝面前发表的一篇讲道的内容；由于圣盎博罗削在得到皇帝将撤销那令人反感的命令的承诺之前，拒绝举行献祭，最终皇帝做出了让步。}\par

\Signature{弟弟致姐妹}\par

1. 承蒙你写信告知，你的圣德仍感焦虑，因为我在信中曾表示自己也是如此，因此我很惊讶你没有收到我的信，在那信中我已写明我已获得满足。因为当有人报告说，一座犹太人的会堂和一处瓦伦廷派的集会所已被基督徒在主教唆使下烧毁时，在我在\Place{阿奎莱亚}期间，便发出了一道命令：重建会堂，并惩罚那些烧毁瓦伦廷派建筑物的隐修士。后来，因我屡次努力却收效甚微，我便写信并送了一封信给皇帝，当他来到教堂时，我就发表了这篇讲道。\par

2. 在先知书中记载着：\BibleQuote{取一根杏树枝，拿在你手中。} \BibleRef{耶1:11} 我们应当思考主为何对先知说这话，因为这话并非无故写下的，因为在\BibleBook{梅瑟}{五书}中我们也读到，司祭亚郎的杏树枝，在长期存放后，竟开了花。因为主似乎用这树枝象征，先知或司祭的权威应当是正直的，且劝告的不应多是悦耳之言，而应是合宜之事。\par

3. 所以，先知被命取一根杏树枝，因为这树的果实外皮苦涩，内壳坚硬，而里面却是甘甜的，为使他效法这果实，宣讲苦涩和坚硬的事，不畏惧宣讲严厉的话。司祭也是如此；因为他的教训，虽然一时在有些人看来是苦涩的，并像亚郎的树枝一样，长久存放在那些伪善者的耳中，然而过了一段时间，当人以为它已枯干时，它却开了花。\par

4. 为此，宗徒也说：\BibleQuote{你们愿意怎样呢？愿意我带着棍棒到你们那里去呢？还是怀着慈爱和温柔的心情呢？} \BibleRef{格前4:21} 他先提到了棍棒，像杏树枝一样击打那些走入歧途的人，为以后能以温柔的心安慰他们。于是，温柔使那曾被棍棒剥夺了天上圣事的人得以恢复。他也给自己的门徒类似的嘱咐，说：\BibleQuote{责备、劝勉、斥责。} \BibleRef{格后2:10} 这三者中有两者是严厉的，一者是温和的，但这严厉是为了能使之柔和；就如对于胆汁过盛的人来说，苦涩的食物或饮品仿佛是甜的，反之甜的食物却是苦的，同样，当灵魂受伤时，在谄媚的愉悦下会恶化，而在纠正的苦涩中却能得以治愈。\par

5. 这些是从先知经文中所汲取的，现在让我们思考福音的读经所包含的内容：\BibleQuote{有一个法利塞人请主耶稣同他吃饭，他便进了那法利塞人的家，在席上坐下。看，有个女人，是城中的罪妇，她听说耶稣在那法利塞人家里吃饭，就带来一玉瓶香液，站在他背后，靠近他的脚，开始用眼泪洗他的脚。} 然后他读到这个地方：\BibleQuote{你的信德救了你，平安回去吧！} 说得多简单啊，我接着说，这福音的教训在文字上多么简单，在计划上又是何等深奥！正因为这些话出自 \BibleQuote{伟大的谋士} \BibleRef{依9:6} 之口，让我们思考它们的深度。\par

6. 我们的主耶稣基督判断，人们更容易通过仁爱而非恐惧被束缚和引导去行正义之事，爱比畏惧更能促使人改正。因此，当他由童贞女诞生而来临时，他施予了他的恩宠，使罪在圣洗中得赦，好让我们对他更加感激。那么，如果我们以感恩者所当有的服务来回报他，他已在这女人身上表明，这恩宠本身对所有的人都将获得赏报。因为如果他只赦免了我们原罪的债，他可能显得更为谨慎而非仁慈，更注重纠正我们，而非慷慨施赏。想要诱引人的，只是狭隘之心的狡计；但天主所当有的，却是以恩宠召唤了人，便以那恩宠的增长来引导他们。因此，他先藉圣洗赐予我们恩惠，然后再对那些忠心侍奉他的人更丰盛地施予。这样，基督的恩惠既是德行的激励，也是德行的赏报。\par

7. 谁也不要因 \BibleQuote{债主} \BibleRef{路7:41} 这个词而惊慌。我们从前处在一个严厉的债主之下，他除非使负债者死亡，否则无法得到满足和全额的偿还。主耶稣来了，他看见我们被沉重的债务所束缚。没有人能以自己的纯洁遗产偿还债务。我没有自己的东西可以赎身。他给了我一种新的清偿方式，更换了我的债主，因为我无力偿还我的债务。但是，使我们成为负债者的，是罪，而非本性，因为我们因罪恶而背负了沉重的债务，使我们这些原为自由的人成为奴隶；因为一个负债者就是从债主那里接受了某种财物的人。而罪是属于魔鬼的；那邪恶者可以说拥有这些财富。正如基督的财富是德行，魔鬼的财富就是罪恶。他由继承而来的沉重债务，将人类置于永久的囚禁中，这债务是我们的祖先身负重债，以遗产的方式传给后代的。主耶稣来了，他为众人死，以自己的死亡代替了众人的死亡，倾流了他的血以代替众人的血。\par

8. 因此，我们更换了债主，但并非完全逃脱，不，或者说我们逃脱了，因为债务虽存留，但利息已被取消，因为主耶稣说过：\BibleQuote{向被囚的人说：出来吧！向在黑暗中的人说：露面吧！} \BibleRef{依49:9} 你们的罪就这样被赦免了。所有的人都被赦免了，没有一个人他没有释放的。因为经上这样记载，他赦免了 \BibleQuote{一切过犯，涂抹了那反对我们，告发我们对戒命负责的债卷。} \BibleRef{哥2:13} 那么，我们为什么还紧紧抓住别人的债，想要追索别人的债，而我们自己却享受宽免呢？那位宽恕了所有人的，要求每个人也应宽恕别人，因为每个人记得自己也曾蒙受了宽恕。\par

9. 要注意，你作为债权人的处境比作为债务人时更糟，就像福音中的人，他的主人赦免了他所有的债，但后来他却向自己的同伴索要他还没有偿还的债，他的主人因此发怒，以最严厉的斥责向他索要之前已被赦免的债。因此，我们要警惕，不要因为不宽恕别人欠我们的债，而导致自己必须偿还被赦免的债，因为这是照着主耶稣的话所记的：\BibleQuote{如果你们不各自从心里宽恕自己的弟兄，我的天父也要这样对待你们。}\BibleRef{玛 18:35} 那么，我们既然蒙受了宽赦，就当宽恕那些欠我们微不足道的人，并且明白我们越宽恕别人，就越蒙天主悦纳，因为我们蒙受的赦免越多，就越得享天主的喜悦。\par

10. 最后，当主问法利塞人：\BibleQuote{他们哪一个更爱他呢？}\BibleRef{路 7:42} 他回答说：\BibleQuote{我想是那多得赦免的人。}主对他说：\BibleQuote{你判断得正确。}\BibleRef{路 7:43} 法利塞人的判断受到了称赞，但他的情感却受到了责备。他对他人的判断很好，但自己却不相信自己在判断他人时所认为好的事。你听见一个犹太人称赞教会的纪律，颂扬它的真实恩宠，尊敬教会的司铎；但如果你劝他相信，他却拒绝，因此他自己并不遵行他在我们身上所称赞的。因此，他的称赞并不完全，因为基督对他说：\BibleQuote{你判断得正确，}因为加音也曾正确地献祭，却没有正确地分开，所以天主对他说：\BibleQuote{你若献得对，但分得不对，你就犯了罪，要安静。}所以，这个人献得对，因为他判断基督应当更受基督徒的爱戴，因为基督赦免了我们许多的罪；但他分得不对，因为他以为那赦免人罪的主可能不知道人的罪。\par

11. 因此，他对西满说：\BibleQuote{你看这女人。我进入你的家，你没有给我水洗脚，她却用眼泪洗了我的脚。}\BibleRef{路 7:44} 我们所有人都是基督同一的身体，头是天主，我们是肢体；有些或许是眼，如先知；有些是牙，如宗徒们，他们把所传的福音之粮送入我们的胸中，所以经上正确地写着：\BibleQuote{他的眼因酒而发红，他的牙因乳而雪白。}\BibleRef{创 49:12} 他的手是那些施行善工的人，他的腹是那些分施养料力量给穷人的人。同样，有些人是他的脚，但愿我配作他的脚跟！那么，浇在基督脚上的水，就是宽恕那些最卑微者的过犯，并且在拯救卑微者的同时，是在洗基督的脚。\par

12. 浇在基督脚上的水，是人净化自己的良心，使之脱离罪污，因为基督行走在每个人的心中。因此，要当心，不要让你们的良心被污染，从而开始玷污基督的脚。要当心，以免他在你们内遇见罪恶的荆棘，从而当他行走在你们内时刺伤他的脚跟。这正是为什么法利塞人没有给基督倒水洗脚的原因，因为他没有一颗从不信的污秽中清洁的灵魂。因为一个没有接受基督之水的人，怎能洁净自己的良心呢？但是教会既有这水，也有眼泪。因为痛悔先前罪过的信德，习惯于防范新的罪过。因此，那个没有水的法利塞人西满，当然也没有眼泪。因为没有悔改的心，他怎么会有眼泪呢？因为他既然不信基督，就没有眼泪。因为如果他有眼泪，他早就洗净了自己的眼睛，以便能够看见基督，而他虽然与他同席，却看不见他。因为如果他看见了他，就不会怀疑他的能力了。\par

13. 法利塞人没有头发，因为他不能认出拿齐尔人\OriginalTerm{拿齐尔人}{Nazarite}；教会有头发，并且她寻求拿齐尔人。头发被算作身体的多余之物，但如果它们被油膏抹，就会发出馨香，成为头上的装饰；如果不被油膏抹，便成为重担。同样，财富也是重担，如果你不知道如何使用它们，不把基督的馨香洒在它们上面。但如果你养育穷人，如果你清洗他们的伤口，擦去他们的污秽，你就确实洗了基督的脚。\par

14. \BibleQuote{你没有给我亲嘴，但她自从进来，就不停地亲我的脚。}\BibleRef{路 7:45} 亲嘴是爱的标记。那么，一个犹太人怎能有亲嘴呢？因为他既没有认识和平，也没有从基督那里接受和平，当基督说：\BibleQuote{我把我的平安赐给你们，我将我的平安留给你们。}\BibleRef{若 14:27} 会堂\OriginalTerm{会堂}{Synagogue}没有亲嘴，但教会有，她等候他，爱慕他，她说：\BibleQuote{愿他用自己的口与我亲嘴。}\BibleRef{歌 1:2} 因为她想用他的亲嘴逐渐熄灭那长久渴望的燃烧，这渴望随着期待主的来临而增长，以此恩赐来解她的渴。因此，圣先知说：\BibleQuote{你要开启我的口，我的口要宣扬你的赞美。} 那么，赞美主耶稣的人就是在亲他，赞美他的人无疑就相信。最后，达味自己也说：\BibleQuote{我信，所以我说；} 以及之前：\BibleQuote{愿我的口充满你的赞美，让我歌颂你的光荣。}\par

15. \OriginalTerm{特殊恩宠}{special grace}的注入，圣经教导你，凡领受圣神的人就是亲吻基督，正如圣先知所说：\Quote{我张开口吸入圣神。}那么，凡承认基督的人就是亲吻祂：\Quote{因为心里相信，可使人成义；口里承认，可使人获得救恩。}\BibleRef{罗 10:10}再者，凡在阅读福音时认出了主耶稣的作为，并以虔敬的情感钦佩它们的人，就是亲吻基督的双足；如此，他以虔敬之情亲吻了主耶稣行走的足迹。那么，我们是以\OriginalTerm{共融之吻}{kiss of communion}亲吻基督：\Quote{读的人应领悟。}\BibleRef{玛 24:15}\par

16. 犹太人从何能有此亲吻呢？因为那相信祂来临的人，并不相信祂的受难。人若不相信祂已来临，怎能相信祂受了苦呢？所以，法利塞人没有亲吻，除非也许有叛徒犹达斯的那种亲吻。然而犹达斯也没有这亲吻；因此，当他想要向犹太人表示他曾应许作为出卖标记的亲吻时，主对他说：\Quote{犹达斯，你用亲吻来出卖人子吗？}\BibleRef{路 22:48}这就是说，你这没有亲吻所标记的爱的人，却献上亲吻。你这不懂亲吻奥迹的人却献上亲吻。所寻求的不是嘴唇的亲吻，而是心灵与灵魂的亲吻。\par

17. 但是你说，他亲吻了主。是的，他确实用嘴唇亲吻了祂。犹太民族有这样的亲吻，因此经上说：\Quote{这民族用嘴唇尊敬我，心却远离我。}\BibleRef{玛 15:8}所以，凡没有信德和爱德的人便没有这亲吻，因为亲吻铭刻着爱的力量。没有爱时，没有信德，没有亲情，亲吻中能有什麽甜蜜呢？\par

18. 然而教会并不停止亲吻基督的双足，因此，她在\BibleBook{雅}{歌}中渴望不只一个亲吻，而是许多亲吻\BibleRef{歌 1:2}；并且像圣母玛利亚一样专心于祂的一切言语，当诵读\WorkTitle{福音}或\WorkTitle{先知书}时，她接受祂的一切话，并\Quote{把这一切事默存在自己心中。}\BibleRef{路 2:51}所以，惟有教会像新娘一样拥有亲吻，因为亲吻有如婚约的保证和婚姻的特权。犹太人从何能有亲吻呢？他们不信新郎；犹太人从何能有亲吻呢？他们不知道新郎已经来临。\par

19. 而且，他不但没有亲吻，也没有用来傅抹基督双足的油，因为如果他有油，他必早已软化自己的颈项。\par

梅瑟说：\Quote{这百姓是执拗的百姓，}\BibleRef{出 34:9}主曾说司铎和肋未人从旁边走过去，他们两人都没有向被强盗打伤的人伤口注油或酒；\BibleRef{路 10:31}因为他们没有什么可注的，因为如果他们有油，他们必注在自己的伤口上了。而依撒意亚宣告说：\Quote{他们不能敷药膏，不能注油，也不能缠绷带。}\BibleRef{依 1:6}\par

20. 然而教会有油，用来包扎她子女的伤口，免得伤口硬化蔓延；她拥有秘密领受的油。正如经上记载，阿协尔用这油洗了脚：\Quote{阿协尔是受祝福的儿子，他必蒙兄弟们悦纳，必将脚浸在油里。}\BibleRef{申 33:24}所以，教会用这油傅抹她子女的颈项，使他们能负起基督的轭；她用这油傅抹殉道者，使他们洁净脱离世俗的尘土；她用这油傅抹精修者，使他们不屈服于劳苦，不因疲倦而沉沦，不被世俗的灼热所胜；她傅抹他们，是为以\OriginalTerm{属神的油}{spiritual oil}使之恢复精力。\par

21. \OriginalTerm{会堂}{Synagogue}却没有这油，因为她没有橄榄树，也不明白洪水后那叼回橄榄枝的鸽子。\BibleRef{创 8:11}因为那鸽子后来在基督受洗时降下，停留在祂身上，正如\BibleBook{若望}{福音}所证实的，说：\Quote{我看见圣神像鸽子从天降下，停在祂身上。}\BibleRef{若 1:32}然而，那看不见圣神像鸽子降在其上的那一位的，又怎能看见鸽子呢？\par

22. 所以，教会既洗基督的双足，又用自己的头发擦干，又用油傅抹，又倾注香液；因为她不仅照顾受伤者，安慰疲惫者，而且也以恩宠的芬芳洒在他们身上；她将这同一恩宠不仅倾注于富贵和有权势的人，也倾注于卑微的人。她用同一天平衡量所有人，将所有人聚集在同一怀抱中，在同一膝上眷顾他们。\par

23. 基督曾一次而死，一次而被埋葬，然而祂愿意每日有人在祂的脚上倾注香液。那么，我们倾注香液的那双基督的脚是什么呢？基督的双足就是祂自己所说的那些：\Quote{你们对我这些最小兄弟中的一个所做的，就是对我做了。}\BibleRef{玛 25:40}福音中的那位妇女所滋润的，就是这双脚；她以眼泪浸湿它们；当罪过向最小的那一位得到宽恕时，罪债便被洗除，宽恕便赐下了。那爱护圣洁子民中最小者的人，就是亲吻这双脚。那将自己温和的慈爱甚至施予软弱者的人，就是在用香液傅抹这双脚。在这些最小的人身上，殉道者，宗徒，以及主耶稣自己都宣称受到了尊敬。\par

24. 你可以看到，主是多么乐于教导，祂用自己的榜样激发你的虔诚，因为祂在谴责时也乐于教导。所以，当祂指责犹太人说：\BibleQuote{我的人民！我对你做了什么？我在什么事上叫你厌烦？我在什么事上使你疲乏？请答复我！难道是因为我由埃及地领了你上来，从奴隶之家救出了你？} 又加上：\BibleQuote{我在你面前派遣了\Person{梅瑟}、\Person{亚郎}和\Person{米黎盎}。} 你应记得\Person{巴郎}企图陷害你的事，\BibleRef{户 23:2}，他寻求巫术的帮助，但我没有容他伤害你。你们确曾在异域充军受压，背负沉重负担。我派遣了梅瑟、亚郎和米黎盎在你们面前，而那劫掠充军者的人，自己先受了劫掠。你们失去了自己的所有，却获得了别人的财产；你们脱离了围困你们的敌人，在海水中间安全地行走，看见敌人被消灭，因为那环绕并载运你们前行的波浪倒流回来，淹没了敌人。\BibleRef{出 14:29} 当你们经过旷野缺少食物时，我岂不是给你们降下食物之雨，并在你们所到之处供给你们食粮吗？我岂不是在征服了你们的一切敌人之后，领你们进入了\Place{厄市苛耳}地区吗？\BibleRef{户 13:24} 我岂不是将阿摩黎人的王\Person{息红}，\BibleRef{户 21:24}（即那骄傲者，激励你们的人的首领），交给了你们吗？我岂不是将那\Place{哈依}王活活地交给了你们，你们按古时的诅咒，将他悬在木架上，并高举在十字架上？我还需要提起那五个王的军队吗？他们因试图拒绝将赐给你们的土地交给你们而被杀。如今，在这一切之后，所要求于你们的，无非是履行公义，爱好慈仁，虚心与你的天主同行。\BibleRef{米 6:8}\par

25. 那位虔诚而又温和的人，先知\Person{纳堂}对\Person{达味}王所发出的谴责又是怎样的呢？祂说：\BibleQuote{我在你兄弟中拣选了你最年幼的，以温和的精神充满了你，借\Person{撒慕尔}的手给你傅油，立你为王，而我和我的名也居住在他内。我除去了那先前的王，就是那被恶神激动去迫害主的司祭的，使你在流亡之后得胜。我使你的后裔继承你的宝座，他不仅是继承人，更是同侪。我甚至使外方人屈服于你，使那些攻击你的人服事你；而你竟将我的仆人交于我的敌人手中，夺去我仆人的所有，以致你自己会因罪恶而受耻辱，我的仇敌也有欢庆的因由。}\par

26. 因此，皇帝陛下，我现在不但论及你，而且也直接向你讲话：既然你看到主习惯于多么严厉地谴责，那么，你越是荣耀，就应越发彻底地服从你的造主。因为经上记载说：\BibleQuote{当主你的天主领你进入异域，你要吃别人的果实时，切不可说：这是我的能力和我的正义给了我这一切，因为这是主你的天主赐给你的；} 因为这是基督按祂的慈爱赐给你的，因此，你要为了对祂的身体——即教会的爱，给祂的脚浇水，亲吻祂的脚，这样你不仅宽恕那些陷于罪恶的人，而且还要以你的和平使他们恢复和睦，赐给他们安息。你把香液倒在祂的脚上，使基督所坐的整个房屋充满你的香液，所有与祂同坐的人都因你的芬芳而喜乐，就是说，你要尊重卑微的人，使天使为他们的宽恕而欢乐，如同为一个罪人悔改而欢乐一样，\BibleRef{路 15:10}宗徒们喜悦，先知们欢欣。因为眼睛不能对手说：\BibleQuote{我不需要你，} 头也不能对脚说：\BibleQuote{你对我毫无用处。}\BibleRef{格前 12:21}因此，既然众肢体都是必要的，你就要保护主耶稣的整个身体，好使祂也因祂天上的俯就而保全你的王国。\par

27. 当我从讲道台下来时，他对我说：\Quote{你讲论了我。} 我回答说：\Quote{我讨论的是为你有益的事。} 然后他说：\Quote{关于主教修缮会堂的事，我确实决定得过于严厉，但那事已经纠正了。隐修士们犯下许多罪行。} 于是\Person{提马修斯}将军开始猛烈攻击隐修士们，我回答他说：\Quote{我与皇帝来往，是合宜的，因为我知道他敬畏天主，但对你这样说话粗暴的人，就必须另当别论。}\par

28. 然后，我站立了一会儿，就对皇帝说：\Quote{请让我无忧地为你献祭，使我的心安定。} 他继续坐着点头，但没有公开承诺，我仍站立着，他说他会修正诏令。我立即接着说，他必须停止整个调查，免得伯爵利用调查的机会伤害基督徒。他应允了。我对他说：\Quote{我凭你的承诺行事，} 又重复说：\Quote{我凭你的承诺行事。} 他说：\Quote{凭我的承诺行事吧。} 于是我便走向祭台，如果他没有给我明确的承诺，我是不会走到那里去的。的确，伴随这祭献的恩宠如此之大，我自己感觉到，皇帝所赐的那恩惠为我们的天主极为悦纳，天主的临在毫不欠缺。于是，一切都照我所愿地成就了。\par

\SourceRecord{Translated by H. de Romestin, E. de Romestin and H.T.F. Duckworth.；Nicene and Post-Nicene Fathers, Second Series；Vol. 10.；Edited by Philip Schaff and Henry Wace.；Buffalo, NY: Christian Literature Publishing Co.,；1896.；<http://www.newadvent.org/fathers/340941.htm>.}

% structure: p0001 source=h1 target=section reason=page-title

\section{书信五十一}

\EditorialNote{致皇帝\Person{狄奥多西}，于\Place{得撒洛尼}屠杀之后。圣\Person{安博}首先说明了他未在皇帝返回\Place{米兰}时与其会面的理由。接着他提及主教们对于\Place{得撒洛尼}屠杀的看法，并指出必须为此事悔改，以求得赦免并战胜那煽动此罪的魔鬼。圣\Person{安博}不能在皇帝面前举行祭献，作为真诚爱戴皇帝之人，他悲伤却仍怀希望。}\par

我愉快地回忆起你我昔日的友谊，并衷心感激你因我的屡次代祷而极其仁慈地赐予他人的恩惠。由此可见，我并非出于忘恩负义，才在你到来时回避与你相见——对此相见，我以往总是迫切期盼的。现在我将简要说明我如此行事的原因。\par

我发觉在您的朝廷中，唯独我被剥夺了聆听的自然权利，以至于我也被剥夺了发言的职分；因为您时常因某些在您的御前会议中决定的事情传到了我耳中而烦恼。因此，我无法享有这一共同权利，因为主耶稣说：\BibleQuote{因为隐藏的事没有不显露出来的。} \BibleRef{路 8:17} 因此，我极尽恭敬地遵从了皇帝的旨意，注意做到既不让您有任何不悦的理由——因为我想方设法不让任何关于皇帝敕令的事传给我；也不让我自己——即使在场——由于对众人的畏惧而听不到，从而落得纵容之嫌；或者即使听到，也要以这样的方式：我的耳朵虽敞开，我的口舌却被封住，使我不可能说出我所听到的，以免伤害和危及那些被疑为泄露机密的人。\par

那么，我能做什么呢？我难道不应该听吗？但我不能用古老传说中的蜡封住我的耳朵。我难道应该说出我所听到的吗？但我在言语上不得不小心提防您命令中我所害怕的事——即可能发生流血事件。我难道应该保持沉默吗？但那样我的良心就会被束缚，我的话语权被剥夺，这将是最可悲的境况。那么，那经文又当如何理解呢？\BibleQuote{如果司铎不警诫犯错的人，那犯错的人必将死于自己的罪恶中，而司铎也因未警诫犯错者而须受罚。} \BibleRef{则 3:18}\par

请听我说，尊贵的皇帝。我不能否认您对信仰的热忱；我确实承认您有敬畏天主的心。但是您天生性情急躁，如果有人设法安抚，您很快就会转为仁慈；如果有人煽动，您就会越发激怒，几乎难以自制。但愿即使无人安抚，至少也不要有人去煽动它！我情愿将这性情托付给您自己，您会自我克制，并以虔敬之爱来克服这急躁的天性。\par

我宁愿私下将您的这种急躁托付给您自己斟酌，而不想可能因我在公开场合的任何举动而激起它。因此，我宁可显得有些失职，也不愿失礼；宁可让人认为我缺少司铎的权威，也不愿让您觉得我缺乏最敬爱的尊敬，以免您的急躁被激怒后，您的决策能力受到削弱。我以身体患病为借口，这病确实很重，除非悉心调养，几乎难以减轻。然而，我宁死也不愿不等两三天迎接您的到来。但我实在做不到。\par

在\Place{得撒洛尼}人的城里，发生了一件史无前例的事，是我无力阻止的；其实，我先前多次为此求情时，就曾说过那将是一桩极其残暴的事，而您自己后来虽为时已晚地撤销了命令，也表明您认为此事严重——事情既已发生，我也无法减轻其罪过。消息初传时，正逢因高卢主教们的到来而召开了一次公会议。没有一个人不为此哀叹，没有一个人不认为事态严重；您与\Person{安博}的友谊也不能为您的行为开脱。若没有一个人说您必须与我们的天主修和，那么对于已发生之事的谴责，必将越来越多地落到我的头上。\par

皇帝啊，难道您羞于效法那按血统作基督祖先的先知\Person{达味}王所做的事吗？有人告诉\Person{达味}，一个拥有许多羊群的富人，因客人到来，竟夺走并宰杀了穷人的唯一一只小母羊；\Person{达味}明白了那个比喻正是指责他自己，因为他正是做了这事，他便说：\BibleQuote{我得罪了主。} 那么，皇帝啊，不可急躁，如果对您说：\BibleQuote{您做了先知曾对\Person{达味}王所说的事。} 如果您顺从地听这话，并说：\BibleQuote{我得罪了主，} 如果您重复先知\Person{达味}王的话：\BibleQuote{请大家前来，我们要屈身朝拜，在创造我们的主面前，悲哀痛哭，} 那么，也必对您说：\BibleQuote{既然你悔改了，主就除去你的罪，你必不死。}\par

再者，\Person{达味}在命令统计百姓之后，心中感到愧疚，对主说：\BibleQuote{我犯了重罪，因为我这样命令了；现在，主，求你除去你仆人的罪孽，因为我实在做了糊涂事。} 于是先知\Person{纳堂}又奉命来见\Person{达味}，给他三个选择，任他挑选其一：国内遭遇三年的饥荒；或是在敌人面前逃难三个月；或是国内发生三天的致命瘟疫。\Person{达味}回答说：\BibleQuote{这三者都使我极为难堪，但我宁愿落在主的手里，因为他的仁慈极多，而不愿落在人的手里。} 他的过错在于他想知道自己手下全体百姓的数目，而这种知识他本应只留给天主。\par

9. 并且，有人告诉我们，当死亡临到百姓，就在当天正午时分，\Person{达味}看见天使击杀百姓，便说：\BibleQuote{我犯了罪，身为牧者的我行了邪恶，但这羊群作了什么？愿祢的手攻击我和我的父家。} 于是主后悔了，就命令天使宽赦百姓，并要\Person{达味}奉献\OriginalTerm{祭献}{sacrifice}，因为那时的人以祭献赎罪；而如今的祭献则是\OriginalTerm{补赎}{penitence}。他就这样藉着谦卑自下，成了天主更悦纳的人，因为人犯罪原不足为奇，可责的是，若不承认己过，并在天主前自谦自卑。\par

10. \Person{圣约伯}，在世上也是有权势的，说：\BibleQuote{我并没有遮掩我的过犯，而是向全体百姓陈明了。} 他的儿子\Person{约纳堂}对那位暴戾的\Person{撒乌耳}王说：\BibleQuote{不要对你的仆人达味犯罪；} 又说：\BibleQuote{你为什么要犯罪，流无辜人的血，无缘无故杀害达味呢？} 因为，\Person{撒乌耳}虽是君王，若他杀害了无辜者，他也就犯了罪。再者，\Person{达味}掌国之后，听闻无辜的\Person{阿贝乃尔}被他的军长\Person{约阿布}所杀，便说：\BibleQuote{我和我的国，对\Person{乃尔}的儿子\Person{阿贝乃尔}的血，从今时直到永远，都是无罪的；} 他并且哀伤禁食。\par

11. 我写这些，不是为叫你羞愧，而是愿这些君王的榜样能激励你，将这罪从你的国中除掉；因为你若在天主前谦卑你的灵魂，就能除掉它。你是一个人，这罪临到了你，就要战胜它。罪过非赖眼泪与补赎不能除尽。连天使也不能，总领天使也不能。主自己，唯有祂能说：\BibleQuote{我同你们天天在一起，} \BibleRef{玛 28:20} 我们若犯了罪，祂只赦免那悔改的人。\par

12. 我恳切地劝勉、哀求、叮咛、警告，因为令我悲伤的是：你曾是异乎寻常的虔诚表率，以仁慈著称，连单个罪犯你都不忍使之陷入危险，如今竟不为这许多人的丧亡哀恸。虽然你作战所向披靡，虽然你在其他事上也值得称扬，但虔诚始终是你行事最华美的冠冕。魔鬼嫉妒你这项最卓越的禀赋。你既还拥有能战胜他的力量，就去战胜他吧。不要再因那已经伤害多人的行为，在你的罪上增添另一项罪了。\par

13. 我诚然亏欠你的仁慈，对此我不能不感念；你的仁慈超越了许多皇帝，仅有一个人可比。我，我说，我没有理由被指控违抗你，却有理由畏惧；倘若你打算亲临，我不敢祭献。流了一个无辜人的血之后尚不可行的事，难道流了许多人的血之后就可以行吗？我不以为然。\par

14. 最后，我亲笔写下这些，只让你一人阅读。正如我希望主将救我脱离一切困境，我受到了警告，并非由人，也非透过人，而是主亲自清楚地明示，这事对我是不许可的。因为我焦虑不安，就在我准备动身的那一夜，你在梦中显现给我，像是进了圣堂，而我竟不被允许祭献。我略过其他我原可避开的事，但我因爱你的缘故忍受了，我如此相信。愿主使万事平稳地过去。我们的天主以多种方式给予警告，藉天上的征兆，藉先知的训令；甚至藉罪人的异象，祂愿意我们领悟，我们应当恳求祂消除一切纷扰，为你们皇帝保持和平，好使教会的信仰与和平得以延续；皇帝身为基督徒且虔诚，正有益于此。\par

15. 你必定切望蒙天主悦纳。\BibleQuote{事事有时节，} \BibleRef{训 3:1} 经上又载：\BibleQuote{主，现在是祢行动的时刻。} \BibleQuote{上主，这是悦纳的时候。} 当你获准祭献，而你的奉献能为天主悦纳时，你才可以献上。若情况允许，能得皇帝恩宠，照你的心意行事，难道不令我欢欣吗？祈祷本身即是祭献，它能求得赦免，当祭品反致冒犯时；因为前者是谦卑的表记，后者是轻慢的表示。因为\OriginalTerm{圣言}{Word of God}亲自告诉我们，祂喜爱人们遵行祂的诫命，胜于奉献祭献。天主这样昭示，\Person{梅瑟}如此向百姓宣告，\Person{保禄}这样向外邦人宣讲。你当行你认为对此时最有益的事。经上说：\BibleQuote{我喜爱仁慈，胜过祭献。} \BibleRef{玛 9:13} 那么，难道那些谴责自己罪恶的人，不比那些想为自己的罪辩护的人，更配称为真正的基督徒吗？\BibleQuote{义人首先控告自己。} 人若犯了罪而控告自己，才是义人，而不是那自夸的人。\par

16. 皇帝啊，我愿自己在此之前信赖的是我自己，而不是你的习性。当我想到你常常迅速宽恕，迅速收回判决；当前我已先发制人，也就没有逃避那原不必畏惧的事。然而感谢主，祂愿意惩戒祂的仆人，为使他们不致丧亡。这一点我与先知们相通，而你必将与圣人们相通。\par

17. 难道我对\Person{格拉齐安}的父亲的爱惜，不胜过我的眼珠吗？你其余神圣的亲人也都要求宽恕。我预先已将美名赠予那些我以共同的爱所拥抱的人。我以我的爱、我的情感和我的祈祷追随你。你若相信我，就听从我的引导；我再说，你若相信我，就承认我所说的；你若不信我，就请原谅我所行的，因我将天主置于你之前。至尊的皇帝，愿你及你圣善的后裔，享永久的平安，满备幸福与昌盛。\par

\SourceRecord{Translated by H. de Romestin, E. de Romestin and H.T.F. Duckworth.；Nicene and Post-Nicene Fathers, Second Series；Vol. 10.；Edited by Philip Schaff and Henry Wace.；Buffalo, NY: Christian Literature Publishing Co.,；1896.；<http://www.newadvent.org/fathers/340951.htm>.}

% structure: p0001 source=h1 target=section reason=page-title

\section{第57封信}

\EditorialNote{圣\Person{安波罗修}告知\Person{欧根尼乌斯}皇帝他为何不在\Place{米兰}，然后继续就皇帝在异教崇拜方面的行为进行责备。他说，这就是他没有早些写信的原因，并承诺今后他将像对待其他皇帝一样自由地对待他。}\par

\Person{安波罗修}主教致最仁慈的皇帝\Person{欧根尼乌斯}。\par

1. 我离去的原因是敬畏主，我习惯于将我的一切行为都归向祂，尽我所能，从不让我的心偏离祂，也不看重任何人胜过\Person{基督}的恩宠。因为我不伤害任何人，如果我把天主置于一切之上，并信赖祂，我就不怕把像我这样的想法告诉你们这些皇帝。所以，皇帝啊，对于我在其他皇帝面前没有保持沉默的事情，在您面前我也不会保持沉默。为了保持事情的条理，我将逐一提及与此事相关的事项。\par

2. 显赫的\Person{西马库斯}在担任城市执政官时，曾上书可敬的、年少时的\Person{瓦伦提尼安}皇帝，请求他下令将被夺走的东西归还给各庙宇。他出于他的热忱和他的宗教，履行了他的职责。而作为主教，我也理应认识到自己的职责。我向皇帝们呈交了两份请愿书，其中我指出，一个基督徒不能为祭献费用出资；我确实不是废除这些祭献的原因，但我当然极力主张不应下令恢复它们；最后，他本人看起来似乎是在将这些款项给予那些偶像，而不是归还它们。因为不是他自己取走的东西，他不能算是归还，而是出于自己的意愿给予迷信的开支。最后，如果他这样做，要么他就不能来教会，要么如果他来了，他要么在那里找不到司铎，要么他会找到一位在教会中反对他的司铎。也不能以他是\OriginalTerm{慕道者}{catechumen}为借口，因为也不允许慕道者为偶像的开支出资。\par

3. 我的信件在御前会议中被宣读。享有最高军事权威的\Person{鲍托}伯爵在场，还有\Person{鲁莫里杜斯}，他自己也享有同等尊位，从少年时代起就沉迷于外邦民族的崇拜。那时\Person{瓦伦提尼安}听从了我的建议，除了我们信仰的规则所要求的，他什么也没做。于是他们就服膺于他的军官。\par

4. 后来，我坦率地进言于最仁慈的\Person{狄奥多西}皇帝，毫不犹豫地当面进言。他收到了来自元老院的类似信息，尽管这不是整个元老院的请求，最终还是同意了我的建议，因此我有几天没有去见他，他也没有见怪，因为他知道我不是为自己的益处行事，而是不耻于在君王面前说出那对他和我自己的灵魂都有益的事。\par

5. 元老院再次派使团前往\Place{高卢}，觐见可敬的\Person{瓦伦提尼安}皇帝，但一无所获；那时我当然不在场，当时也没有给他写过任何东西。\par

6. 然而，当您的仁慈接掌政府权柄之后，后来人们发现，此类恩惠竟被赐予那些在国家事务上确实出色、但在宗教上却是异教徒的人。或许，奥古斯都皇帝，可以说您并没有向庙宇作出任何归还，而是向那些对您有功的人赠送了礼物。但您知道，我们必须为了天主的事业不断行动，正如为了自由的事业常做的那样，不仅由司铎来做，也由那些在您军队中的人，或被视为居住在各省的人们来做。当您成为皇帝时，有使者请求您向庙宇作出归还，而您并没有这样做；其他人第二次来，您拒绝了，后来您认为应将此恩惠赐予那些提出请求的同一批人。\par

7. 虽然皇权是伟大的，但请想一想，皇帝啊，天主是何其伟大。祂洞悉所有人的内心，祂审察最深处之良心，祂在万事发生之前便已知道一切，祂知道你胸中最隐秘之事。你们不容自己受欺骗，难道你们还想向天主隐瞒什么吗？难道这没有进入你的心思吗？虽然那些人的行动如此坚持不懈，但皇帝啊，你难道不应该出于对至高的、真实且生活的天主的崇敬，以更大的毅力加以抵制，并拒绝那冒犯祂神圣法律的事吗？\par

8. 谁嫉妒您将您想要的东西给予他人呢？我们不是您慷慨行为的审查者，也不嫉妒他人的利益，而是信仰的解释者。您将如何向\Person{基督}呈献您的礼物呢？并非很多人会按你所作的作出自己的评价，但所有人都会按你的意图这么做。他们所做的任何事都将归于您；他们不做的任何事，则将归于他们自己。虽然您是皇帝，您更应当服从天主。\Person{基督}的仆人将如何分发您的礼物呢？\par

9. 早年也曾有过此类问题，然而迫害本身却屈服于我们祖先的信仰，异教也退让了。因为在\Place{提洛}城举行五年一度的竞技会时，极恶的安提约基王前来观看，\Person{雅松}委派了圣职官员，就是安提约基人，从\Place{耶路撒冷}带去三百\OriginalTerm{达玛}{didrachms}银子，交给\Person{赫拉克勒斯}的祭献。但祖先们并没有把钱交给异教徒，而是派了忠实的人去声明，这笔钱不应用于祭神的祭献，因为这不合适，而应用于其他开支。于是颁布命令，因为他说那笔钱是为\Person{赫拉克勒斯}的祭献而送的，就应按照它被送的目的来使用；但是，带钱来的人出于他们的热忱和宗教，恳求不要用于祭献，而用于其他开支，那笔钱就被拨给了造船之用。他们被迫送去了钱，但这钱并未用于祭献，而是用于国家的其他开支。\par

10. 无疑，那些携款而来的人原可保持沉默，但这会违背他们的信德，因为他们知道那款项正被送往何处；因此他们差遣敬畏天主的人，设法使所送之物不用于\Place{圣殿}，而用于船只的费用。因为，他们把这些钱托付给了应为神圣法律辩护的人，而那赦免良心的那一位成了此事的审判者。如果他们在受制于他人时都如此谨慎，那么，皇帝啊，毫无疑问您本应怎样做。您，无论如何，无人强迫，无人能掌控您，本应向司铎寻求劝言。\par

11. 当然，当我那时抗拒时，虽是我独自一人在抗拒，但我所愿的并非独我一人，我所劝言的亦非独我一人。因此，既然我在天主和众人面前都受自己言语的约束，我觉得除了自行其是之外，既无其他许可，也无必要，因为我不能十分信任你。我长期抑制和隐藏自己的悲伤；我认为向任何人透露什么是不对的，如今我既不能再掩饰，也无权保持缄默。也正因如此，在你即位之初，当你写信给我时，我没有答复，因为我预见到会发生这样的事。最后，当你需要一封信，因为我没有写回信时，我说：这就是我觉得这件事会从我这里被强求出来的原因。\par

12. 但当我履行职责的理由出现时，我既写信又为那些为自己焦虑的人恳求，为要表明在天主的事上我感到正义的敬畏，并且我不看重阿谀奉承过于我自己的灵魂；但在那些适合向您陈情的事上，我也向权威致以应有的尊敬，正如经上所载：\BibleQuote{该尊敬的，就尊敬；该纳税的，就纳税。}\BibleRef{罗 13:7} 因为我既然从心底里尊重一个普通人，怎能不尊重皇帝呢？但您既渴望别人尊敬您，就请容许我们将尊敬献给那一位，就是您渴望被证明是您权力之源的那一位。\par

\SourceRecord{Translated by H. de Romestin, E. de Romestin and H.T.F. Duckworth.；Nicene and Post-Nicene Fathers, Second Series；Vol. 10.；Edited by Philip Schaff and Henry Wace.；Buffalo, NY: Christian Literature Publishing Co.,；1896.；<http://www.newadvent.org/fathers/340957.htm>.}

% structure: p0001 source=h1 target=section reason=page-title

\section{书信六十一}

\EditorialNote{圣安博罗削解释自己在狄奥多西皇帝战胜尤金尼乌斯后抵达米兰时缺席的原因，在表达对那次胜利的感激之后，他许诺服从皇帝的意愿，并在赞扬其虔诚的同时，敦促他对被征服者施以仁慈。}\par

\Person{安博罗削}致\Person{狄奥多西}皇帝\par

1. 您认为，至福的皇帝，——我从您的来信中得知——我远离米兰城，是因为相信您的事业已被天主所弃。但我既非如此短视，也未在远离之际忘记您的德行与功德，以至于不能预先想到上天的助佑必与您的虔诚同在，您将借此从蛮族盗匪的残暴和不配篡位者的统治下拯救罗马帝国。\par

2. 因此，我一得知那位我认为应躲避的人已经离去，便急速返回那里；我并未离弃天主审判所托付给我的米兰教会，而是躲避了那陷于亵圣罪行的人。所以，我大约在\OriginalTerm{八月朔日}{Calends of August}返回，并从那天起一直住在这里。奥古斯都啊，您的信也在这里找到了我。\par

3. 感谢我们的主天主，祂回应了您的信德和虔诚，并恢复了古代圣德的样式，使我们在我们的时代得以目睹我们在阅读圣经时所惊叹的事：即在战斗中天主的助佑如此临在，以致没有高山能延缓您前进的步伐，也没有敌军的武器构成任何障碍。\par

4. 您认为我理应为这些恩惠向主我们的天主致谢，而鉴于您的功德，我乐意如此行。那以您的名义献上的祭品必然为天主所悦纳，这是何等信德与虔敬的标记！别的皇帝在胜利之后，立即下令建造凯旋门或其他纪念其胜利的纪念碑；而您的仁慈却为天主预备祭品，并希望司铎们向主献上奉献和感恩祭。\par

5. 虽然我不配且不堪担当此职，也不堪献上这样的谢恩，但我仍要述说我所行的。我将您虔敬之函带到祭台前，放在祭台上，在献圣祭时手持此信；好使您的信德借我的声音说话，并使皇帝的信函履行司铎祭献的职责。\par

6. 确实，主恩待罗马帝国，因为祂拣选了这样一位君王和君王之父，他的德行与权能，立于如此崇高而凯旋的统治顶峰，却安于如此的谦卑；他在勇力上超越了皇帝，在谦卑上超越了司铎。我还能期望什么？还能渴求什么？您拥有一切，因此我将努力达成我渴望的总和。皇帝啊，您是慈悲的，且极具仁慈。\par

7. 至于您本人，我愿一再祈求您虔敬的增长，天主所赐没有比这更卓越的；愿您的仁慈使天主的教会，既因无辜者的和平与安宁而欢欣，也因有罪者的得蒙宽恕而喜乐。尤其要宽恕那些以前未曾冒犯的人。愿主保全您的仁慈。阿们。\par

\SourceRecord{Translated by H. de Romestin, E. de Romestin and H.T.F. Duckworth.；Nicene and Post-Nicene Fathers, Second Series；Vol. 10.；Edited by Philip Schaff and Henry Wace.；Buffalo, NY: Christian Literature Publishing Co.,；1896.；<http://www.newadvent.org/fathers/340961.htm>.}

% structure: p0001 source=h1 target=section reason=page-title

\section{书信 62}

\EditorialNote{圣\Person{盎博罗削}为失去了一次给皇帝写信的机会而道歉，但现在他派遣一位执事送去一封信，请求宽恕一些曾寻求教会庇护的\Person{尤金尼乌斯}的追随者，特别是考虑到天主赐予皇帝的神奇援助。}\par

\Person{盎博罗削}致\Person{狄奥多西}皇帝\par

1. 虽然我最近甚至第二次写信给您的仁慈，但在我看来，似乎我还未能通过类似往返回复的方式充分尽到交往之礼，因为我因您的仁慈屡赐恩惠而深感亏欠，以至我无法以任何服事来偿还所亏欠于您的，最蒙福且威严的皇帝。\par

2. 所以，正如那个最初的良机我不可错失——那时借由您的内侍，我得以向您的仁慈谢恩并履行致辞之职——尤其免得我之前未曾写信看似出于懈怠而非必要，同样，我也必须设法以某种方式向您的虔敬致以应尽的致意。\par

3. 我转而理当派遣我的儿子，执事\Person{斐理斯}，携带我的信函，同时代我向您呈上应尽之礼，并为那些逃至教会——您虔敬之母——寻求怜悯的人呈递请愿书。我无法忍受他们的眼泪，不得不在您的仁慈驾临之前前来恳求。\par

4. 我所求的是一项大恩惠，但我是向那位祈求，因主曾赐予他伟大而前所未闻之事，我深知他的仁慈，并以其虔敬为凭。因为您的胜利被认为是按古法并伴随旧时的奇迹而赐予您的，就如在圣\Person{梅瑟}、圣\Person{若苏厄}——\OriginalTerm{纳韦}{Nave}之子——、\Person{撒慕尔}和\Person{达味}身上所赐的胜利一样，不是凭人的谋算，而是由于天上恩宠的倾注。如今我们期待一份同等的仁厚，正是借着这仁厚，如此伟大的胜利才得以赢获。\par

\SourceRecord{Translated by H. de Romestin, E. de Romestin and H.T.F. Duckworth.；Nicene and Post-Nicene Fathers, Second Series；Vol. 10.；Edited by Philip Schaff and Henry Wace.；Buffalo, NY: Christian Literature Publishing Co.,；1896.；<http://www.newadvent.org/fathers/340962.htm>.}

% structure: p0001 source=h1 target=section reason=page-title

\section{第63封书信}

\EditorialNote{\Place{韦尔切利}主教\Person{利米纽斯}去世后，由于内部派系纷争，该主教职位长期空缺。因此，\Person{圣安博罗削}作为总主教，写信给\Place{韦尔切利}的基督徒，首先提及前主教\Person{优西比乌}迅速且一致当选的事，提醒他们基督同在是和睦的理由。接着他提到两个背教的修士\Person{萨尔玛提奥}和\Person{巴尔巴提安}，并痛斥纵欲使人沦为禽兽不如。进而论及主教应具备的美德，再次提到\Person{优西比乌}和\Place{米兰}主教\Person{狄奥尼修斯}，将圣职生活与隐修生活加以比较，最后以基督徒德行之劝勉结束。此信似乎写于主历396年。}\par

\Person{安博罗削}，基督的仆人，奉召作主教，致书\Place{韦尔切利}的教会，以及所有呼求吾主\Person{耶稣基督}之名的人：愿恩宠在圣神内，由天主父及祂的独生圣子，圆满地赐予你们。\par

1. 我因忧伤而耗尽，因为你们中间的天主的教会至今仍无主教，在整个\Place{利古里亚}、\Place{艾米利亚}、\Place{威尼提亚}以及\Place{意大利}邻近地区中，唯独它需要其他教会通常向它请求的那种照料；而更令我羞愧的是，你们中间造成障碍的紧张局势竟归咎于我。既然你们中间纷争不休，我们怎能作出任何决定？你们怎能推选？又有谁愿意在彼此不和的人中间接受这项职务，而即使是在和睦的人中间也难以承担？\par

2. 难道这就是\OriginalTerm{精修者}{confessor}的操练吗？这就是那些义德的祖先的后裔吗？他们一看见圣\Person{优西比乌}——他们以前从未认识他——就立刻认可他，宁可选择他而不要自己的同乡；而且他一来到他们中间，他们就认可了他，观察他之后更是如此。他成为如此伟大的人物，正是理所当然的，因为整个教会都推选了他；众人所一致要求的，正是被天主的判断所选定，这也是理所当然的。因此，你们理当效法你们祖先的榜样，尤其因为你们既受教于一位圣洁的精修者，就应当比你们的父辈更加优秀，正如一位更优秀的教师教导和训诲了你们，并且应当通过同意你们对主教的一致要求，以显示你们的节制与和睦。\par

3. 因为按照主的教导，在地上两个人若同心合意，无论求什么，我在天之父必要给他们成就，正如祂所说：\BibleQuote{因为无论在哪里，有两三个人因我的名字聚在一起，我就在他们中间。}\BibleRef{玛 18:21}那么，更何况是全体会众都奉主的名聚集呢？当众人的要求完全一致时，我们又岂能怀疑主耶稣就在那里，祂是这愿望的创始者、祈求的倾听者、\OriginalTerm{祝圣}{ordination}的主持者和恩宠的施予者呢？\par

4. 所以，你们要使自己配得上让基督在你们中间。因为哪里和平，哪里就有基督，因为基督就是和平；哪里有公义，哪里就有基督，因为基督就是公义。让祂在你们中间吧，好使你们能看见祂，免得也对你们说：\BibleQuote{在你们中间站着一位，是你们所不认识的。}\BibleRef{若 1:26}犹太人没有看见他们所不信的那一位；我们却借着虔诚瞻仰祂，借着信德注视祂。\par

5. 让祂站在你们中间，使那宣扬天主光荣的天，为你们敞开，使你们能奉行祂的旨意，作成祂的工程。谁看见耶稣，天就为他开了，就像为\Person{斯德望}打开一样，当时他说：\BibleQuote{看，我见天开了，并见人子站在天主的右边。}\BibleRef{宗 7:56}耶稣站着，是他的辩护者；祂急切地站着，为要在祂的竞技者\Person{斯德望}的竞赛中帮助他；祂站着，仿佛准备给祂的殉道者戴上荣冠。\par

6. 那么，就让祂为你们而站着，这样你们就不必害怕祂的坐姿；因为祂坐着时是审判者，正如\Person{达尼尔}所说：\BibleQuote{宝座已安放，案卷已展开，万古常存者坐下了。}\BibleRef{达 7:9}然而在\BibleBook{圣咏}{第八十二篇}中却记着：\BibleQuote{天主在神会中站起来，在众神中施审判。}因此，祂坐着时施行审判，祂站着时施行判断：对不完善者，祂施行审判；而在众神中间，祂施行判断。让祂为你们站着，作为辩护者，如同一位善牧，免得凶残的豺狼袭击你们。\par

7. 我的警告转向这一点并非徒然；因为我听说\Person{萨尔玛提奥}和\Person{巴尔巴提安}到了你们那里，他们是胡言乱语的人，宣称节欲没有功劳，节俭的生活没有恩宠，童贞毫无价值，所有人都被定在同一价格上，说那些用斋戒克苦肉身，为使肉身顺服灵魂的人是疯狂的。然而，如果\Person{保禄}宗徒认为这是疯狂，就绝不会自己这样做，也不会写信教导别人。因为他在此上夸耀，说：\BibleQuote{我痛击我身，使它为奴，免得我给别人报捷，自己反而\OriginalTerm{被弃绝}{reprobate}。}\BibleRef{格前 9:27}所以，那些不克苦肉身，又想向别人宣讲的人，他们自己正该被看作被弃绝的。\par

8. 有什么比激发奢华、腐化、淫荡，如同情欲的刺激者、享乐的引诱者、放纵的燃料、私欲的火把更该被弃绝的呢？是哪个新学派遣出这些\OriginalTerm{伊壁鸠鲁派}{Epicureans}？他们自称不是哲学学派，而是一群宣扬享乐、劝导奢华、认为贞洁无用的人。他们曾与我们在一起，却不属于我们，\BibleRef{若一 2:19}因为我们不怕重复圣史\Person{若望}所说的话。但当初他们住在这里时，曾守斋戒，被禁闭在隐修院内，没有奢华的空间，也没有嘲笑和争论的机会。\par

9. 这些娇气的人无法忍受这一切。他们离开了，后来想回来时，却不被接纳；因为我已听到许多事，使我不得不谨慎；我劝诫了他们，却毫无效果。于是他们怒不可遏，开始散布这些言论，使自己成为万恶的可悲引诱者。他们全然丧失了守斋的益处，也失去了暂时\OriginalTerm{节欲}{continence}的果实。如今，他们怀著撒殚般的狂热，嫉妒他人的善工，而那些果实他们自己未能持守。\par

10. 哪一位贞女听到她的\OriginalTerm{贞洁}{chastity}没有赏报而不叹息呢？她断然不会轻易相信，更不会放弃她的热忱，或改变她心灵的主意。哪一位寡妇，当她得知守寡毫无益处时，会选择保持婚姻的信约而生活在悲伤中，而不将自己投入更幸福的状态呢？哪一位受婚姻约束的人，若听到贞洁毫无光荣，会不被身心的轻率诱惑呢？为此，教会在神圣的读经中，在司铎的讲道中，宣扬贞洁的赞美与\OriginalTerm{童贞}{virginity}的荣耀。\par

11. 所以，宗徒的话是徒然的：\Quote{我在信上写给你们，不可与淫荡的人交往；} \BibleRef{格前 5:9} 为免他们有人说：我们不是指世上一切的淫荡者，而是说那在基督内受洗的人，如今不应被看作淫荡者，他的生活无论如何都为天主所接纳，宗徒接着补充说：\Quote{并不是指这世上所有淫荡的人，} 再往后说：\Quote{若有称为弟兄的，是淫荡的，或贪婪的，或拜偶像的，或辱骂人的，或醉酒的，或勒索人的，就不要与他交结；并且同这样的人，连一起吃饭也不可。审判教外的人，与我何干？} \BibleRef{格前 5:10} 又\BibleBook{厄弗所}{书}：\Quote{至于邪淫，一切不洁和贪婪之事，在你们中间，连提也不要提：如此才合乎圣人的身分。} \BibleRef{弗 5:3} 他随即加上：\Quote{因为你们应该清楚知道：不论是淫乱的，是污秽的，是贪婪的——即拜偶像的——在基督和天主的国内，都不得承受产业。} \BibleRef{弗 5:5} 显然这是对受洗者说的，因为他们承受产业，他们受洗归于基督的死 \BibleRef{罗 6:3} 并与他一同埋葬，为能与他一同复活。因此，他们是天主的承继者，是与基督同承继者：\BibleRef{罗 8:17} 天主的承继者，因为基督的恩宠通传给他们；与基督同承继者，因为他们被更新进入他的生命；也是基督的承继者，因为他如同立遗嘱者，藉着他的死将产业授给了他们。\par

12. 因此，那些有可能会失去什么的人，比没有的人更应当小心。他们应当更谨慎行事，应当防备邪恶的诱惑，或错误的煽动，这些都主要来自饮食。因为 \Quote{百姓坐下吃喝，起来玩乐。} \BibleRef{格前 10:7}\par

13. \Person{伊壁鸠鲁}本人，这些人认为他们应当追随他而非宗徒，他是快乐的辩护者，尽管他否认快乐带来邪恶，但并不否认由快乐产生某些事物，由此生出邪恶；他最终断言，充满快乐、生活奢侈的人似乎也不应受谴责，除非被对痛苦或死亡的恐惧所扰乱。然而，他离真理有多远，由此可知：他断言快乐起初是由天主，其创造者，在人类中创造的，正如其追随者\Person{菲洛马鲁斯}在其\WorkTitle{概要}中所论述的，他断言斯多葛派是这种意见的创始者。\par

14. 但圣经驳斥了这一点，因为它教导我们，快乐是由蛇的诡计和诱惑暗示给亚当和厄娃的。既然蛇本身就是快乐，因此快乐的偏情多种多样且难以捉摸，仿佛感染了败坏的毒液，那么，亚当因被快乐的欲望所欺骗，背离了天主的诫命和恩宠的享受，这是确定的。那么，快乐如何能召我们回归乐园，既然它独自剥夺了我们乐园？\par

15. 因此，主耶稣也愿意使我们面对魔鬼的诱惑时更加坚强，他在将要与之争战时便守了斋，叫我们明白，我们无法以其他方式战胜邪恶的诱惑。而且，魔鬼自己从快乐的箭囊中射出了第一枝试探的箭，说：\Quote{你若是天主子，就命这些石头变成饼吧！} \BibleRef{玛 4:3} 此后，主说：\Quote{人生活不只靠饼，也靠天主口中所发的一切言语。} \BibleRef{玛 4:4} 虽然他能，却没有这样做，为的是用有益的训诫教导我们，要更专心致志于研读天主的言语，而非快乐。既然他们说我们不应守斋，就让他们证明基督何以守斋，除非他的斋戒是为我们作榜样。最后，在他后来的话中，他教导我们，邪恶非因我们的斋戒便不能轻易战胜，他说：\Quote{这类魔鬼非用祈祷和斋戒，是不能赶出去的。} \BibleRef{玛 17:21}\par

16. 圣经教导我们伯多禄守了斋，正在他守斋祈祷时，获得了关于外邦人受洗的启示 \BibleRef{宗 10:10}，其用意岂不是要显示圣人们守斋时便有所长进？最后，梅瑟也是在守斋时领受了法律；\BibleRef{出 34:28} 同样，伯多禄守斋时也领受了\WorkTitle{新约}的恩宠训诲。达尼尔也因行斋戒堵住了狮子的口，看见了未来时代的事。\EditorialNote{达尼尔书六至七章} 除非我们以斋戒洗净罪恶，我们又怎能安全呢？因为圣经上说：斋戒和施舍能消除罪过。\BibleRef{多 12:8}\par

17. 那么，这些拒绝斋戒功劳的新教师是谁呢？岂不是那些说\Quote{我们吃喝吧}的外教人的声音吗？宗徒极好地嘲笑了他们，说：\Quote{我若照人的样，在\Place{以弗所}和野兽搏斗，为我有什么益处呢？如果死人不能复活，我们就吃喝吧，因为明天就要死了。} \BibleRef{格前 15:32} 就是说，我为至死之争战有什么益处，岂不为赎回我的身体吗？如果复活的希望全无，这赎回也是徒然。因此，如果复活的一切希望都已失去，我们就吃喝吧，莫要错失现世的享受，因为我们并无来世的希望。因此，那些对死后毫无希望的人，自可以沉湎于饮食。\par

18. 既然如此，宗徒为反驳这些人，便理当警告我们，不要被此类见解所动摇，他说：\BibleQuote{你们不要受迷惑：恶交败坏善德。你们应当清醒，行事端正，不要犯罪，因为有些人不认识天主。}\BibleRef{格前 15:33}因此，清醒是好的，因为酗酒是罪。\par

19. 至于那位享乐辩护者\Person{伊壁鸠鲁}本人，我们之所以屡次提到他，原是为了证明这些人若非异教徒的门徒和伊壁鸠鲁派的追随者，便是他自己，而就连哲学家们也因他是奢靡的袒护者而将其摈于群体之外——我们若能证明他甚至比这些人更可容忍，又将如何呢？据\Person{德马尔库斯}记述，他宣称：并非饮酒、宴乐、子嗣、拥抱妇人、鱼类丰盛，以及为奢宴而备的其他诸类事物，使生命甘饴，乃是清醒的言谈。最后他又补充说，那些不过度享用社交宴会的人，乃是适度地享用宴乐。凡甘愿仅以草木的汁液，伴以饼与水为饮食的人，将藐视珍馐的盛宴，因为许多不便皆由此等事物而来。在另一处，他们又说：带来享乐的并非过度的盛宴与饮酒，而是一种节制的生活。\par

20. 那么，哲学既然已经弃绝了这些人，教会岂能不排斥他们呢？再者，他们因自己的主张不义，常在自己的断言中自相矛盾。因为，尽管他们的主要见解是：除非来自饮食，便无快乐可言；但他们明白，若坚持如此可耻的定义，必蒙极大的羞辱，且遭众人遗弃，于是他们试图用一种似是而非的言论之色彩加以粉饰；以致其中一人曾说：当我们借宴会与歌曲追求快乐时，我们已失落了那借着领受\OriginalTerm{圣言}{Word}而浇灌在我们内的，唯有借此我们才能得救的恩赐。\par

21. 他们借着这些五花八门的议论，岂不是向我们显明他们彼此分歧、互不一致吗？而圣经也定了他们的罪，并未略去宗徒所驳斥的那些人，正如\Person{路加}——以历史家身份撰写那书的人——在\WorkTitle{宗徒大事录}中告诉我们：\BibleQuote{有几个伊壁鸠鲁派和斯多葛派的哲士也同他辩论；有的说：\InnerQuote{这个饶舌汉想说什么？}还有的说：\InnerQuote{他看来是在传说外邦的神明。}}\BibleRef{宗 17:18}\par

22. 然而，宗徒从这场交手也并非无所收成，因为就连\Person{亚略巴古的狄约尼削}，同他的妻子\Person{达玛黎}，以及许多别的人，都信了。如此，那群最具学识与辩才的人，竟在一次简朴的谈论中，因着那些信了的人的榜样而显得被战胜了。那么，那些力图阻挠宗徒所赢得的人，以及基督以自己的血所救赎的人，又存何居心呢？他们竟主张：受洗者不应致力于德行的操练，狂欢宴乐对他们并无害处，享乐的丰盛也无妨碍；那些放弃这些事的人才是愚昧的；贞女应当嫁人、生育子女，同样，寡妇也应当重复那曾使她吃过苦头的与男子的交往；又说，即使她们能够节欲，若不愿再次进入婚姻的纽带，便是有过错的。\par

23. 那么，你们要怎样呢？难道你们愿意我们脱去人，而穿上兽；脱去基督，而穿上魔鬼的衣裳，甚或披上更多吗？然而，既然异教徒的师傅们也认为荣誉与享乐不能并存——因为那样似乎就会将人与兽类混为一谈——我们难道也要把兽类的习性注入人的胸怀，将野兽的非理性方式，铭刻在理性的人心之上吗？\par

24. 然而，有许多种类的动物，一旦失去伴侣，便不再交配，度着仿佛是独居的生活；也有许多动物仅以素草为食，除非清泉，绝不饮用；人们也常能见到狗戒避禁食，甚至当人被禁制时，它们虽忍饥挨饿，也紧闭其口。那么，对于畜生都已学会不去触犯的事，人难道还需要被警诫吗？\par

26. 然而，有什么比节制更可钦佩的呢？节制甚至能使青春年华成熟，以致形成一种品性的老成。因为正如过度的食物与醉酒能激动已成熟的年岁一样，简朴的饮食与奔流的溪水也能减消少年的狂野。外在的火焰一经浇水便会熄灭；因此，若内体的热力因饮用溪水而得以冷却，也就不足为奇了，因为火焰的滋养或熄灭，全在于燃料。干草、麦秆、木柴、油及诸如此类之物，是供应火焰的养料；你若取走它们，或不予供应，火就熄灭了。同样地，身体的热力也因食物而得以维持或削弱，它因食物而兴奋，也因食物而减弱。因此，奢侈乃是欲望之母。\par

27. 节制岂不是合乎自然，也合乎那在创世之初便将泉水赐作饮料、树木果实赐作食物的神圣法律吗？洪水以后，那位义人发现酒成了诱惑他的根由。\BibleRef{创 9:20}那么，让我们使用节制之自然饮料吧，但愿我们人人都能如此。然而，因为不是人人都强壮，宗徒才说：\BibleQuote{为了你常有的病痛，你可以稍微用一点酒。}\BibleRef{弟前 5:23}所以，我们饮酒不应为了享乐，而是因着病弱的缘故，因此要像服药一般稀少而有限，不可过量以恣情逸乐。\par

28. 最后，\Person{厄里亚}——主所训练以臻于全德者——在头旁发现了一个饼和一瓶水；然后他仗着那食物的力量，守斋四十天四十夜。我们的祖先，当他们步行过海时，\BibleRef{出 17:6}饮的是水，不是酒。\Person{达尼尔}和希伯来青年，食用他们特别的食物，\BibleRef{达 1:8}并以水为饮料，前者胜过了狮子的凶暴；\BibleRef{达 6:22}后者则见烈火在他们身上戏耍，却毫发未伤。\BibleRef{达 3:27}\par

29. 我何必多谈男人？\Person{友弟德}丝毫不为\Person{敖罗斐乃}的奢华宴席所动，单凭她的节制，便获得了人的武力所无望取得的胜利；她救自己的国家免于占领，亲手杀死了远征军的首领。这清楚地证明了，他的放纵使那令万民畏惧的勇士软弱无力，而节制使这位女子比男人更强。在此事上，并非她的性别超越了本性，而是她以饮食之道得胜了。\Person{艾斯德尔}以禁食感动了骄傲的君王。 \BibleRef{艾 4:16}\Person{亚纳}，守寡八十四年，日夜在圣殿中以禁食和祈祷事奉天主， \BibleRef{路 2:37}她认出了基督，祂是由禁欲的大师、可说是人间新天使的\Person{若翰}所宣告的。\par

30. 愚蠢的\Person{厄里叟}啊，竟用野苦瓜喂养了先知们！健忘圣经的\Person{厄斯德拉}啊，虽然他的确凭记忆恢复了圣经！ \BibleRef{厄上 7:6}愚蠢的\Person{保禄}啊，他竟以禁食夸耀， \BibleRef{格后 11:27}如果禁食毫无益处的话。\par

31. 但若以此可洗净我们的罪，又怎能说它无益呢？如果你怀着谦卑和怜悯献上此事，正如\Person{依撒意亚}所说，你的骨骸必得滋润，你必成为一座灌溉充沛的园圃。 \BibleRef{依 58:11}这样，你的灵魂及其德行，也必因禁食的精神富饶而得以丰盈；由于你心田的肥沃，你的果实必增多，以致你内充满清醒的酣醉，就像先知所说的那杯：\BibleQuote{你的杯爵使我酣畅，何其美妙！}\par

32. 然而，不仅那节制饮食的节制堪受赞美，那节制情欲的节制亦然。因为经上记载：\BibleQuote{不要随从你的欲念而行，要克制你的贪欲。你若放任你的灵魂贪求所好，你必成为你敌人的笑柄；}\BibleRef{德 18:30}又往后说：\BibleQuote{醇酒与美女，令智慧的人跌倒。}\BibleRef{德 19:2}所以，\Person{保禄}甚至教导在婚姻中也要节制；因为那在婚姻中放纵的人，在某种意义上是奸夫，违背了宗徒的法律。\par

33. 我又何必述说童贞的恩宠何其伟大呢？它堪当为基督所拣选，甚至成为天主有形可见的殿宇，我们念到，天主性的一切完满都真实地居住其中。 \BibleRef{哥 2:9}一位童贞女孕育了世界的救恩，一位童贞女生下了万民的生命。因此，童贞不应被遗弃不顾，因它在基督内使众人获益。一位童贞女孕育了这位世界不能容纳、不能支撑的祂。当祂由母亲胎中诞生时，祂仍保守了她贞洁的围墙，以及她童贞的不可侵犯的印记。就这样，基督在那童贞女身上找到祂愿据为己有的，也就是万有的主愿意取为己有的。再者，我们的血肉因一男一女而被逐出乐园，却经由一位童贞女而与天主联结。\par

34. 关于另一位玛利亚，\Person{梅瑟}的姊妹，身为妇女的领袖徒步走过海的窄处，我又能说什么呢？ \BibleRef{出 15:20}因着同一恩赐，\Person{特克拉}也受众狮敬畏，以致那些未进食的猛兽在猎物脚前伸展，延长一次圣洁的禁食，既没有以淫邪的目光、也没有以爪牙伤害童贞女，因为童贞甚至可因一个目光而受损。\par

35. 再者，圣宗徒论及童贞时，是以何等尊敬的语气说：\BibleQuote{论到童身的人，我没有主的命令，只就我蒙主的仁慈，作为一个忠信的人，提出我的意见。}\BibleRef{格前 7:25}他没有获得命令，而是获得一项建议，因为那超越法律的事不是命令，而是以建议的方式给予劝勉。权柄未被擅自把持，但恩宠得以彰显，且这恩宠并非由任何人彰显，而是由那位蒙主仁慈的人彰显。那么，这些人的建议难道比宗徒们的建议更好吗？宗徒说：\BibleQuote{我提出我的意见，}但他们却认为宜于劝阻任何人培养童贞。\par

36. 我们应当认清，先知，或更好说在先知内的基督，以简短的诗歌对童贞作了何等称赞：\BibleQuote{我的妹妹，我的新娘，是关闭的花园，是关锁的花园，是封闭的泉源。}\BibleRef{歌 4:12}基督对教会说这话，祂愿意教会是童贞的，没有瑕疵，没有皱纹。童贞是肥沃的花园，能结出许多芬芳的果实。是关闭的花园，因为处处被贞洁的围墙封锁。是封闭的泉源，因为童贞是端庄的源头和起始，必须保持纯洁的印记不可侵犯；在这泉源中映出天主的肖像，因为单纯的洁净也与肉身的贞洁相通。\par

37. 没有人能怀疑教会是童贞女，她也在致格林多人后书中被许配并呈现为贞洁的童贞女归与基督。 \BibleRef{格后 11:2}所以，他在前书中提出他的意见，认为童贞的恩赐是好的，因为它不因现今的艰难而困扰，不被任何玷污所染污，不被任何风暴所动摇；在后书中，他引新娘归于基督，因为他能在那人民的纯洁中证实教会的童贞。\par

38. \Person{保禄}啊，请你现在回答我，你因着现今的艰难，是以何种方式提出意见。 \BibleRef{格前 7:26}他说：\BibleQuote{没有妻子的，所挂虑的是主的事，想怎样悦乐天主；}他接著又说：\BibleQuote{没有丈夫的妇女和童女，所挂虑的是主的事，一心使身心圣洁。}\BibleRef{格前 7:32}如此，她有了抵御尘世风暴的围墙，因天主护佑的防御而巩固，不受任何尘世疾风所动摇。所以，意见是好的，因为意见中有益处，而命令中则有束缚。意见吸引自愿者，命令约束不自愿者。因此，若有人接受了意见，且不后悔，她便获得了益处；但若后悔了，她也没有理由责怪宗徒，因为她本应自己判断自己的软弱；这样，她对自己的意志负责，因为她用一个超越自己能力的束缚和结索捆绑了自己。\par

39. 因此，正如一位好医生，渴望在强者中保持德行的稳固，并赐予弱者健康，他向一方提供劝告，向另一方指明补救方法：\BibleQuote{那软弱的吃蔬菜，}\BibleRef{罗 14:2}就让他娶妻吧；那更有能力的人，让他寻求德行更坚固的食粮。而且他恰当地补充道：\BibleQuote{因为那心里坚定，没有不得已的事，且能由自己作主，心里又决定了保存自己的童贞的人，作得好。所以，谁若叫自己的童贞女出嫁，作得好；谁若不叫她出嫁，作得更好。妻子是受法律束缚的，只要丈夫活着。但如果丈夫安眠了，她就自由了，可以随意嫁人，但应在主内。然而，如果她按我的劝告保持现状，她更为有福；因为我也认为，我具有主的神。}\BibleRef{格前 7:37} 这就是拥有天主的劝告：殷勤考察一切事，劝告那最好的，并指明那最稳妥的。\par

40. 谨慎的向导指出多条道路，好让每个人沿着自己喜爱并认为合宜的道路前行，只要他能够找到一条通往营地之路。\OriginalTerm{童贞}{virginity}的道路是好的，但高峻而陡峭，需要更强健的行人。\OriginalTerm{寡居}{widowhood}的道路也是好的，不如前者困难，但崎岖不平，需要更谨慎的旅人。\OriginalTerm{婚姻}{marriage}的道路同样是好的；它平坦而平直，以较长的环绕途径抵达圣人的营地。这条路为大多数人所行。因此，有童贞的赏报，有寡居的功绩，也有\OriginalTerm{夫妻贞洁}{conjugal modesty}的位置。在每一种德行中，都有阶梯和进步。\par

41. 所以你们心里要坚定站稳，免得有人推翻你们，或使你们跌倒。宗徒已教导我们\Quote{站立稳固}是什么，那就是对梅瑟所说的：\BibleQuote{你所站立的地方是圣地；}\BibleRef{出 3:5}因为没有人能站立，除非凭\OriginalTerm{信德}{faith}站立，除非在内心的决定中坚定不移地站立。在另一处经上也记载说：\BibleQuote{你且同我站在这里。}主对梅瑟所说的话，每一句都指出了\Quote{你站立的地方是圣地}和\Quote{你且同我站在这里}，意思是，你若在教会内坚定站立，便是同我站立。因为这地方本是圣的，这地面富有圣德，丰产美德的果实。\par

42. 所以你们要站立在教会中，站立在我显现给你们的地方，就是我与你们同在之处。哪里有教会，哪里就是你心灵最坚固的安息之所，哪里就有你灵魂的支撑，就是我在荆棘丛中显现给你们的地方。你们是荆棘，我是火；火在荆棘中，我在肉身内。我之所以是火，是要光照你们，是要烧尽你们的荆棘，就是你们的罪，并向你们显示我的\OriginalTerm{恩宠}{grace}。\par

43. 你们心里既然坚定站稳，就应将那企图掳掠的豺狼赶出教会。你们内不要有懈怠；你们的口不可邪恶，你们的舌也不可刻薄。不要坐在虚荣者的会中；因为经上记载说：\BibleQuote{我没有与虚荣的人同坐。}不要听那些毁谤邻人的人说话，免得你们听了别人，自己也被激动去诋毁邻人，并对你们每人都说：\BibleQuote{你坐着毁谤你的兄弟。}\par

44. 人们在诋毁别人时坐着，在赞美主时站立，对这些人说：\BibleQuote{看，请赞美主，所有主的仆人，你们这些站在主殿中的人。}坐着的人，就身体的姿态而言，仿佛因安逸而松懈，精神的能力也随之弛缓。但谨慎的守望者，警觉的探察者，护卫营垒前哨的看守者，却是站立着。那热忱的战士渴望先发制敌，更是在敌人预料之前就已列阵站立。\par

45. \BibleQuote{凡自以为站得稳的，务要小心，免得跌倒。}\BibleRef{格前 10:12} 那站立的人不屈服于毁谤，因为毁谤正流传于安逸者的话语中，并在那里暴露恶意。因此先知说：\BibleQuote{我痛恨恶人的集会，誓不与恶人同坐。}而在第三十六篇圣咏中，充满伦理规诫的那一篇，他一开头就说：\BibleQuote{不要在恶人中作恶，也不要嫉妒那些作恶的人。}\OriginalTerm{恶意}{malignancy}比单纯的\OriginalTerm{邪恶}{malice}更有害，因为恶意既无纯粹的单纯，也无公开的邪恶，而是隐藏的敌意。隐藏之事比已知之事更难防范。正因如此，我们的救主也警告我们要提防\OriginalTerm{恶神}{malignant spirits}的计谋，因为它们会以甜蜜快乐的表象和其他事物的伪装来诱捕我们，当它们展示荣誉以诱惑我们寻求野心，财富以激发贪婪，权力以引致骄傲时。\par

46. 因此，在一切行动中，尤其在拣选主教时——因这主教要作为一切生活的典范，恶意必须完全排除；好使那位由众人中拣选出来、要治愈众人的，能藉着平静而和平的决断被推举在众人之上。因为\Quote{温良的人是心灵的医生。}主在福音中也这样称呼自己，说：\BibleQuote{健康的人不需要医生，有病的人才需要。}\BibleRef{玛 9:12}\par

47. 祂是良医，背负了我们的软弱，医治了我们的疾病；然而，如经上记载，祂并未自取尊荣作大司祭，而是那向祂说话的圣父。圣父说：\BibleQuote{你是我的儿子，我今日生了你。} \BibleRef{希 5:5} 在另一处又说：\BibleQuote{你照默基瑟德的品位，永为司祭。} 祂既是一切未来司祭的预像，便取了我们的血肉，为能在祂的血肉之身时，\BibleQuote{以大声哀号和眼泪，献上了祈祷和恳求；并因所受的苦难，虽然身为天主之子，却学习了服从，且以此教导我们，好成为我们救恩的创始者。} 最后，当祂的苦难完成时，好像自身也成全完备了，便将健康赐予众人，承担了所有人的罪过。\par

48. 因此，祂自己也拣选了亚郎为司祭，为使在司祭的拣选中，不是人的意愿，而是天主的恩宠起主导作用；\BibleRef{户 16:40} 不是出于自愿的祭献，也不是自己擅取，而是从天而来的召叫，好使他能为罪人献上赎罪祭，因他自己也能同情那些犯罪的人，因为圣经上说，祂自己担荷了我们的软弱。\BibleRef{希 5:2} 谁也不应擅取这尊位，而应像亚郎一样，蒙天主召叫，\BibleRef{希 5:4} 因此，基督也没有自求，而是接受了司祭职。\par

49. 最后，由于出自亚郎家族世系的继承，所传递的更多是家族的继承者，而非他正义的分享者，于是，照着我们在旧约中所读到的那个默基瑟德的样式，真正的默基瑟德来临了，就是那真正的和平之王，真正的正义之王，因为这就是这名字的含义：\BibleQuote{无父、无母、无族谱，既没有起始之日，也没有生命之终。} \BibleRef{希 5:3} 这也指天主之子，祂在天主性的生发中无母，在由童贞玛利亚诞生时无父；在万世之前，独由圣父所生，在此世，独由童贞女出生；既然祂\BibleQuote{在起初}就有，\BibleRef{若 1:1} 自然不可能有起始之日。祂是万有的生命之源，又怎么可能有生命之终呢？祂是\BibleQuote{元始与终末。} \BibleRef{默 1:8} 但这作为典范也指向祂：司祭应当无父、无母，因为在司祭身上，蒙拣选的不是家族的高贵，而是品格的圣洁和德行的卓越。\par

50. 愿他兼具信德和品格的成熟，不偏废其一，且使二者借着善工与行为，结合为一。为此，宗徒保禄愿意我们效法那些——如他所说——\BibleQuote{凭信德和忍耐} \BibleRef{希 6:12} 承受了向亚巴郎所许诺恩许的人；亚巴郎因忍耐，堪当领受并持守了向他预许的祝福之恩。先知达味劝诫我们效法圣亚郎，并将他列入天主的圣者之中，供我们效法，说：\BibleQuote{梅瑟和亚郎列于祂的司祭中，撒慕尔列于呼求祂圣名的人中。}\par

51. 他的确堪称楷模，堪当让众人效法：当因叛逆者而招致可怕的死亡在民众中蔓延时，他站在死者与生者之间，为要阻止死亡，以免更多人丧亡。\BibleRef{户 16:48} 他真是一个怀有司祭心神的灵魂，如同善牧，以虔诚的爱，为主羊群舍命。如此，他折断了死亡的刺，遏止了它的狂暴，阻挡了它的蔓延。爱心助长了他的功绩，因为他为自己作对的人献上了自己。\par

52. 因此，那些持异议者当学会敬畏，不可激怒上主，并要平息祂的司祭。怎么！达堂、阿彼兰和科辣黑岂非因分裂而被大地吞没了吗？\BibleRef{户 16:32} 因为当科辣黑、达堂和阿彼兰煽动二百五十人起来攻击梅瑟和亚郎，欲与他们分离时，他们起来反抗说：\BibleQuote{你们太过分了！全体会众都是圣的，上主也在他们中间。} \BibleRef{户 16:3}\par

53. 于是上主动怒，向全会众发言。上主省察并认识属于祂的人，将祂的圣者引到自己身边；凡祂所没有拣选的，祂并未引近。上主吩咐科辣黑和所有随同他起来反对上主的司祭梅瑟和亚郎的人，各拿自己的香炉，加上乳香，\BibleRef{户 16:17} 好使上主所拣选的人，在上主的肋未人中，被确立为圣。\par

54. 梅瑟对科辣黑说：\BibleQuote{肋未的子孙，请听我说：天主将你们从以色列会众中分开，使你们亲近祂，在上主的会幕内供职服务，这对你们还算小事吗？} \BibleRef{户 16:8} 稍后又说道：\BibleQuote{你们还要谋求司祭职吗？以致你和你的同党联合起来反抗上主。亚郎算得了什么，你们竟抱怨他？} \BibleRef{户 16:9}\par

55. 于是，考虑到这罪过的原因：那些不配的人妄想占据司祭的职分，因而制造纷争，并抱怨指责天主拣选司祭的判决，全体百姓充满了极大的恐惧，对惩罚的畏惧临到众人。但当众人恳求不要因少数人的放肆而全都灭亡时，那些作恶的人被标记出来；那二百五十人连同他们的首领，被从全体民众中分开；随后，大地在民众中间呻吟着裂开，一个极深的深渊张开了口，罪犯被吞没，就这样从今世的一切元素中被彻底移除，既不呼吸空气而玷污空气，也不仰望天而玷污天，不触摸海而玷污海，也不以坟墓占据土地而玷污大地。\par

56. 惩罚停止了，但邪恶并未止息；因为正是从这事上，他们中起了怨言，认为百姓是因\OriginalTerm{司铎}{priests}们而丧亡。为此事，主在震怒中本来要消灭他们所有的人，若不是祂先因\Person{梅瑟}和\Person{亚郎}的祈祷而受感动，之后又因祂的司铎\Person{亚郎}的转求（这样一来，他们获得宽恕的羞辱就更大），祂愿意把生命赐给那些他们曾拒绝其特恩的人。\par

57. 女先知\Person{米黎盎}自己，曾与她的兄弟们赤足踏过海的狭处，却因仍不明白那\OriginalTerm{厄提约丕雅女子}{Ethiopian woman}的奥秘，而对她的兄弟\Person{梅瑟}口出怨言，便长出了癞病的斑点\BibleRef{户 12:10}，以致若非\Person{梅瑟}为她祈祷，她几乎难以从如此大的灾病中得解脱。虽然这抱怨所预表的是\OriginalTerm{会堂}{Synagogue}的典型，即不明了那厄提约丕雅女子的奥秘——也就是由万民聚集而成的教会，且以日常的指责抱怨不已，并嫉妒那个藉着其\OriginalTerm{信德}{faith}，它自己也必将从\OriginalTerm{不信}{unbelief}的癞病中得救的民族，正如我们所读到的：\BibleQuote{部分心硬临于以色列，直到外邦人的数目满了，这样，全以色列都将得救。}\BibleRef{罗 11:25}\par

58. 为使我们能观察到，在司铎中是天主的\OriginalTerm{恩宠}{grace}而非人的作为起作用，在\Person{梅瑟}按支派取得的许多棍杖中，只有\Person{亚郎}的杖开了花。因此百姓看出，在司铎身上应寻求的是天主召选的恩赐，于是不再要求将同等的恩宠归于人的选择，尽管他们先前认为同样的特权属于他们自己。但这根杖还显示什么，岂不正是\OriginalTerm{司铎的恩宠}{priestly grace}永不衰败，在最深的卑微中仍在其职务上具有所托付权能的花朵？抑或这件事也蕴含着奥秘？我们也不认为在司铎\Person{亚郎}生命的尽头发生此事是毫无目的的。这似乎显示，那因长久不信守司铎职务而衰残腐朽的旧民，在末期中将藉教会的榜样，在信德和热忱上重新塑造，并将以复活的恩宠，再度发出那在如此多世纪中已死去的花朵。\par

59. 但这有何意义？在\Person{亚郎}死后，主没有命令全体百姓，只命令了那在司铎中的\Person{梅瑟}一人，去给\Person{亚郎}的儿子\Person{厄肋阿匝尔}穿上\OriginalTerm{司铎的礼服}{priestly garments}，这岂不是要我们明白，\OriginalTerm{司铎}{priest}必须祝圣司铎，并亲自给他穿上礼服，即以司铎的\OriginalTerm{德行}{virtues}为衣；然后，若见他不缺司铎礼服中的任何一件，且都齐备，才可领他上神圣的\OriginalTerm{祭台}{altars}。因为那将为民恳求的人，应蒙天主拣选，并经司铎们认可，以免在他身上有任何严重冒犯，因他的职责正是为别人的过犯转求。因为司铎的德行必须非同寻常，他不仅应避免接近较重大的过失，就连极小的也应警戒。他又必须敏于怜悯，不背弃诺言，扶助跌倒者，同情痛苦，保持温和，热爱虔敬，拒绝或抑制愤怒，应像号角般激励人民虔敬，或安抚他们趋于宁静。\par

60. 有句古话说：\Quote{要习惯保持一贯，使你的生活如同一幅画，始终保持着它所得到的同一面貌。}一个人如果一时因愤怒而燃烧，又一时因暴怒而发作，一会儿满脸通红，一会儿又变为苍白，时刻变化着颜色，他怎能算是一贯的呢？但即便这是自然的事，即是说人天生易怒，或通常有原因，人的本分是要抑制愤怒，不可像狮子般被狂怒所驱使，以至不知平息，也不可传播谗言，或激化家庭纷争；因为经上记载着：\BibleQuote{易怒的人，挖掘罪恶。}\BibleRef{箴 15:18} 心怀二意的人，不能成为一贯的；在发怒时不能自制的人，更不能是一贯的，关于这一点\Person{达味}说得好：\BibleQuote{你们可以生气，但不可犯\OriginalTerm{罪}{sin}。}他不驾驭自己的怒气，而是放纵本性，这虽是人无法预防但可节制的。因此即使我们发了怒，也要让我们的情绪只容纳合乎本性的感受，而不要犯违背本性的罪。因为谁若自己都不能约束，又怎能承担约束他人的职责呢？\par

61. 因此，\OriginalTerm{宗徒}{Apostle}立下了典范，说\OriginalTerm{主教}{bishop}应是无可指摘的\BibleRef{弟前 3:2}，又在另一处说：\BibleQuote{主教应是无可指摘的，做天主的\OriginalTerm{管家}{steward}，不骄傲，不轻易发怒，不嗜酒，不打架，不贪不义之财。}\BibleRef{铎 1:7} 因为施舍者的慈悲和贪婪者的吝啬怎能相容呢？\par

62. 我记下了这些别人告诉我的应避免的事，但\OriginalTerm{宗徒}{Apostle}是德行的导师，他教导说，应以耐心驳斥反对者\BibleRef{铎 1:9}；他也规定，作主教的应只作过一个妻子的丈夫\BibleRef{铎 1:6}，这并不是要剥夺他婚姻的权利（因为这超出了诫命的能力），而是要他藉夫妇的贞洁，保存他\OriginalTerm{圣洗}{baptismal washing}的\OriginalTerm{恩宠}{grace}；也并非因宗徒的权威而让他在司铎职中生育子女；因为他所说的，是拥有子女，而非生育子女或再婚。\par

63. 我认为这一点不可略过，因为许多人主张，\Quote{只做一个妻子的丈夫}是指圣洗之后的事；这样，任何障碍的过错都将在圣洗中被洗去。事实上，所有过犯和罪过都被洗去；所以，如果有人与许多未通过婚姻律法约束的人玷污了自己的身体，一切罪过都蒙赦免；但如果有人缔结了第二次婚姻，这并不会被消除；因为\OriginalTerm{洗礼之水}{laver}解除的是罪，而不是律法，就圣洗而言，涉及的不是罪，而是律法。那么，与律法有关的事，不被看作罪而赦免，却被保留。宗徒制定了律法，说：\BibleQuote{如果谁无可指责，只做一个妻子的丈夫。} \BibleRef{弟前 3:2} 因此，无可指责、只做过一个妻子的丈夫的人，符合领受司铎职务的规定；然而，再婚者虽无玷污之罪，却丧失了司铎的特恩。\par

64. 我们已经说明了符合律法的事，现在再说明符合理性的事。但首先我们必须注意，不仅宗徒为主教或司铎定下了这条规则，而且尼西亚大公会议的教父们也补充说，凡缔结第二次婚姻的人，根本不得被接纳进入圣职界。因为，他自己在先前婚姻中都没有守住的信约，他如何能安慰或尊敬一位寡妇，或劝勉她持守寡居之身，并对丈夫保持忠信呢？或者，如果民众和司铎受同样的律法约束，他们之间又有什么区别呢？司铎的生活应超越他人，如同他的恩宠超越他人一样，因为他用规诫要求别人，自己就应当遵守律法的规诫。\par

65. 我如何抗拒我的\OriginalTerm{晋秩}{ordination}，最后，当我被迫接受时，力图至少能推迟它，但规定的准则未能胜过民众的热情。然而，西方的主教们以他们的决定赞同了我的晋秩，东方的主教们则以类似的例子予以认可。然而，禁止\OriginalTerm{新归化者}{neophyte}受晋秩，以免他因骄傲而自大。\BibleRef{弟前 3:6} 如果晋秩没有被推迟，那是由于强迫；并且，如果与司铎职相称的谦逊并不缺少，既然没有理由，就不会归咎于他。\par

66. 但是，如果说在其他教会中，为了主教的祝圣尚需如此深思熟虑，那么在\Place{韦尔切利}教会中，又需要多么谨慎呢？在那里，对主教似乎有两项同等的要求：隐修院规和教会纪律。因为，值得怀念的\Person{优西比乌}在西方土地上率先将这两种不同的事务结合了起来：他既居住在城中，遵守隐修士的规条，又藉斋戒和节制治理教会。因为，如果主教能约束青年人行节制之道、持守贞洁之规，禁止他们虽居城中却沾染都市的习气和生活方式，那么司铎的恩宠便会大大增长。\par

67. 由这样的规则，诞生了那些伟人：\Person{厄里亚}、\Person{厄里叟}、\Person{依撒伯尔}的儿子\Person{若翰}。他们身披绵羊皮，贫穷困苦，身受磨难，在旷野中、山洞和丛林间、无路的岩石间、粗糙的洞穴、深坑与沼泽中流浪，\BibleRef{希 11:37} 这世界原配不上他们。同样，\Person{达尼尔}、\Person{阿纳尼雅}、\Person{阿匝黎雅}和\Person{米沙耳}，\BibleRef{达 1:16} 他们在王宫中长大，却如在旷野中一般，以粗食淡饮为生。这些王家的奴隶，理当战胜列国，鄙视囚禁，挣脱其轭，屈服权势，征服元素，熄灭火的本性，使烈焰减弱，使刀剑锋刃变钝，堵塞狮子的口；\BibleRef{希 11:33} 当他们被认为最软弱时，却显得最为刚强，并不因人的嘲笑而退缩，因为他们仰望天上的赏报；他们也不畏惧监狱的黑暗，因为永光的荣美照耀着他们。\par

68. 效法这些人，圣\Person{优西比乌}离开了自己的家乡和亲属，宁愿异乡漂泊，也不愿在家安享。为了信德，他也宁愿并选择了流亡的艰辛，与值得怀念的\Person{狄奥尼修斯}一同如此，后者视自愿的流亡胜过皇帝的友谊。于是，这些杰出的人，虽被武装包围，被士兵围困，但当他们从大教会中被强行拖走时，却战胜了皇帝的权力，因为他们以尘世的羞辱换取了灵魂的刚毅和王者的权能；那群士兵和兵器的喧嚣不能夺走他们的信德，他们反而制服了蛮横心灵的狂暴，这狂暴无法伤害圣徒们。因为，正如你在\BibleBook{箴言}{}中所读到的，\BibleQuote{君王的震怒，有如狮子的怒吼。} \BibleRef{箴 19:12}\par

69. 他承认自己被战胜了，当他请求他们改变决心时，他们却认为自己的笔比铁剑更为有力。那时，受伤倒下的是不信，而非圣徒们的信德；他们不渴求在本土拥有坟墓，因为天上有为他们存留的家园。他们走遍大地，\BibleQuote{一无所有，却拥有一切。} \BibleRef{格后 6:10} 无论被遣往何处，他们都视之为充满喜乐的地方，因为对于拥有丰厚信德财富的人，什么都没有缺乏。最后，他们使别人富足，自己却在这世上穷乏，而在恩宠上富足。他们受试探却未被杀死，在斋戒、劳苦、警醒、守夜中，出于软弱变得坚强。因斋戒而饱饫的人不冀求享乐的诱惑；永恒的恩宠希望使他们清爽，酷暑不能烤焦他们，冰封地区的寒冷也不能摧垮他们，他们的热忱因炽热的虔诚而常新；被耶稣释放的人不畏惧人的锁链；期待基督使人复活的人不渴望脱离死亡。\par

70. 最后，圣\Person{狄约尼削}在祈祷中祈求能在流放中结束生命，因为他担心若返回家园，会发现信众或圣职人员的内心因不信者的教导或行为而受到扰乱。他获得了这一恩宠，内心平静地带着主的平安离去。因此，正如圣\Person{欧瑟伯}率先高举了\OriginalTerm{宣认信仰}{confessorship}的旗帜，真福\Person{狄约尼削}在流放中交出了生命，其荣耀甚至高于殉道者。\par

71. 圣\Person{欧瑟伯}内心的忍耐，因隐修院的纪律而变得坚固，他从艰苦忍耐的习惯中获得了承受困苦的力量。因为在较严格的基督徒虔敬中，有哪两件事是最卓越的呢？无非是圣职人员的职务和隐修士的规矩。前者是训练和蔼与善行的纪律，后者是训练节欲和忍耐的纪律；前者恰如公开舞台上的角色，后者则处于隐秘中；一个是可见的，另一个是隐藏的。因此，一如那优良的竞技者所说：\BibleQuote{我们成了供此世和天使观赏的一出戏。}\BibleRef{格前 6:9} 他确实值得天使注目观看，因为他努力争取基督的奖赏，努力在地上度天使般的生活，并要克胜天上的恶神，因为他正与\OriginalTerm{属灵的邪恶势力}{spiritual wickedness}搏斗。\BibleRef{弗 6:12} 世界理当注视他，好能效法他。\par

72. 因此，一种生活是在公开的竞技场上，另一种则如同隐藏在洞穴中；一种对抗世俗的混乱，另一种对抗肉身的欲望；一种克制，另一种回避身体的享乐；一种更令人称心，另一种更为稳妥；一种是管理，另一种是自我约束，为能完全属于基督，因为对\OriginalTerm{成全}{perfect}的人说：\BibleQuote{谁若愿意跟随我，该弃绝自己，背着自己的十字架来跟随我。}\BibleRef{玛 17:24} 如今，那能说：\BibleQuote{我生活已不是我生活，而是基督在我内生活。}\BibleRef{迦 2:20} 的人，便是跟随基督者。\par

73. \Person{保禄}弃绝了自己，当他明知在\Place{耶路撒冷}有锁链和患难等待着他时，却自愿置身于危险，说：\BibleQuote{我却不以性命为念，只要我能跑完我的路程，完成主耶稣交给我的，为天主恩宠的福音作证的\OriginalTerm{职分}{ministry}。}\BibleRef{宗 20:24} 最后，虽有许多人围着他流泪恳求，他也没有改变主意，如此，热忱的信德对自己是多么严厉的审判者。\par

74. 于是，一个去争战，另一个则退隐；一个战胜外来的挑动，另一个则躲避它们；一个藉着克胜世俗而凯旋，另一个则因世俗而欢跃；对一个而言，世界被钉在十字架上，或自己为世界被钉在十字架上，\BibleRef{迦 6:14} 对另一个而言，世界是陌生的；一个人忍受更频繁的诱惑，从而获得更大的胜利，另一个人较少跌倒，守护也更容易。\par

75. \Person{厄里亚}本人，为使由他口中所说的话得应验，奉主命去藏在\Place{革黎特}小溪旁。\Person{阿哈布}威胁他，\Person{依则贝尔}威胁他，\Person{厄里亚}害怕，就起身，随后\BibleQuote{靠着那神粮的力量，走了四十天四十夜，直走到天主的山\Place{曷勒布}；}进入一个洞穴，在那里休息；之后又被派去给君王傅油。他因居于独处而惯于忍耐，又仿佛因那简朴的食物滋养了丰盛的德行，变得更加强壮。\par

76. \Person{若翰}也是在旷野中长大，并为上主施洗，他在那里先操练了恒心，日后才能斥责君王。\par

77. 由于谈及圣\Person{厄里亚}在旷野居住时，我们未加注意地忽略了那些并非毫无寓意而命名的地名，现在回顾一下它们的含义似乎是好的。\Person{厄里亚}被差往\Place{革黎特}小溪，在那里乌鸦供养他，早上为他衔来饼，因为这饼\BibleQuote{能坚固人心。} 如果不用奥秘的食物，先知怎能被供养呢？而在晚上则有肉供给。要悟明你所读的：\Place{革黎特}意为\Quote{明悟，} \Place{曷勒布}意指\Quote{心}或\Quote{如同心，} \Place{贝尔舍巴}也被解释为\Quote{第七口井，}或\Quote{誓约的井。}\par

78. \Person{厄里亚}先去\Place{贝尔舍巴}，即神圣法律的奥迹和圣事那里，接着又被差往溪边，走向那使天主的城欢乐的河流岸。你便领悟到同一天主所立的两部约；旧约圣经如同深昧难明的井，你只能费力地从其中汲水；它尚不满全，因为那要来使之满全的还没有来到，他后来说：\BibleQuote{我来不是为废除，而是为成全法律。}\BibleRef{玛 5:17} 于是，上主吩咐圣人渡往溪流，因为凡饮于新约的，不仅是一条溪流，而且\BibleQuote{从他的腹中要流出活水的江河，}\BibleRef{若 7:38} 流出明悟的江河、默想的江河，属灵的江河，然而这些江河在不信仰的时代干涸了，免得亵渎和不信者饮用。\par

79. 在那个地方，乌鸦认出了上主的先知，而犹太人却没有认出。乌鸦喂养了他，而那王室贵胄的种族却在迫害他。迫害他的\Person{依则贝尔}不是什么，正是会堂，徒然地夸夸其谈，徒然地拥有丰富的圣经，却既不遵守，也不理解？什么是喂养他的乌鸦，不正是那些幼雏呼求上主的鸟吗？经上论到上主赐食物给它们的牲畜时如此说：\BibleQuote{向那呼求祂的乌鸦幼雏。} 这些乌鸦知道它们在喂养谁，它们接近明悟，将食物带到那神圣知识的溪流旁。\par

80. 祂养育那明了并遵守所记载之事的先知。我们的信德给他食粮，我们的进步给他养分；他以我们的心神和理智为食，他的话语因我们的理解而得到滋养。清晨我们给他饼，这饼是在福音的光照下，我们奉献给祂的我们心中坚定的力量。借着这些，他得到滋养，靠着这些，他强壮有力，用这些，他填满了那些守斋者的口，犹太人的不信未给他们提供任何信德的食粮。对他们而言，一切先知的话语不过是斋戒的饮食，他们看不见其中的内在丰盛；空虚而贫乏，无法使他们的口腹肥壮。\par

81. 或许他们在傍晚给他带来肉，作为更强固的食粮，如同格林多人，他们心智软弱，不能领受，因此宗徒用奶喂养他们。\BibleRef{格前 3:2} 如此，更强的食物在世界之傍晚被带来，在早晨则有饼。因此，因为主命令供应这食粮，那段预言的话可以恰当地在这里指向祂：\BibleQuote{祢使晨昏的出产带来喜乐；} 又接着说：\BibleQuote{祢预备了他们的食物，因为如此便是它的预备。}\par

82. 但我想关于师主已经说得足够多了，现在让我们转向门徒们的生活，他们献身于赞美祂，日夜以圣咏庆祝。因为这是天使的职务，时常专注于赞颂天主，以不断的祈祷求天主垂怜并恳求。他们专心阅读，或以持续的劳作占据心神，并且远离女性的交往，彼此提供安全的保护。这是何等的生活，其中无可怕之事，却有众多可效法之处！守斋的痛苦由内心的宁静得到补偿，因习练而减轻，因闲暇得以缓和，或以工作消磨；不为世俗挂虑所累，不卷入不合性情的烦恼，也不因城市的纷扰而沉重。\par

83. 你们明白必须找到何种教师来保存或教导这恩赐，我们若能同心合意，若有人自认受他人的伤害，你们能彼此宽恕，便能找到他。因为正义不单是不伤害那未伤害我们的人，还要宽恕那最伤害我们的人。我们常被他人的欺诈、邻人的诡计所伤；我们认为以诡计报复诡计、以欺诈回报欺诈就是美德吗？如果正义是德行，它就不应冒犯他人，也不应以恶报恶。因为，你所惩罚他人的事，自己却也同样去做，这算什么德行呢？这是恶的蔓延而非惩处，因为伤害谁并无区别，无论义人或是不义之人，人本不该伤害任何人。心中怀怨的方式也无区别，无论是出于报复的欲望，还是想要伤人，因为无论哪种情况，怀怨都难辞其咎。因为怀怨就是不义，所以经上说：\BibleQuote{不要与那怀怨的人一同怀怨，不要效法那行不义的人；} 前面又说：\BibleQuote{我憎恶那怀怨之人的集会。} 祂清楚指出一切，毫无例外，祂针对怀怨而不问原因。\par

84. 然而，有什么典范能比天主的正义更好呢？因为天主之子说：\BibleQuote{要爱你们的仇敌；} \BibleRef{玛 5:44} 又说：\BibleQuote{为那迫害你们和毁谤你们的人祈祷。} \BibleRef{玛 5:44} 祂从成全之人身上如此彻底除去报复的欲望，竟命令他们以爱德对待那伤害他们的人。并且，既然祂在旧约中说过：\BibleQuote{复仇是我的事，我必要报复。} \BibleRef{申 32:35} 祂在福音中却说，我们应为那伤害我们的人祈祷，好使那说过要报复的那一位不实行报复；因为祂愿意依照你的愿望而宽恕，这愿望与祂的应许相符。但若你寻求，你知道不义之人受自己良心的惩罚，比受司法的严惩更重。\par

85. 既然无人能免于逆境，我们应当小心，不要让逆境因我们自身的过错而临到我们。因为没有人比愚昧人更受他人严厉的审判，愚昧人既是自己不幸的肇因，便受自己的审判。为此，我们应避开那些充满烦恼和争竞、毫无益处反成阻碍的事。虽然我们应当谨慎，不使自己的决定或行为导致悔恨。因为明智的人思虑未来，以免常需后悔，而永不后悔唯独属于天主。但正义的果实，岂不是内心的宁静吗？正义地生活，岂不是宁静地生活吗？主人的榜样如何，全家的状况便如何。但若这些事在一个家庭中尚且必要，在教会中岂不更当如此，\BibleQuote{在那里，我们无论贫富、为奴或自主、希腊人或斯基泰人、尊贵或卑贱，都在基督耶稣内合而为一。} \BibleRef{哥 3:11}\par

86. 不要让人以为因他富有，就当受人更多的尊重。在教会中，富有的是那在信德上富足的人，因为有信德的人拥有全世界的财富。有信德的人拥有世界，这有何奇怪呢？他拥有基督的嗣业，其价值远超世界。\BibleQuote{你们是用宝血被赎出来的，} \BibleRef{伯前 1:18} 这话无疑是对所有人说的，不单是对富人。但若你愿意富有，就当听从那说：\BibleQuote{你们在一切举止上要圣善。} \BibleRef{伯前 1:15} 祂不是在单对富人说话，而是对所有人；因为祂不以貌取人，正如祂忠信的见证者宗徒所说的。因此他说：\BibleQuote{你们寄居此世的时日，} \BibleRef{伯前 1:17} 不要活在宴乐、挑剔或心高气傲中，而要活在敬畏中。在这地上，你们有的是时间而非永恒，当使用这时间，如同必须路过的人。\par

87. 不要倚靠财富；因为这一切事物都留在这里，唯有信德会伴随你们。如果信德引领前路，义德也必将与你们同行。财富为何诱惑你们？\BibleQuote{你们不是用金银赎回来的，} 用财产，或丝绸衣裳， \BibleQuote{脱离你们虚妄的生活，而是用基督的宝血。}\BibleRef{伯前 1:18} 那么，那继承天主产业、与基督同为继承人的才是富足的。不要轻视穷人，他使你们富足。\Quote{这穷人一呼号，主便俯听了他。} 不要拒绝一个穷人，基督原是富足的，却成了贫穷的，为你们的缘故成为贫穷的，好使你们因祂的贫穷而成为富足。\BibleRef{格后 8:9} 因此，不要如同富足的人自高，祂派遣宗徒们出去不带金钱。\par

88. 他们中的第一位说：\BibleQuote{银子和金子，我都没有。}\BibleRef{宗 3:6} 他以贫穷为荣，仿佛避免玷污。\BibleQuote{银子和金子，}他说\BibleQuote{我都没有，}——不是金银。那不知其用处的人，也不晓其价值的高低。\BibleQuote{银子和金子我都没有，}但我有信德。我因耶稣的名足够富足，\BibleQuote{那超越一切名号的名。}\BibleRef{斐 2:9} 我没有银钱，也不需任何银钱；我没有金子，也不渴望金子，但我有你们富人所没有的，我甚至有你们认为更有价值的，并且我将它施与穷人，即我因耶稣的名说：\BibleQuote{你们软弱的手坚强起来，你们颤抖的膝稳固起来。}\BibleRef{依 35:3}\par

89. 但若你们愿意成为富足的，必须成为贫穷的。那时你们若神贫，便会在一切事上富足。使人富足的不是财产，而是心神。\par

90. 有人在大量的财富中自谦自卑，他们做得好且明智，因为自然律为所有人足够丰富，据此人很快便能找到有余的；但对贪欲而言，任何财富的充裕仍是不足。再者，没有人天生贫穷，而是变得贫穷。因此，贫穷不在于本性，而在于我们自己的感受，所以为本性而言，发现自己是富足的很容易，但对贪欲却很难。因为人获得的越多，就越渴望获得，仿佛被贪欲的某种沉醉所焚烧。\par

91. 你们为何寻求财富的堆积，仿佛那是必需的？没有什么比知道这并非必需更必需的了。你们为何归咎于肉身？使人不知足的不是体内的肚腹，而是心中的贪婪。难道肉身夺去了对将来的希望吗？难道肉身摧毁了超性的恩宠的甘饴吗？难道肉身阻碍了信德吗？是肉身把虚荣的意见，如同发疯的主人一般，加在我们身上吗？肉身喜爱节制的简朴，藉此摆脱重担，披上健康，因为它放下了挂虑，获得了安宁。\par

92. 然而财富本身并非可责怪的。因为\BibleQuote{人的生命的赎价就是他的财富，}\BibleRef{箴 13:8}因为施舍给贫穷人的，赎回自己的灵魂。因此，即使在物质的财富中，也有修德之处。你们就像汪洋大海中的舵手。人若善掌航向，便迅速渡过海洋，到达港口；但那不知如何导向自己财产的人，连同他的货物一同沉没。经上如此记载：\BibleQuote{富人的财富是他的坚城。}\BibleRef{箴 10:15}\par

93. 但那城是什么，岂不是天上的耶路撒冷，在那里有天主的国吗？这份好的产业，结出永生的果实。好的产业不是留在这里，而是在那里被拥有。拥有这产业的人说：\BibleQuote{主是我的份。}他没有说，我的份从这边界延伸到那边界。他也没有说，我的份在某个邻居之间，除非是在宗徒们、先知们、主的圣人们中间，因为这就是义人的份。他没有说，我的份在草地上，或在树林中，或平原上，除非是那拥有教会的树木繁茂的平原，经上记载说：\BibleQuote{我们在树木繁茂的平原中找到了它。}他没有说，我的份在于成群的马匹，因为\BibleQuote{马是徒劳的，不能救人。}他没有说，我的份在于牛群、驴群或羊群；除非他把自己算在那认识自己主人的牲畜中，并愿意与那不避开基督驴棚的驴为伴\BibleRef{依 1:3}；那份是他的羊，就是被牵去宰杀的羊，以及那在剪毛的人前缄口无声的羔羊\BibleRef{依 53:7}，在祂的卑微中，审判已被高举。他说\BibleQuote{在剪毛的人前}的确恰当，因为祂在十字架上舍弃了附加之物，而非祂的\OriginalTerm{性体}{essence}，当时祂舍弃了身体，却没有失去天主性。\par

94. 因此，并不是人人都能说：\BibleQuote{主是我的份。}贪婪的人不能，因为贪婪走近说：你是我的份，我使你臣服于我，你曾服事我，你藉那黄金把自己卖给我，藉那财产你已将自己判给了我。奢侈的人不说：基督是我的份，因为奢侈来到说：你是我的份，我在那宴席中使你成为我的，我以那盛宴的网罗抓住了你，我以你贪食的锁链捆绑了你。你岂不知你的宴席比你的生命更被你珍视吗？我凭你自己的判断反驳你，否认吧，若你能，但你却不能。总之，你未曾为你的生命保留什么，你已为宴席耗尽了一切。奸淫的人不能说：\BibleQuote{主是我的份；}因为情欲来到说：我是你的份，你藉对那少女的爱情，藉与那娼妓共度的一夜，将自己与我捆绑，归入我的法律和权下。叛徒不能说：\BibleQuote{基督是我的份，}因为他罪恶的邪恶立刻冲向他，说：他欺骗了祢，主耶稣，他是我的。\par

95. 我们对此有一个例子，因为当\Person{犹达斯}从\Person{基督}手中接过饼时，魔鬼就进入了他的心，仿佛在索取自己的财产，仿佛在保留对自己份额的权利，仿佛在说：他不是你的，而是我的；他显然是我的仆人，是你的出卖者，明明属于我。他与你同席，却服事我；他与你一同宴饮，却由我喂养；从你手中领受饼，从我手中领受钱；他与你同饮，却把你的血卖给了我。他证明了他所说的是多么真实。于是\Person{基督}离开了他，\Person{犹达斯}也离开了\Person{耶稣}，跟随了魔鬼。\par

96. 离弃了那唯一一位的人，有多少主人啊！但我们不要离弃祂。谁会离弃祂呢？他们跟随祂，虽被锁链所捆，却是爱的锁链，这锁链释放人而不捆绑人，被这锁链捆锁的人夸耀说：\BibleQuote{耶稣基督的仆人\Person{保禄}和\Person{弟茂德}}。\BibleRef{斐 1:1} 被祂捆绑，比被他人释放和解开更为光荣。那么谁还会逃避平安？谁还会逃避救恩？谁还会逃避怜悯？谁还会逃避救赎？\par

97. 我的孩子们，你们看，那些随从这些事的人结局如何，他们虽死犹生。让我们努力学习那些我们赞赏其德行之人的勤奋，我们在他处所称赞的，在自己身上要默默认出。柔弱萎靡不能获得称赞。\BibleQuote{天国是以猛力夺取的，以猛力夺取的人，就攫取了它。}\BibleRef{玛 11:12} 祖先们仓促地吃羔羊。\OriginalTerm{信德}{faith}急速，虔敬敏捷，\OriginalTerm{望德}{hope}积极，它不爱理性的推敲，而愿从无果的闲适转向辛劳的果实。为什么拖延到明天？今天你就可获得；必须提防既不能获得明天，又失去了今天。即使一个小时的损失也不是轻微的，一个小时是我们整个生命的一部分。\par

98. 有些年轻人渴望尽快步入老年，以便不再受长辈意愿的约束；也有些老年人，如果可以，希望重返青春。这两种愿望我都不赞成，因为年轻人轻视现实，仿佛忘恩负义地渴望改变生活方式，老年人则希望延长生命，而青年的品格可以成熟，老年可以因行动而焕发青春。因为修养与年龄同样带来品格的改善。因此，我们更应高举希望于天主的国，那里有生命的新颖，有\OriginalTerm{恩宠}{grace}而非年龄的转变！\par

99. 赏报不是靠安逸或睡眠获得的。睡觉的人不做工，安逸无益，反致损失。\Person{厄撒乌}因贪安逸而失去长子的祝福，因为他宁要人把食物给他，而不愿去寻找。勤劳的\Person{雅各伯}却获得双亲的欢心。\par

100. 然而，尽管\Person{雅各伯}在德行与\OriginalTerm{恩宠}{favour}上更优越，他却向哥哥的愤怒让步，因为弟弟比他更受偏爱，哥哥为此恼怒。因此经上记载：\BibleQuote{给忿怒留有余地，}\BibleRef{罗 12:19} 免得当你想要反抗并报复时，别人的愤怒也把你卷入罪恶。你若愿意退让，便能消除他和你自己的罪。效法这位圣祖吧，他听从母亲的劝告远走他方。这母亲是谁？是\Person{黎贝加}，即\OriginalTerm{忍耐}{Patience}。除了\OriginalTerm{忍耐}{Patience}，谁能给出这样的劝告呢？母亲爱自己的儿子，却宁愿他与自己分离，而不愿他与天主分离。因此，因为母亲是善的，她使两个儿子都受益，但给了幼子一个他能持守的祝福；然而，她不是把两个儿子当作儿子而偏爱某一个，而是偏爱勤勉的胜过闲散的，忠信的胜过不信的。\par

101. 因此，他既因\OriginalTerm{虔诚}{piety}而非不敬而与父母分离，便与天主交谈，财富、子女和\OriginalTerm{恩宠}{favour}日益增长。当他遇见哥哥时，并不因这些而骄傲；反而谦卑地向哥哥下拜，他并非认为哥哥是那无怜悯、狂暴、堕落的人，而是敬畏在他内的那位。于是他下拜七次，这正是宽恕的数字，因为他不是向人下拜，而是向他在\OriginalTerm{圣神}{Spirit}内预见的、将来要取肉身除去世罪的那位下拜。\BibleRef{若 1:29} 这个奥迹在对\Person{伯多禄}的答复中向你们揭示，当\Person{伯多禄}说：\BibleQuote{如果我的弟兄得罪了我，我该宽恕他多少次？直到七次吗？}\BibleRef{玛 18:21} 你们看，罪的赦免正是那伟大\OriginalTerm{安息日}{Sabbath}、那永恒\OriginalTerm{恩宠}{grace}安息的预像，因此借\OriginalTerm{默观}{contemplation}而赐予。\par

102. 但他安排自己的妻子、儿女和所有仆人，命令他们都俯伏于地，这是什么意思？当然不是向大地的元素下拜，因为大地常沾满鲜血，是一切罪行的工厂，时常因巨石或断崖而崎岖，或土壤贫瘠饥渴；而是向为拯救我们而成为\OriginalTerm{血肉}{Flesh}的那位下拜。这或许正是主所教导的奥迹，祂说：\BibleQuote{不但是到七次，而是到七十个七次。}\BibleRef{玛 18:22}\par

103. 所以，你们要宽恕对你们所做的伤害，使你们成为\Person{雅各伯}的子女。不要像\Person{厄撒乌}那样被激怒。要效法圣\Person{达味}，他作为好老师，给我们留下了当效法的榜样，说：\BibleQuote{他们以诬告回报我的恩爱，我却只有祈祷，}当他受辱骂时，他却祈祷。祈祷是良好的盾牌，能抵挡凌辱，驱散诅咒，并常使诅咒返回到发出者身上，使他们被自己的武器所伤。他说：\BibleQuote{任凭他们咒骂，惟愿你赐福。} 人的诅咒是值得寻求的，因为它能获得\OriginalTerm{主}{Lord}的祝福。\par

104. 至于其余的事，最亲爱的弟兄们，要想到耶稣在城门外受了苦难，你们也要出离这地上的城，因为你们的城是那在天上的耶路撒冷。你们的家乡在那里，所以你们可以说：\Quote{但我们的家乡原是在天上。}\BibleRef{斐 3:20}因此，耶稣出了城，为使你们出离这世界，而能超升于世界之上。唯有梅瑟看见了天主，当他与天主说话时，他的帐幕却在营外；\BibleRef{出 33:7}并且，为罪所献祭品的血被带到祭台，而尸身却在营外焚烧；\BibleRef{出 29:12}因为凡置身于这世界的邪恶中的人，不能撇开罪恶，他的血也不蒙天主悦纳，除非他出离这肉身的污秽。\par

105. 要款待旅客，圣亚巴郎因此蒙受恩宠，接待了基督为客，已衰老的撒辣也因而得了儿子；罗特也逃脱了索多玛毁灭的烈火。你们若款待外方人，也能接待天使。对于辣哈布，她以此法获得安全，我又该说什么呢？\par

106. 要同情那些带有锁链的人，如同与他们同受捆绑。要安慰忧苦的人；因为，\BibleQuote{进入居丧的家，胜过进入宴会的家。}\BibleRef{训 7:2}从前者可得善工的功德，从后者却跌入罪恶。最后，在前者你仍盼望赏报，在后者你已得了赏报。要与受苦的人感同身受，如同自己也同受苦难。\par

107. 作妻子的要表示尊敬，勿作丈夫的奴隶；要表示自己愿意被领导，而非被强迫。该受责斥的妻子，不配婚姻。丈夫也要像舵手引导妻子，要尊敬她为自己生命的伴侣，视她为恩宠的共同继承人。\par

108. 为母亲的，要使你们的孩子断奶，爱他们，但要为他们祈祷，愿他们能长久地活在世上，但不是在地上，而是在天上，因为在地上没有什么长久的，那看似长久的事物，其实短暂且极为脆弱。要教导他们背起主的十字架，胜于贪爱现世的生命。\par

109. 主的母亲玛利亚，站在她圣子的十字架旁；这事除了圣史圣若望外，没有人告诉我。\BibleRef{若 19:25}别的圣史记述了主受难时地震、天昏地暗、太阳隐蔽；\BibleRef{玛 27:45}记述了右盗因忠信的宣认而被接入乐园。\BibleRef{路 23:43}若望却告诉我们别的圣史未曾记述的事：主如何在十字架上呼唤祂的母亲，认为祂在战胜苦难之后，向她尽孝爱之责，比赐给她天国更为重要。因为，如果按照信仰赦免右盗是虔诚之举，那么圣子以如此深情尊敬母亲，乃是更深厚孝爱的标记。祂说：\Quote{看，你的儿子}……\Quote{看，你的母亲。}\BibleRef{若 19:27}基督在十字架上作了证，并把孝爱之责分给了母亲和门徒。主不仅立了公开的遗嘱，也立了私下的遗嘱，而若望签署了祂的遗嘱，是一位无愧于如此伟大立遗嘱者的证人。这是一份美好的遗嘱，不是关乎金钱，而是关乎永生，不是用墨水写成，而是用永生天主的圣神写成；祂说：\Quote{我的舌头犹如敏捷书记的笔。}\par

110. 玛利亚也不负基督母亲的身份。宗徒们逃跑时，她站在十字架旁，以虔诚的目光注视她圣子的伤痕，因为她所期待的，不是她儿子的死亡，而是世界的得救。或者，因为那\Quote{王宫}知道，世界的救赎将借她圣子的死亡来完成，她或许以为借自己的死亡也能为公益有所贡献。但耶稣为救赎万民，并不需要助手，祂无需助手而拯救了万民。因此祂也说：\Quote{我好像成了无依无靠的人，在死人中自由自在。}祂的确接受了母亲的孝爱之情，却未寻求他人的帮助。\par

111. 圣善的母亲们，请效法她吧！她在她独一的、极可爱的圣子身上，树立了如此伟大的母德表率；因为你们的孩子并不比她更甘饴，童贞圣母也未曾寻求能再生一个儿子作为安慰。\par

112. 作主人的，命令你们的仆人，不要视他们为等级低于你们者，却要记得他们是与你们有同样本性的人。\BibleRef{伯前 2:18}作仆人的，要甘心服事你们的主人，因为各人都当耐心承受自己生来的状况，不但要服从良善的主人，也要服从乖戾的主人。如果你们热心服事良善的主人，你们有什么酬谢呢？但你们若也这样服事那乖戾的，你们就有了功绩。因为自由人若犯罪而被法官惩罚，也没有赏报；他们的功德在于无罪而受苦。所以你们，若瞻仰主耶稣，以耐心服事甚至难以相处的主人，就必有赏报。既然主自己本身是义者，却受了不义之人的苦，并以祂奇妙的忍耐，将我们的罪钉在祂的十字架上，使凡效法祂的人，能在祂的血中洗净自己的罪。\par

113. 最后，把一切转向主耶稣。让你们对今生的享受，出于纯洁的良心；你们的死亡忍耐，怀着永生的希望；你们的复活确证，藉着基督的恩宠；让真理伴同纯朴，信德伴同依恃，节德伴同圣洁，勤劳伴同节制，言谈伴同端庄，学识不沾虚荣；愿教义清醒，信德不带异端的醉态。愿主耶稣基督的恩宠与你们众人同在。阿们。\par

\SourceRecord{Translated by H. de Romestin, E. de Romestin and H.T.F. Duckworth.；Nicene and Post-Nicene Fathers, Second Series；Vol. 10.；Edited by Philip Schaff and Henry Wace.；Buffalo, NY: Christian Literature Publishing Co.,；1896.；<http://www.newadvent.org/fathers/340963.htm>.}
